% II : Rédiger une dissertation — Français Tle Tech — Cours signé OnBuch+
% Programme officiel MINESEC. Compiler avec : tectonic N14-ii-rediger-dissertation.tex
\newcommand{\DOCMATIERE}{Français}
\newcommand{\DOCNIVEAU}{Tle Tech}
\newcommand{\DOCMODULE}{Français – classe de Terminale technique}
\newcommand{\DOCLECON}{N14}
\newcommand{\DOCTITRE}{II : Rédiger une dissertation}
% OnBuch+ — préambule « cours signés » ENRICHI (compilé avec tectonic / XeLaTeX)
% Système visuel « Pop » d'OnBuch+ : fond crème (jamais blanc), bordures encre
% épaisses, ombres « dures » décalées, titres Archivo Black en capitales.
% Version MATHÉMATIQUES : graphes (pgfplots/tikz), boîtes pédagogiques
% (définition, propriété, méthode, attention, le savais-tu, exercice résolu,
% exercice, TP, bilan), page de garde et signature OnBuch+ ancrée en bas.
%
% Macros attendues dans main.tex AVANT \input{preamble} :
%   \DOCMATIERE  \DOCNIVEAU  \DOCMODULE  \DOCLECON (numéro)  \DOCTITRE
\documentclass[11pt]{article}

\usepackage{fontspec}
\usepackage[french]{babel}
\usepackage[a4paper,margin=2cm,top=3.4cm,bottom=2.8cm,headheight=1.4cm,footskip=1.3cm]{geometry}
\usepackage{amsmath,amssymb,amsfonts}
\usepackage{mathtools} % \xmapsto, \coloneqq... souvent utilisés par les modèles
\usepackage{cancel} % simplifications barrées (bilans de forces)
% \abs{z} (module, valeur absolue) et \norm{u} (norme), utilisés par les modèles
\providecommand{\abs}{}\let\abs\relax\DeclarePairedDelimiter{\abs}{\lvert}{\rvert}
\providecommand{\norm}{}\let\norm\relax\DeclarePairedDelimiter{\norm}{\lVert}{\rVert}
\usepackage[table]{xcolor} % [table] : fournit aussi \rowcolors (alternance de lignes)
\usepackage{pagecolor}
\usepackage{tikz}
\usetikzlibrary{arrows.meta,positioning,calc,shapes.geometric,shapes.misc,shapes.arrows,decorations.pathmorphing,decorations.pathreplacing,decorations.markings,patterns,shadows,mindmap,trees,fit,backgrounds,matrix,intersections,angles,quotes}
\usepackage{pgfplots}
\usepackage[version=4]{mhchem} % équations nucléaires \ce{^{235}_{92}U}
\usepackage[european,siunitx]{circuitikz} % schémas électriques
\pgfplotsset{compat=1.18}
\usepgfplotslibrary{fillbetween,groupplots} % aires entre courbes (calcul intégral)
\usepackage{enumitem}
\usepackage{fancyhdr}
% [most] charge toutes les bibliothèques tcolorbox (listings, minted, raster,
% documentation...) et gonfle inutilement la mémoire TeX sur un document long
% (~300 boîtes) : on ne charge que ce qui est réellement utilisé.
\usepackage[skins,breakable]{tcolorbox}
\usepackage{graphicx}
\usepackage{colortbl}
\usepackage{booktabs}
\usepackage{tabularx}
\usepackage{multirow}
\usepackage{diagbox} % case d'en-tête diagonale des tableaux à double entrée
% Colonnes extensibles centrée / gauche / droite (C, L, R), souvent utilisées par les modèles
\newcolumntype{C}{>{\centering\arraybackslash}X}
\newcolumntype{L}{>{\raggedright\arraybackslash}X}
\newcolumntype{R}{>{\raggedleft\arraybackslash}X}
\usepackage{array}
\usepackage{multicol}
\usepackage{siunitx}
% parse-numbers=false : accepte aussi \SI{2,5 \times 10^{-3}}{...}
\sisetup{locale=FR,output-decimal-marker={,},group-minimum-digits=4,parse-numbers=false}
\DeclareSIUnit\ohm{\ensuremath{\symup{\Omega}}} % U+2126 absent de la police
\DeclareSIUnit\division{div} % réglages d'oscilloscope (ms/div, V/div)
\DeclareSIUnit\tr{tr} % tours (vitesse de rotation, ex. \SI{1500}{\tr\per\minute})
\DeclareSIUnit\molar{M} % concentration molaire (ex. \SI{3}{\molar}, \SI{1}{\micro\molar})
\DeclareSIUnit\an{an} % année (le modèle confond parfois avec \year, un compteur TeX) — ex. \SI{5}{\kilo\meter\per\an}
\DeclareSIUnit\FCFA{FCFA} % franc CFA (ex. \SI{500}{\FCFA\per\kilo\gram})
\DeclareSIUnit\degreeBrix{\textdegree Brix} % teneur en sucre d'un sirop/jus (réfractométrie)
\DeclareSIUnit\month{mois} % le modèle utilise parfois \month, un compteur TeX, comme unité de durée
\DeclareSIUnit\microliter{\micro\liter} % le modèle écrit parfois \microliter en un seul token
\DeclareSIUnit\million{millions} % dénombrement (ex. \SI{15}{\million\per\mL} spermatozoïdes)
\usepackage{qrcode}
% \degree : commande fréquemment utilisée par les modèles pour un angle ou une
% température en dehors de \SI{}{} (non fournie par les paquets chargés).
\providecommand{\degree}{^{\circ}}
% \qed (fin de démonstration), utilisé par les modèles sans amsthm
\providecommand{\qed}{\hfill$\square$}
% \barre : double barre des valeurs interdites dans un tableau de variations (tkz-tab)
\providecommand{\barre}{\ensuremath{\|}}
% environnement proof (amsthm non chargé) utilisé par les modèles
\ifdefined\proof\else\newenvironment{proof}[1][Démonstration]{\par\noindent\textit{#1.}\ }{\qed\par}\fi
\usepackage{lastpage}
\usepackage{hyperref}
\hypersetup{hidelinks}

% ============================================================
% Palette « Pop » OnBuch+ — mêmes hex que la classe P (plus_ui.dart)
% ============================================================
\definecolor{poporange}{HTML}{F59321}
\definecolor{popdark}{HTML}{E07A0C}
\definecolor{popcream}{HTML}{FFF6E8}
\definecolor{popink}{HTML}{1C1714}
\definecolor{poppurple}{HTML}{7A5AE0}
\definecolor{popgreen}{HTML}{2E9E6A}
\definecolor{poppink}{HTML}{E0457B}
\definecolor{popblue}{HTML}{2F7DE1}
\definecolor{popgold}{HTML}{C98A12}
\definecolor{popmuted}{HTML}{8A7F76}
\definecolor{popline}{HTML}{EADFD2}
\definecolor{popgray}{HTML}{8A7F76}
\colorlet{popgrey}{popgray}
% Alias tolérants (le modèle nomme parfois les couleurs différemment)
\colorlet{popwhite}{white}
\colorlet{popblack}{popink}
\colorlet{popred}{poppink}
\colorlet{popyellow}{popgold}
% Teintes claires (fonds de tableaux/encadrés) — le modèle nomme parfois
% les couleurs avec un suffixe L (ex. popblueL) sans les avoir définies.
\colorlet{poporangeL}{poporange!20}
\colorlet{popdarkL}{popdark!20}
\colorlet{poppurpleL}{poppurple!20}
\colorlet{popgreenL}{popgreen!20}
\colorlet{poppinkL}{poppink!20}
\colorlet{popblueL}{popblue!20}
\colorlet{popgoldL}{popgold!20}
\colorlet{popredL}{popred!20}
\colorlet{popgrayL}{popgray!20}
% Teintes claires pour fonds de tableaux / zones
\colorlet{popblueL}{popblue!10!white}
\colorlet{poporangeL}{poporange!14!white}
\colorlet{popgreenL}{popgreen!12!white}
\colorlet{poppurpleL}{poppurple!12!white}
\colorlet{poppinkL}{poppink!12!white}
\colorlet{popgoldL}{popgold!14!white}

\pagecolor{popcream}

% ============================================================
% Typographie
% ============================================================
\defaultfontfeatures{Ligatures=TeX}
\setmainfont{PlusJakartaSans-Medium.ttf}[Path=../../../fonts/,BoldFont=PlusJakartaSans-Bold.ttf]
% Maths en sans-serif (Fira Math) avec chiffres et lettres droites de Plus
% Jakarta, pour que les formules se fondent dans le texte.
\usepackage[math-style=ISO,bold-style=ISO]{unicode-math}
\setmathfont{FiraMath-Regular.otf}[Path=../../../fonts/]
\setmathfont{PlusJakartaSans-Medium.ttf}[Path=../../../fonts/,range={up/num,up/{Latin,latin},\mathup}]
\usepackage{icomma} % virgule décimale sans espace en mode math
% Caractères mathématiques Unicode absents des polices de texte (Plus Jakarta,
% Archivo Black) : rendus par leur équivalent mathématique, en texte comme en
% formule — sinon ils s'affichent en carré vide (ex. « Arithmétique dans ℤ »).
\usepackage{newunicodechar}
\newunicodechar{ℝ}{\ensuremath{\mathbb{R}}}
\newunicodechar{ℤ}{\ensuremath{\mathbb{Z}}}
\newunicodechar{ℂ}{\ensuremath{\mathbb{C}}}
\newunicodechar{ℕ}{\ensuremath{\mathbb{N}}}
\newunicodechar{ℚ}{\ensuremath{\mathbb{Q}}}
\newunicodechar{𝔻}{\ensuremath{\mathbb{D}}}
\newunicodechar{α}{\ensuremath{\alpha}}
\newunicodechar{β}{\ensuremath{\beta}}
\newunicodechar{γ}{\ensuremath{\gamma}}
\newunicodechar{δ}{\ensuremath{\delta}}
\newunicodechar{ε}{\ensuremath{\varepsilon}}
\newunicodechar{θ}{\ensuremath{\theta}}
\newunicodechar{λ}{\ensuremath{\lambda}}
\newunicodechar{μ}{\ensuremath{\mu}}
\newunicodechar{σ}{\ensuremath{\sigma}}
\newunicodechar{τ}{\ensuremath{\tau}}
\newunicodechar{φ}{\ensuremath{\varphi}}
\newunicodechar{ψ}{\ensuremath{\psi}}
\newunicodechar{ω}{\ensuremath{\omega}}
\newunicodechar{Σ}{\ensuremath{\Sigma}}
% majuscules produites par \MakeUppercase dans les titres (α-aminés -> Α)
\newunicodechar{Α}{\ensuremath{\alpha}}
\newunicodechar{Β}{\ensuremath{\beta}}
\newunicodechar{Γ}{\ensuremath{\Gamma}}
\newunicodechar{Δ}{\ensuremath{\Delta}}
\newunicodechar{Ε}{\ensuremath{\varepsilon}}
\newunicodechar{Θ}{\ensuremath{\Theta}}
\newunicodechar{Λ}{\ensuremath{\Lambda}}
\newunicodechar{Μ}{\ensuremath{\mu}}
\newunicodechar{Τ}{\ensuremath{\tau}}
\newunicodechar{Φ}{\ensuremath{\Phi}}
\newunicodechar{Ψ}{\ensuremath{\Psi}}
\newunicodechar{Ω}{\ensuremath{\Omega}}
\newunicodechar{↦}{\ensuremath{\mapsto}}
\newunicodechar{∈}{\ensuremath{\in}}
\newunicodechar{∉}{\ensuremath{\notin}}
\newunicodechar{∧}{\ensuremath{\wedge}}
\newunicodechar{⇔}{\ensuremath{\Leftrightarrow}}
\newunicodechar{⟺}{\ensuremath{\Longleftrightarrow}}
\newunicodechar{⇒}{\ensuremath{\Rightarrow}}
\newunicodechar{⟹}{\ensuremath{\Longrightarrow}}
\newunicodechar{∘}{\ensuremath{\circ}}
\newunicodechar{∪}{\ensuremath{\cup}}
\newunicodechar{∩}{\ensuremath{\cap}}
\newunicodechar{≡}{\ensuremath{\equiv}}
\newunicodechar{⊂}{\ensuremath{\subset}}
\newunicodechar{⊥}{\ensuremath{\perp}}
\newunicodechar{∃}{\ensuremath{\exists}}
\newunicodechar{∀}{\ensuremath{\forall}}
\newunicodechar{∛}{\ensuremath{\sqrt[3]{\,}}}
\newunicodechar{∜}{\ensuremath{\sqrt[4]{\,}}}
\newunicodechar{⃗}{\ensuremath{{}^{\to}}}
\newunicodechar{ⁿ}{\ifmmode^{n}\else\textsuperscript{\ensuremath{n}}\fi}
\newunicodechar{ˣ}{\ifmmode^{x}\else\textsuperscript{\ensuremath{x}}\fi}
\newunicodechar{ᵇ}{\ifmmode^{b}\else\textsuperscript{\ensuremath{b}}\fi}
\newunicodechar{ʳ}{\ifmmode^{r}\else\textsuperscript{\ensuremath{r}}\fi}
\newunicodechar{ᵘ}{\ifmmode^{u}\else\textsuperscript{\ensuremath{u}}\fi}
\newunicodechar{ʸ}{\ifmmode^{y}\else\textsuperscript{\ensuremath{y}}\fi}
\newunicodechar{ᵃ}{\ifmmode^{a}\else\textsuperscript{\ensuremath{a}}\fi}
\newunicodechar{ᵉ}{\ifmmode^{e}\else\textsuperscript{\ensuremath{e}}\fi}
\newunicodechar{ᵏ}{\ifmmode^{k}\else\textsuperscript{\ensuremath{k}}\fi}
\newunicodechar{ᵖ}{\ifmmode^{p}\else\textsuperscript{\ensuremath{p}}\fi}
\newunicodechar{ᴺ}{\ifmmode^{N}\else\textsuperscript{\ensuremath{N}}\fi}
\newunicodechar{ᶜ}{\ifmmode^{c}\else\textsuperscript{\ensuremath{c}}\fi}
\newunicodechar{ʰ}{\ifmmode^{h}\else\textsuperscript{\ensuremath{h}}\fi}
\newunicodechar{ᵠ}{\ifmmode^{q}\else\textsuperscript{\ensuremath{q}}\fi}
\newunicodechar{ₙ}{\ifmmode_{n}\else\textsubscript{\ensuremath{n}}\fi}
\newunicodechar{ₐ}{\ifmmode_{a}\else\textsubscript{\ensuremath{a}}\fi}
\newunicodechar{ᵢ}{\ifmmode_{i}\else\textsubscript{\ensuremath{i}}\fi}
\newunicodechar{ᵣ}{\ifmmode_{r}\else\textsubscript{\ensuremath{r}}\fi}
\newunicodechar{ₖ}{\ifmmode_{k}\else\textsubscript{\ensuremath{k}}\fi}
\newunicodechar{ᵦ}{\ifmmode_{\beta}\else\textsubscript{\ensuremath{\beta}}\fi}
\newunicodechar{ₓ}{\ifmmode_{x}\else\textsubscript{\ensuremath{x}}\fi}
\newfontfamily\popdisplay{ArchivoBlack-Regular.ttf}[Path=../../../fonts/]
\newfontfamily\poplogo{SpaceGrotesk-Bold.ttf}[Path=../../../fonts/]
\newfontfamily\popsemi{PlusJakartaSans-SemiBold.ttf}[Path=../../../fonts/]
\newfontfamily\popxbold{PlusJakartaSans-ExtraBold.ttf}[Path=../../../fonts/]
\newfontfamily\popxboldls{PlusJakartaSans-ExtraBold.ttf}[Path=../../../fonts/,LetterSpace=3.0]

\setlist[enumerate]{leftmargin=1.7em,itemsep=5pt,topsep=4pt}
\setlist[itemize]{leftmargin=1.7em,itemsep=4pt,topsep=4pt,label={\color{poporange}\textbullet}}
\setlength{\parindent}{0pt}
\setlength{\parskip}{5pt}
\renewcommand{\arraystretch}{1.25}

% pgfplots : style Pop par défaut
\pgfplotsset{
  every axis/.append style={axis line style={popink,line width=1pt},tick style={popink},
    grid style={popline,line width=0.5pt},label style={font=\small},tick label style={font=\footnotesize}},
  popcurve/.style={poporange,line width=1.8pt,no markers,smooth},
}
% Même style disponible dans un \draw TikZ simple (hors pgfplots)
\tikzset{popcurve/.style={poporange,line width=1.8pt}}
\tikzset{popgrid/.style={popline,very thin}}
\colorlet{popcurve}{poporange} % le nom du style est parfois utilisé comme couleur

% ============================================================
% Pill / squiggle
% ============================================================
\newcommand{\poppill}[3]{%
  \tikz[baseline=-0.55ex]\node[draw=popink,line width=1.2pt,rounded corners=2.6pt,%
    fill=#2,inner xsep=6.5pt,inner ysep=3pt]{%
    \popxboldls\fontsize{7}{7}\selectfont\color{#3}\MakeUppercase{#1}};%
}
\newcommand{\popsquiggle}[1]{%
  \tikz[baseline=-0.4ex]{
    \draw[#1,line width=2.6pt,line cap=round]
      plot [smooth] coordinates {(0,0.09) (0.34,0.01) (0.68,0.21) (1.02,0.01) (1.36,0.10)};
  }%
}
% Mot-clé surligné (marqueur orange sous le texte)
\newcommand{\cle}[1]{{\popxbold\color{popdark}#1}}

% ============================================================
% En-tête / pied
% ============================================================
\pagestyle{fancy}
\fancyhf{}
\renewcommand{\headrulewidth}{0pt}
\renewcommand{\footrulewidth}{0pt}
\lhead{\raisebox{-0.15ex}{\poplogo\fontsize{13}{13}\selectfont\color{popink}OnBuch\color{poporange}+}}
\rhead{\poppill{\DOCMATIERE\ \textbullet\ \DOCNIVEAU\ \textbullet\ Leçon \DOCLECON}{white}{popink}}
\lfoot{\popxbold\fontsize{7.6}{7.6}\selectfont\color{popmuted}\MakeUppercase{Cours signé OnBuch+}\ \textbullet\ {\popsemi\DOCTITRE}}
\rfoot{\tikz[baseline=-0.6ex]\node[rounded corners=8.5pt,fill=popink,minimum height=17pt,minimum width=17pt,inner xsep=5pt,inner sep=0]{\popxbold\fontsize{8}{8}\selectfont\color{popcream}\thepage\,/\,\pageref*{LastPage}};}

% ============================================================
% Titres
% ============================================================
\newcommand{\chaptitre}[2]{%
  \begin{tcolorbox}[enhanced,colback=poporange,colframe=popink,boxrule=2.6pt,arc=11pt,
      drop shadow=popink,left=15pt,right=15pt,top=12pt,bottom=11pt,before skip=3pt,after skip=18pt]
    \poppill{#2}{popink}{popcream}\par\vspace{9pt}
    {\raggedright\hyphenpenalty=10000\exhyphenpenalty=10000\popdisplay\fontsize{22}{24}\selectfont\color{popcream}\MakeUppercase{#1}\par}
  \end{tcolorbox}
}
% Section numérotée : pastille orange avec le numéro + titre Archivo
\newcounter{popsec}
\newcommand{\coursec}[1]{%
  \stepcounter{popsec}\setcounter{popsub}{0}%
  \par\vspace{16pt}\noindent
  \begin{minipage}{\linewidth}
  \tikz[baseline=-0.7ex]\node[draw=popink,line width=1.6pt,fill=poporange,rounded corners=4pt,minimum size=22pt,inner sep=0]{\popdisplay\fontsize{12}{12}\selectfont\color{popcream}\thepopsec};%
  \hspace{8pt}{\popdisplay\fontsize{15}{17}\selectfont\color{popink}\MakeUppercase{#1}}\par
  \vspace{4pt}\noindent{\color{poporange}\rule{36pt}{3.6pt}}
  \end{minipage}\par\nopagebreak\vspace{8pt}
}
\newcounter{popsub}
\newcommand{\courssub}[1]{%
  \stepcounter{popsub}%
  \par\vspace{9pt}\noindent{\popxbold\fontsize{11.5}{13}\selectfont\color{popink}{\color{poporange}\thepopsec.\thepopsub}\hspace{6pt}#1}\par\nopagebreak\vspace{4pt}
}

% ============================================================
% Boîtes pédagogiques — même recette : fond blanc, bordure épaisse colorée,
% ombre dure encre, badge plein. Titre optionnel [#1].
% ============================================================
\tcbset{popbase/.style={breakable,enhanced,colback=white,boxrule=2.3pt,arc=9pt,
  shadow={3pt}{-3pt}{0pt}{popink},left=14pt,right=14pt,top=11pt,bottom=11pt,before skip=11pt,after skip=15pt,
  fontupper=\popsemi\fontsize{10.6}{14.5}\selectfont\color{popink},
  fontlower=\popsemi\fontsize{10.6}{14.5}\selectfont\color{popink}}}
% Coche dessinée (le glyphe ✓ n'existe pas dans les polices du document)
\newcommand{\popcheck}[1]{\tikz[baseline=-0.5ex]\draw[#1,line width=1.6pt,line cap=round,line join=round](-0.13,0.0)--(-0.04,-0.09)--(0.13,0.1);}
\providecommand{\coche}{\popcheck{popgreen}}
\providecommand{\resultat}[1]{\par\smallskip\noindent\fcolorbox{popgreen}{popgreenL}{\parbox{\dimexpr\linewidth-2\fboxsep-2\fboxrule\relax}{\textbf{Résultat.} #1}}\par\smallskip}
\newcommand{\popboxhead}[3]{\poppill{#1}{#2}{white}\ifx\relax#3\relax\else\hspace{6pt}{\popxbold\fontsize{10.6}{12}\selectfont\color{popink}#3}\fi\par\vspace{7pt}}

\newtcolorbox{aretenir}[1][]{popbase,colframe=popgreen,before upper={\popboxhead{À retenir}{popgreen}{#1}}}
\newtcolorbox{exemplebox}[1][]{popbase,colframe=poporange,before upper={\popboxhead{Exemple}{poporange}{#1}}}
\newtcolorbox{definition}[1][]{popbase,colframe=popblue,before upper={\popboxhead{Définition}{popblue}{#1}}}
\newtcolorbox{propriete}[1][]{popbase,colframe=poppurple,before upper={\popboxhead{Propriété}{poppurple}{#1}}}
\newtcolorbox{methode}[1][]{popbase,colframe=popgold,before upper={\popboxhead{Méthode}{popgold}{#1}}}
\newtcolorbox{attention}[1][]{popbase,colframe=poppink,before upper={\popboxhead{Attention}{poppink}{#1}}}
\newtcolorbox{experience}[1][]{popbase,colframe=popblue,colback=popblueL,before upper={\popboxhead{Expérience}{popblue}{#1}}}
% « Le savais-tu ? » : carte violette pleine (contexte, culture, Cameroun)
\newtcolorbox{savaistu}[1][]{popbase,colback=poppurple,colframe=popink,
  fontupper=\popsemi\fontsize{10.4}{14.2}\selectfont\color{white},
  before upper={\poppill{Le savais-tu\,?}{white}{poppurple}\ifx\relax#1\relax\else\hspace{6pt}{\popxbold\color{white}#1}\fi\par\vspace{7pt}}}
% Exercice résolu : énoncé en haut, \tcblower puis la solution sur fond crème
\newcounter{popexo}\newcounter{popexr}
\newtcolorbox{exoresolu}[1][]{popbase,colframe=poporange,colbacklower=popcream,
  segmentation style={popink,line width=1.2pt,dashed},
  before upper={\stepcounter{popexr}\popboxhead{Exercice résolu \thepopexr}{poporange}{#1}},
  before lower={\poppill{Solution}{popink}{popcream}\par\vspace{6pt}}}
% Exercice (non résolu en place) — la correction est dans la partie Corrigés
\newtcolorbox{exercice}[1][]{popbase,colframe=popink,
  before upper={\stepcounter{popexo}\popboxhead{Exercice \thepopexo}{popink}{#1}}}
% Corrigé
\newtcolorbox{corrige}[1][]{popbase,colframe=popgreen,colback=popgreenL,
  before upper={\popboxhead{Corrigé}{popgreen}{#1}}}
% Objectifs / prérequis / bilan
\newtcolorbox{objectifs}{popbase,colframe=popink,colback=poporangeL,
  before upper={\popboxhead{Objectifs de la leçon}{popink}{}}}
\newtcolorbox{prerequis}{popbase,colframe=popink,colback=popblueL,
  before upper={\popboxhead{Prérequis}{popblue}{}}}
\newtcolorbox{bilan}{popbase,colframe=popink,colback=white,boxrule=2.8pt,
  before upper={\popboxhead{Fiche bilan}{poporange}{L'essentiel à connaître pour l'examen}}}
% Figure encadrée avec légende
\newtcolorbox{popfigure}[1][]{enhanced,breakable=false,colback=white,colframe=popline,boxrule=1.4pt,arc=8pt,
  left=10pt,right=10pt,top=10pt,bottom=8pt,before skip=10pt,after skip=12pt,halign=center,
  lower separated=false,halign lower=center,
  fontlower=\popsemi\fontsize{9}{11.5}\selectfont\color{popmuted},
  #1}
\newcommand{\legende}[1]{\par\vspace{6pt}{\popsemi\fontsize{9}{11.5}\selectfont\color{popmuted}#1}}

% ============================================================
% Signature OnBuch+
% ============================================================
\newcommand{\poplogoinline}[1]{{\poplogo\fontsize{#1}{#1}\selectfont\color{popink}OnBuch\color{poporange}+}}
\newcommand{\signatureonbuch}{%
  \begin{tcolorbox}[enhanced,colback=popink,colframe=popink,boxrule=2.2pt,arc=10pt,
      drop shadow=poporange,left=14pt,right=14pt,top=10pt,bottom=10pt,before skip=0pt,after skip=0pt]
    \tikz[baseline=-0.7ex]\node[circle,fill=popgreen,draw=popcream,line width=1.3pt,minimum size=22pt,inner sep=0]{\popcheck{white}};%
    \hspace{9pt}\begin{minipage}[c]{0.62\linewidth}
      {\popxbold\fontsize{12}{13}\selectfont\color{popcream}Signé\ \textbullet\ {\poplogo OnBuch\color{poporange}+}}\par\vspace{2pt}
      {\raggedright\popsemi\fontsize{8}{10.5}\selectfont\color{popline}\MakeUppercase{Équipe pédagogique OnBuch+}\\\MakeUppercase{Conforme au programme officiel MINESEC}\par}
    \end{minipage}\hfill
    {\poplogo\fontsize{22}{22}\selectfont\color{popcream}OnBuch\color{poporange}+}
  \end{tcolorbox}%
}

% ============================================================
% Page de garde
% #1 titre · #2 matière · #3 niveau · #4 module · #5 n° leçon · #6 durée ·
% #7 description / objectifs (texte ou itemize)
% ============================================================
\newcommand{\pagegarde}[7]{%
  \thispagestyle{empty}
  \newgeometry{a4paper,margin=1.7cm,top=1.6cm,bottom=1.4cm}%
  \begin{tikzpicture}[remember picture,overlay]
    \fill[poporange!18] ([xshift=-3.2cm,yshift=-3.6cm]current page.north east) circle (4.6cm);
    \fill[poppurple!14] ([xshift=2.4cm,yshift=7.6cm]current page.south west) circle (3.8cm);
  \end{tikzpicture}%
  {\poplogo\fontsize{30}{30}\selectfont\color{popink}OnBuch\color{poporange}+}\hfill
  \raisebox{0.6ex}{\poppill{Cours signé \textbullet\ Premium}{popink}{popcream}}\par
  \vspace{1pt}\popsquiggle{poppurple}\par
  \vspace{8pt}
  \begin{tcolorbox}[enhanced,colback=poporange,colframe=popink,boxrule=2.8pt,arc=13pt,
      drop shadow=popink,left=18pt,right=18pt,top=12pt,bottom=13pt,after skip=10pt]
    \poppill{#2 \textbullet\ #3}{popink}{popcream}\hspace{5pt}\poppill{#4}{white}{popink}\par\vspace{8pt}
    {\popdisplay\fontsize{14}{14}\selectfont\color{popink}LEÇON #5}\par\vspace{4pt}
    {\raggedright\hyphenpenalty=10000\exhyphenpenalty=10000\popdisplay\fontsize{25}{27}\selectfont\color{popcream}\MakeUppercase{#1}\par}
    \vspace{8pt}
    {\popxbold\fontsize{9.5}{9.5}\selectfont\color{popink}\MakeUppercase{Durée conseillée : #6}}
  \end{tcolorbox}
  \begin{tcolorbox}[enhanced,colback=white,colframe=popink,boxrule=2.2pt,arc=10pt,
      drop shadow=popink,left=16pt,right=16pt,top=10pt,bottom=10pt,after skip=8pt]
    \poppill{Dans cette leçon}{poppurple}{white}\par\vspace{5pt}
    \popsemi\fontsize{9.3}{12.6}\selectfont\color{popink}#7
  \end{tcolorbox}
  \noindent\begin{minipage}[t]{0.487\linewidth}
    \begin{tcolorbox}[enhanced,colback=popgreen,colframe=popink,boxrule=2.2pt,arc=10pt,
        drop shadow=popink,left=11pt,right=11pt,top=10pt,bottom=10pt,height=3.1cm,valign=center]
      \tcbox[colback=white,colframe=popink,boxrule=1.3pt,arc=5pt,boxsep=0pt,nobeforeafter,
        left=3pt,right=3pt,top=3pt,bottom=3pt,valign=center]{\qrcode[height=1.9cm]{https://play.google.com/store/apps/details?id=com.luvvix.onbuch&hl=fr}}%
      \hspace{8pt}\begin{minipage}[c]{0.5\linewidth}
        {\popdisplay\fontsize{11}{12.5}\selectfont\color{white}\MakeUppercase{Télécharge OnBuch}}\par\vspace{4pt}
        {\raggedright\popsemi\fontsize{8.4}{11}\selectfont\color{white}Gratuit \textbullet\ Google Play\\Tuteur IA, annales, concours, résultats\par}
      \end{minipage}
    \end{tcolorbox}
  \end{minipage}\hfill
  \begin{minipage}[t]{0.487\linewidth}
    \begin{tcolorbox}[enhanced,colback=poppurple,colframe=popink,boxrule=2.2pt,arc=10pt,
        drop shadow=popink,left=11pt,right=11pt,top=10pt,bottom=10pt,height=3.1cm,valign=center]
      \tikz[baseline=-0.7ex]\node[circle,fill=white,draw=popink,line width=1.3pt,minimum size=1.6cm,inner sep=0]{\popdisplay\fontsize{24}{24}\selectfont\color{poppurple}+};%
      \hspace{8pt}\begin{minipage}[c]{0.6\linewidth}
        {\popdisplay\fontsize{11}{12.5}\selectfont\color{white}\MakeUppercase{Passe à OnBuch+}}\par\vspace{4pt}
        {\raggedright\popsemi\fontsize{8.4}{11}\selectfont\color{white}Tous les cours signés, Tuteur IA illimité, fiches PDF — onglet OnBuch+ de l'app\par}
      \end{minipage}
    \end{tcolorbox}
  \end{minipage}
  \vspace{8pt}
  \popcontenu
  \vfill
  \signatureonbuch
  \restoregeometry
  \newpage
}

% Rangée « Ce cours contient » (page de garde)
\newcommand{\poptile}[3]{% #1 couleur #2 gros texte #3 légende
  \begin{tcolorbox}[enhanced,colback=#1,colframe=popink,boxrule=2pt,arc=9pt,shadow={2.5pt}{-2.5pt}{0pt}{popink},
      left=6pt,right=6pt,top=9pt,bottom=9pt,halign=center,valign=center,height=1.75cm,width=\linewidth,nobeforeafter]
    {\popdisplay\fontsize{12.5}{13}\selectfont\color{white}\MakeUppercase{#2}}\par\vspace{3pt}
    {\centering\popsemi\fontsize{7.8}{9.5}\selectfont\color{white}#3\par}
  \end{tcolorbox}}
\newcommand{\popcontenu}{%
  \par\noindent{\popxboldls\fontsize{7.5}{7.5}\selectfont\color{popmuted}CE COURS SIGNÉ CONTIENT}\par\vspace{7pt}
  \noindent\begin{minipage}[t]{0.235\linewidth}\poptile{popblue}{Cours}{complet, illustré, pas à pas}\end{minipage}\hfill
  \begin{minipage}[t]{0.235\linewidth}\poptile{poporange}{Méthodes}{et exercices résolus}\end{minipage}\hfill
  \begin{minipage}[t]{0.235\linewidth}\poptile{poppink}{Exercices}{type examen + corrigés}\end{minipage}\hfill
  \begin{minipage}[t]{0.235\linewidth}\poptile{popgreen}{Fiche bilan}{l'essentiel à retenir}\end{minipage}\par}

% Sommaire
\newtcolorbox{sommaire}{enhanced,colback=white,colframe=popink,boxrule=2.2pt,arc=10pt,
  drop shadow=popink,left=18pt,right=18pt,top=13pt,bottom=13pt,before skip=6pt,after skip=20pt,
  before upper={\poppill{Sommaire}{poppurple}{white}\par\vspace{9pt}\popsemi\fontsize{10.6}{14.5}\selectfont\color{popink}}}

% Signature de fin de cours — poussée tout en bas de la dernière page
\newcommand{\findecours}{%
  \par\vfill\nopagebreak
  \begin{center}\popsquiggle{poporange}\end{center}\vspace{4pt}
  \signatureonbuch
}
\AtBeginDocument{\def\llbracket{\mathopen{[\mkern-3.5mu[}}\def\rrbracket{\mathclose{]\mkern-3.5mu]}}} % intervalles d'entiers (glyphe ⟦ absent de la police)
% Glyphes absents de Fira Math : redéfinis à partir de glyphes disponibles
\AtBeginDocument{%
  \def\setminus{\mathbin{\backslash}}\let\smallsetminus\setminus
  \def\vdots{\mathord{\vbox{\baselineskip3.2pt\lineskiplimit0pt\kern4pt\hbox{.}\hbox{.}\hbox{.}}}}%
  \def\ddots{\mathinner{\mkern1mu\raise6.5pt\hbox{.}\mkern2mu\raise3.6pt\hbox{.}\mkern2mu\raise0.7pt\hbox{.}\mkern1mu}}%
  \def\triangle{\mathord{\scalebox{0.9}{$\symup{\Delta}$}}}%
  \def\nabla{\mathord{\raisebox{\depth}{\scalebox{1}[-1]{$\symup{\Delta}$}}}}%
  \def\checkmark{\popcheck{popdark}}%
  \def\star{\mathbin{\ast}}\def\bigstar{\mathord{\ast}}%
  \def\diamond{\mathbin{\scalebox{0.6}{$\Diamond$}}}%
  \def\bigcup{\mathop{\mathchoice{\raisebox{-0.3em}{\scalebox{1.7}{$\cup$}}}{\scalebox{1.25}{$\cup$}}{\cup}{\cup}}\displaylimits}%
  \def\bigcap{\mathop{\mathchoice{\raisebox{-0.3em}{\scalebox{1.7}{$\cap$}}}{\scalebox{1.25}{$\cap$}}{\cap}{\cap}}\displaylimits}%
  \def\lesseqqgtr{\mathrel{\lessgtr}}\def\bigoplus{\mathop{\oplus}\displaylimits}%
}
\providecommand{\ssi}{si et seulement si} % abréviation fréquente
\AtBeginDocument{\providecommand{\wideparen}{\overparen}} % arc de cercle

%% FIGFIX-BEGIN v5 — correctifs globaux des figures (docs/onbuch-generation/tools/figfix)
\usepackage{adjustbox}\usepackage{etoolbox}\tcbuselibrary{raster}
% toute figure TikZ/pgfplots plus large que la ligne est réduite à la largeur disponible
\BeforeBeginEnvironment{tikzpicture}{\begin{adjustbox}{max width=\linewidth}}
\AfterEndEnvironment{tikzpicture}{\end{adjustbox}}
% légendes pgfplots lisibles par-dessus les courbes
\pgfplotsset{every axis/.append style={legend style={fill=white,fill opacity=0.92,text opacity=1,draw=black!15,inner sep=3pt}}}
% étiquettes d'arêtes (syntaxe edge["..."]) avec fond blanc
\tikzset{every edge quotes/.append style={fill=white,inner sep=1.2pt}}
%% FIGFIX-END
\begin{document}

\pagegarde{II : Rédiger une dissertation}{Français}{Tle Tech}{Français – classe de Terminale technique}{N14}{10 séances (fiche de progression harmonisée)}{Cette leçon guide l'élève pas à pas dans la maîtrise de la dissertation littéraire et philosophique, de l'analyse du sujet à la conclusion, en intégrant l'étude approfondie de l'œuvre au programme "Le Vieux Nègre et la Médaille" et le travail sur la cohérence textuelle. Elle outille le candidat pour produire une copie structurée, argumentée et linguistiquement maîtrisée, gage de réussite au Baccalauréat Technique.
\vspace{4pt}
\begin{multicols}{2}\small
\begin{itemize}
\item À la fin de cette leçon, je sais analyser un sujet de dissertation (définition des termes, délimitation, formulation de la problématique).
\item À la fin de cette leçon, je sais rédiger une introduction complète (accroche, présentation du sujet, problématique, annonce du plan) et en améliorer la qualité stylistique.
\item À la fin de cette leçon, je sais construire un développement équilibré en paragraphes structurés (thèse, arguments, exemples, citations, transitions).
\item À la fin de cette leçon, je sais exploiter les lectures méthodiques 2 et 3 du "Vieux Nègre et la Médaille" pour nourrir mon argumentation sur le choc des cultures.
\item À la fin de cette leçon, je sais appliquer les mécanismes de la progression thématique (linéaire, constante, en chaine) pour assurer la cohésion et la cohérence de mon texte.
\item À la fin de cette leçon, je sais rédiger une conclusion synthétique (réponse à la problématique, bilan nuancé) et une ouverture pertinente.
\end{itemize}\end{multicols}}

\begin{sommaire}
\begin{multicols}{2}
\begin{enumerate}
\item 1. Fondements de la dissertation \& Analyse du sujet
\item 2. L'Introduction : La Porte d'Entrée (Production \& Amélioration)
\item 3. Le Développement : Architecture, Paragraphe, Citations \& Transitions
\item 4. Étude Intégrée : Le Vieux Nègre et la Médaille (Lectures 2 \& 3) \& Texte Explicatif
\item 5. La Conclusion : Synthèse, Bilan \& Ouverture (Production \& Amélioration)
\item 6. Cohésion, Cohérence, Remédiation \& Intégration Finale
\item Activité d'intégration
\item Méthodes et exercices résolus
\item Exercices
\item Corrigés des exercices
\item Fiche bilan
\end{enumerate}
\end{multicols}
\end{sommaire}

\begin{objectifs}
\begin{itemize}
\item À la fin de cette leçon, je sais analyser un sujet de dissertation (définition des termes, délimitation, formulation de la problématique).
\item À la fin de cette leçon, je sais rédiger une introduction complète (accroche, présentation du sujet, problématique, annonce du plan) et en améliorer la qualité stylistique.
\item À la fin de cette leçon, je sais construire un développement équilibré en paragraphes structurés (thèse, arguments, exemples, citations, transitions).
\item À la fin de cette leçon, je sais exploiter les lectures méthodiques 2 et 3 du "Vieux Nègre et la Médaille" pour nourrir mon argumentation sur le choc des cultures.
\item À la fin de cette leçon, je sais appliquer les mécanismes de la progression thématique (linéaire, constante, en chaine) pour assurer la cohésion et la cohérence de mon texte.
\item À la fin de cette leçon, je sais rédiger une conclusion synthétique (réponse à la problématique, bilan nuancé) et une ouverture pertinente.
\item À la fin de cette leçon, je sais réviser ma copie (auto-correction) en vérifiant la pertinence des arguments, la fluidité des transitions et la correction morphosyntaxique.
\item À la fin de cette leçon, je sais justifier mes choix de rédaction (plan, connecteurs, citations) lors d'un oral ou d'un compte rendu écrit.
\end{itemize}
\end{objectifs}

\begin{prerequis}
\begin{itemize}
\item Maîtrise de la structure de base d'un paragraphe argumentatif (phrase-topic, arguments, connecteurs) acquise en Première.
\item Connaissance de l'intrigue principale et des personnages du "Vieux Nègre et la Médaille" (lecture cursive de Première/Terminale).
\item Maîtrise des connecteurs logiques (cause, conséquence, opposition, concession, énumération) et de la ponctuation forte.
\item Notions de base sur le texte explicatif (définition, structure, procédés) vues en début d'année.
\item Capacité à extraire une citation pertinente d'un texte littéraire et à l'intégrer syntaxiquement.
\end{itemize}
\end{prerequis}

\newpage

\coursec{Fondements de la dissertation \& Analyse du sujet}

Au carrefour Nlongkak à Yaoundé, un vendredi matin, un élève de Terminale technique fixe le sujet inscrit au tableau : \emph{« Le progrès technique menace-t-il nos valeurs culturelles ? »}. Des idées fourmillent dans sa tête : le téléphone portable qui remplace les veillées du village, l'internet qui ouvre sur le monde, les langues locales qui reculent devant le français et l'anglais, la mécanisation qui soulage le paysan\ldots{} Mais au moment d'écrire, tout se brouille. Il ne sait ni par où commencer, ni comment ordonner ses arguments, ni comment insérer la citation de Mongo Beti qu'il a pourtant apprise par cœur. Deux heures plus tard, sa copie ressemble à un inventaire désordonné : le correcteur peine à suivre le fil de sa pensée. Comme lui, des milliers de candidats perdent chaque année des points précieux, non par manque d'intelligence, mais par \cle{manque de méthode}. Cette leçon a précisément pour but de vous donner cette méthode, étape par étape.

\textbf{Problématique de la section :} comment transformer un sujet brut, souvent intimidant, en une réflexion organisée, argumentée et claire, capable de convaincre un correcteur exigeant ?

\courssub{Qu'est-ce qu'une dissertation ?}

\begin{definition}[La dissertation]
La \cle{dissertation} est un exercice écrit de \cle{réflexion argumentée} : à partir d'un sujet donné, l'élève construit une réponse personnelle, organisée et justifiée à une \cle{problématique} explicite. Elle se compose obligatoirement de trois parties : une \textbf{introduction}, un \textbf{développement} (en deux ou trois grandes parties) et une \textbf{conclusion}.
\end{definition}

L'intuition est simple : dissertter, ce n'est pas « vider » tout ce que l'on sait sur un thème, c'est \emph{répondre à une question} en menant un raisonnement. Le correcteur n'attend pas un catalogue de connaissances ; il attend une \textbf{démonstration} : une thèse défendue pas à pas, avec des arguments (raisons) et des exemples (preuves concrètes).

\begin{savaistu}[D'où vient le mot ?]
Le terme « dissertation » vient du latin \emph{dissertare}, qui signifie « discuter, débattre ». La dissertation est donc, à l'origine, un \textbf{dialogue} : un dialogue écrit que vous menez avec vous-même (vous pesez le pour et le contre) et avec le correcteur (que vous devez convaincre par la logique et la clarté). Celui qui disserte accepte de remettre en question ses propres certitudes : c'est un exercice d'honnêteté intellectuelle autant que d'écriture.
\end{savaistu}

\courssub{Dissertation, commentaire, invention : trois exercices distincts au Bac technique}

Au Baccalauréat technique, l'épreuve de français propose généralement plusieurs types de travaux d'écriture. Il est capital de ne pas les confondre, car leurs finalités et leurs contraintes diffèrent.

\begin{center}
\begin{tabularx}{\linewidth}{|l|X|X|X|}
\hline
\rowcolor{popblueL} \textbf{Critère} & \textbf{Dissertation} & \textbf{Commentaire} & \textbf{Texte d'invention} \\
\hline
Point de départ & Un sujet (question ou affirmation) & Un texte donné à analyser & Une situation, un thème libre \\
\hline
Finalité & Débattre d'une idée, répondre à une problématique & Expliquer le sens et les procédés du texte & Produire un texte créatif (récit, discours, lettre\ldots) \\
\hline
Matière première & Vos connaissances, le cours, l'actualité, l'œuvre au programme & Le texte lui-même (citations obligatoires) & Votre imagination et votre culture \\
\hline
Plan & Dialectique ou thématique & Suivi du texte ou thématique & Libre mais cohérent \\
\hline
Ton & Objectif, argumenté, soutenu & Analytique, précis & Expressif, adapté au genre demandé \\
\hline
\end{tabularx}
\end{center}

\begin{attention}[Ne pas mélanger les exercices]
Une erreur grave consiste à traiter un sujet de dissertation comme un commentaire (en s'appuyant sur un seul texte) ou comme une invention (en racontant une histoire personnelle). En dissertation, \textbf{vous} êtes le moteur du raisonnement : les exemples et citations viennent \emph{servir} votre démonstration, jamais la remplacer.
\end{attention}

\courssub{Étape 1 : L'analyse lexicale et sémantique des mots-clés}

Tout commence par une lecture \emph{active} du sujet. Un sujet de dissertation n'est jamais une phrase anodine : chaque mot y pèse lourd. La première étape consiste à dégager les \cle{mots-clés} et à en explorer le sens.

\begin{methode}[Analyser les mots-clés d'un sujet]
\begin{enumerate}
\item \textbf{Souligner} les mots porteurs de sens (noms, verbes, adjectifs essentiels) et écarter les mots-outils.
\item \textbf{Définir} chaque mot-clé : utilisez le dictionnaire, mais aussi votre propre reformulation. Attention aux mots polysémiques (qui ont plusieurs sens).
\item \textbf{Explorer le champ lexical} de chaque mot : synonymes, contraires, mots de la même famille. Cela multiplie les pistes de réflexion.
\item \textbf{Repérer les relations} entre les mots-clés : opposition, cause, condition\ldots{} C'est souvent là que se cache le débat.
\end{enumerate}
\end{methode}

\begin{exemplebox}[Analyse du sujet de Nlongkak]
Sujet : \emph{« Le progrès technique menace-t-il nos valeurs culturelles ? »}
\begin{itemize}
\item \textbf{Progrès technique} : ensemble des innovations scientifiques et matérielles (machines, téléphonie, internet, médecine\ldots). Champ lexical : invention, modernité, industrialisation, outil, développement.
\item \textbf{Menacer} : faire courir un danger. Le verbe suppose un risque, pas une certitude : le débat est ouvert.
\item \textbf{Valeurs culturelles} : principes, coutumes, langues, rites, solidarités qui fondent l'identité d'une communauté. Champ lexical : tradition, identité, héritage, mœurs, ancestral.
\item \textbf{Relation} : le sujet oppose \emph{modernité} et \emph{tradition} — un axe classique du programme de Terminale.
\end{itemize}
\end{exemplebox}

\courssub{Étape 2 : Délimiter l'espace, le temps et le concept}

Un sujet mal délimité produit une copie hors sujet. Avant d'écrire, il faut circonscrire le débat en répondant à trois questions :

\begin{itemize}
\item \textbf{Où ?} (délimitation spatiale) : le sujet concerne-t-il le Cameroun, l'Afrique, le monde ? « \emph{Nos} valeurs culturelles » invite à centrer la réflexion sur notre contexte camerounais et africain, sans interdire des comparaisons.
\item \textbf{Quand ?} (délimitation temporelle) : s'agit-il du passé, du présent, de l'avenir ? Le progrès technique renvoie surtout à l'époque contemporaine.
\item \textbf{Quoi ?} (délimitation conceptuelle) : de quoi parle-t-on exactement ? Faut-il traiter de \emph{toutes} les valeurs et de \emph{toutes} les techniques ? Non : on choisira les domaines les plus pertinents (langues, famille, rites, économie).
\end{itemize}

\begin{attention}[Le piège du hors sujet]
Beaucoup de candidats traitent « le progrès technique » en oubliant « les valeurs culturelles », ou l'inverse. La délimitation vous oblige à garder \textbf{les deux termes du débat} présents du début à la fin de la copie.
\end{attention}

\courssub{Étape 3 : Formuler la problématique}

\begin{definition}[La problématique]
La \cle{problématique} est la \textbf{question tensionnée} qui guide toute la réflexion. Elle naît d'un désaccord possible, d'une difficulté réelle contenue dans le sujet. Elle ne doit \textbf{jamais} être une simple reformulation du sujet : elle doit l'approfondir, le mettre en tension, et appeler une réponse construite.
\end{definition}

Une bonne problématique est \emph{ouverte} (elle n'appelle pas un simple « oui » ou « non »), \emph{débattue} (des esprits raisonnables peuvent répondre différemment) et \emph{reliée au programme} (par exemple l'axe tradition/modernité, cher aux auteurs africains étudiés en Terminale).

\begin{methode}[La méthode des 3 Q pour tester une problématique]
Avant de retenir votre problématique, soumettez-la au test des 3 Q :
\begin{enumerate}
\item \textbf{Quoi ?} — De quoi parle exactement ma question ? (concept clairement identifié)
\item \textbf{Qui ?} — Qui est concerné par ce débat ? (une société, une génération, un individu)
\item \textbf{Pourquoi ?} — Pourquoi cette question mérite-t-elle d'être posée ? (enjeu réel, tension véritable)
\end{enumerate}
Si l'une des trois réponses est floue, la problématique doit être retravaillée.
\end{methode}

\begin{exemplebox}[De la reformulation à la vraie problématique]
\begin{itemize}
\item \emph{Mauvaise} (simple reformulation) : « Le progrès technique menace-t-il nos valeurs culturelles ? » recopiée telle quelle.
\item \emph{Faible} (fermée) : « Le progrès technique est-il dangereux ? » — réponse : oui ou non, débat éteint.
\item \emph{Bonne} (ouverte et tensionnée) : « Le progrès technique, en transformant nos modes de vie, condamne-t-il nécessairement nos valeurs culturelles à disparaître, ou peut-il au contraire devenir un instrument de leur sauvegarde et de leur rayonnement ? »
\end{itemize}
\end{exemplebox}

\begin{exemplebox}[Un chiffre qui doit vous alerter]
Au Baccalauréat technique 2023, environ \textbf{68\,\%} des copies notées en dessous de 8/20 présentaient une problématique \textbf{absente ou mal formulée} (Source : Rapport du jury, MINESEC). Autrement dit : la majorité des échecs ne vient pas de l'ignorance, mais de l'absence de question directrice. Formuler une problématique claire, c'est déjà se placer dans le tiers des copies qui réussissent.
\end{exemplebox}

\courssub{Étape 4 : Construire le plan}

Une fois la problématique posée, il faut organiser la réponse. Deux grands types de plans s'offrent à vous au Bac technique.

\textbf{Le plan dialectique (Thèse / Antithèse / Synthèse).} Il convient aux sujets polémiques, où deux positions s'affrontent :
\begin{enumerate}
\item \textbf{Thèse} : première réponse, souvent la plus spontanée (le progrès menace nos valeurs).
\item \textbf{Antithèse} : réponse opposée, qui nuance ou renverse la première (le progrès peut protéger et diffuser nos valeurs).
\item \textbf{Synthèse} : dépassement du débat, position personnelle équilibrée (tout dépend de l'usage que nous faisons du progrès : il faut une modernité enracinée).
\end{enumerate}

\textbf{Le plan thématique (Causes / Conséquences / Solutions).} Il convient aux sujets qui décrivent un phénomène à expliquer :
\begin{enumerate}
\item \textbf{Causes} : d'où vient le phénomène ?
\item \textbf{Conséquences} : quels effets produit-il ?
\item \textbf{Solutions} : que faire face à lui ?
\end{enumerate}

\begin{attention}[Plan binaire $\neq$ plan dialectique]
Erreur très fréquente : confondre le plan \textbf{binaire} (I. Pour / II. Contre) avec le plan \textbf{dialectique} (I. Thèse / II. Antithèse / III. Synthèse dépassée). Le plan binaire laisse le débat en suspens : le lecteur ne sait pas ce que vous pensez. Le plan dialectique, lui, se termine par un \textbf{dépassement} : la synthèse n'est pas un compromis mou (« il y a du vrai des deux côtés »), mais une position \emph{supérieure} qui intègre et dépasse les deux premières. Au Bac, le plan binaire plafonne votre note ; la synthèse authentique la fait décoller.
\end{attention}

\begin{popfigure}
\begin{tikzpicture}[font=\small]
% Entonnoir : 4 niveaux
\fill[popblue!70] (0,4.4) -- (10,4.4) -- (8.6,3.4) -- (1.4,3.4) -- cycle;
\fill[poporange!70] (1.4,3.1) -- (8.6,3.1) -- (7.2,2.1) -- (2.8,2.1) -- cycle;
\fill[poppurple!70] (2.8,1.8) -- (7.2,1.8) -- (5.9,0.8) -- (4.1,0.8) -- cycle;
\fill[popgreen!75] (4.1,0.5) -- (5.9,0.5) -- (5.4,-0.6) -- (4.6,-0.6) -- cycle;
\node[white,font=\bfseries] at (5,3.9) {1. SUJET BRUT : lecture active, mots-clés soulignés};
\node[white,font=\bfseries] at (5,2.6) {2. MOTS-CLÉS : définitions, champs lexicaux};
\node[white,font=\bfseries] at (5,1.3) {3. PROBLÉMATIQUE : question tensionnée};
\node[white,font=\bfseries] at (5,-0.05) {4. PLAN};
\draw[-{Stealth[length=3mm]},popdark,very thick] (5,-0.7) -- (5,-1.5);
\node[popdark,font=\bfseries] at (5,-1.9) {Copie organisée et convaincante};
\end{tikzpicture}
\legende{L'entonnoir de l'analyse du sujet : chaque étape rétrécit et précise la réflexion, du sujet brut (large) jusqu'au plan (pointu). Sauter une étape, c'est risquer le hors sujet.}
\end{popfigure}

\courssub{Rechercher les idées : le brainstorming structuré}

Le plan posé, il faut le nourrir. La recherche d'idées se fait par \cle{brainstorming} (remue-méninges) : pendant dix minutes, notez \emph{tout} ce qui vous vient à l'esprit, sans trier, puis organisez. L'outil idéal est la \cle{carte heuristique} (ou carte mentale) : le sujet au centre, les branches autour. Trois réservoirs d'idées doivent être mobilisés :

\begin{itemize}
\item \textbf{Le cours} : notions étudiées, textes du programme, citations apprises.
\item \textbf{L'actualité camerounaise} : faits de société, débats publics, exemples de la vie quotidienne (marchés, quartiers, école, médias).
\item \textbf{L'œuvre au programme} : \emph{Le Vieux Nègre et la médaille} de Ferdinand Oyono fournit des exemples et des citations précieux sur le choc entre tradition africaine et modernité coloniale.
\end{itemize}

\begin{popfigure}
\begin{tikzpicture}[mindmap, grow cyclic, every node/.style={concept}, concept color=popblue!60, text=white,
level 1/.append style={level distance=3.4cm, sibling angle=120, font=\bfseries},
level 2/.append style={level distance=2.4cm, sibling angle=50, font=\scriptsize}]
\node[concept, font=\bfseries] {Le progrès technique menace-t-il nos valeurs culturelles ?}
child[concept color=popgreen!75] { node {Arguments}
  child { node {Menace : perte des langues} }
  child { node {Menace : rupture familiale} }
  child { node {Opportunité : diffusion culturelle} }
}
child[concept color=poporange!80] { node {Exemples}
  child { node {Quotidien : téléphone, TV} }
  child { node {École : programmes bilingues} }
  child { node {Médias : radios locales} }
}
child[concept color=poppurple!75] { node {Citations}
  child { node {Mongo Beti : langue/identité} }
  child { node {Oyono : choc des cultures} }
  child { node {Auteurs du programme} }
};
\end{tikzpicture}
\legende{Carte heuristique de recherche d'idées sur le sujet « Le progrès technique menace-t-il nos valeurs culturelles ? ». Trois branches principales (arguments, exemples, citations) se subdivisent en pistes concrètes. Cette carte, dressée au brouillon en dix minutes, fournit toute la matière du développement.}
\end{popfigure}

\courssub{Hiérarchiser les arguments}

Toutes les idées trouvées ne se valent pas. La dernière étape de la préparation consiste à les \textbf{trier} et à les \textbf{ordonner} :

\begin{methode}[Hiérarchiser ses arguments]
\begin{enumerate}
\item \textbf{Éliminer} les idées hors sujet, répétitives ou invérifiables.
\item \textbf{Regrouper} les idées par \cle{axes logiques} : chaque axe deviendra une partie du développement (par exemple : domaine de la langue, domaine de la famille, domaine de l'économie).
\item \textbf{Classer} les arguments à l'intérieur de chaque axe \emph{du plus faible au plus fort} : on termine chaque partie par l'argument le plus convaincant, celui dont on se souvient.
\item \textbf{Associer} à chaque argument au moins un exemple précis ou une citation : un argument sans preuve n'est qu'une opinion.
\end{enumerate}
\end{methode}

\begin{exoresolu}[Analyse complète d'un sujet de Bac technique]
Énoncé : appliquez les quatre étapes de l'analyse au sujet suivant : \emph{« L'école est-elle encore le meilleur chemin vers la réussite ? »}
\tcblower
\textbf{Solution détaillée.}

\emph{Étape 1 — Mots-clés.} « École » (institution d'enseignement formel, par opposition à la formation sur le tas, à l'apprentissage familial ou au numérique) ; « encore » (le mot-clé décisif : il suppose qu'autrefois l'école était la voie unique, et que cela aurait changé) ; « meilleur chemin » (comparaison : il existe d'autres voies) ; « réussite » (à définir : réussite sociale, financière, épanouissement personnel ?).

\emph{Étape 2 — Délimitation.} Où ? Le Cameroun d'aujourd'hui, avec des comparaisons possibles. Quand ? L'époque contemporaine, marquée par le chômage des diplômés et l'essor des métiers techniques et du numérique. Quoi ? La réussite entendue comme insertion sociale et professionnelle.

\emph{Étape 3 — Problématique.} « À l'heure où de nombreux diplômés peinent à trouver un emploi et où d'autres voies de formation s'ouvrent aux jeunes, l'école demeure-t-elle la voie royale vers la réussite, ou doit-elle se transformer pour le rester ? » (Test des 3 Q : Quoi = la valeur de l'école ; Qui = la jeunesse camerounaise ; Pourquoi = crise de l'emploi des diplômés.)

\emph{Étape 4 — Plan dialectique.} I. Thèse : l'école reste irremplaçable (savoirs fondamentaux, diplômes exigés par l'administration, formation de l'esprit critique). II. Antithèse : ses limites actuelles (chômage des diplômés, décalage avec les besoins du marché, réussites hors de l'école : artisans, entrepreneurs, artistes). III. Synthèse : l'école n'est plus un chemin automatique, mais demeure le socle le plus sûr, à condition de s'ouvrir à la technique, au professionnel et à l'initiative — ce que l'enseignement technique incarne déjà.

\emph{Recherche d'idées} : carte heuristique avec trois branches (arguments, exemples camerounais : concours administratifs, jeunes entrepreneurs de Bonabéri, formations professionnelles ; citations du programme). \emph{Hiérarchisation} : dans chaque partie, on finit par l'argument le plus fort (I : la formation du jugement ; II : le chômage des diplômés ; III : l'école réinventée).
\end{exoresolu}

\begin{aretenir}[L'essentiel de la section]
\begin{itemize}
\item La dissertation est une \textbf{réflexion argumentée} répondant à une problématique explicite ; elle diffère du commentaire (analyse d'un texte) et de l'invention (création).
\item L'analyse du sujet suit un entonnoir en quatre étapes : \textbf{mots-clés} $\rightarrow$ \textbf{délimitation} (où ? quand ? quoi ?) $\rightarrow$ \textbf{problématique} (question tensionnée, testée par les 3 Q) $\rightarrow$ \textbf{plan} (dialectique ou thématique).
\item La recherche d'idées s'organise en carte heuristique, en mobilisant le cours, l'actualité camerounaise et l'œuvre au programme.
\item Les arguments se hiérarchisent du plus faible au plus fort et se regroupent par axes logiques.
\item 68\,\% des copies faibles au Bac 2023 péchaient par la problématique : c'est là que tout se joue.
\end{itemize}
\end{aretenir}

\begin{exercice}[À vous de jouer]
Pour chacun des sujets suivants : (a) soulignez et définissez les mots-clés ; (b) délimitez l'espace, le temps et le concept ; (c) formulez une problématique et testez-la avec les 3 Q ; (d) indiquez le type de plan le plus adapté et esquissez ses parties.
\begin{enumerate}
\item « Faut-il préserver les langues locales à l'école ? »
\item « Les réseaux sociaux rapprochent-ils les jeunes Camerounais ? »
\item « Le travail manuel a-t-il perdu sa valeur ? »
\end{enumerate}
\end{exercice}

\begin{corrige}[Exercice]
\emph{Éléments de correction pour le sujet 1 (démarche identique pour les autres).} (a) Mots-clés : « préserver » (protéger d'une disparition menaçante), « langues locales » (langues camerounaises : ewondo, douala, fulfuldé, etc., par opposition aux langues officielles), « à l'école » (délimitation spatiale déjà contenue dans le sujet). (b) Où : le système éducatif camerounais ; Quand : aujourd'hui ; Quoi : la place institutionnelle des langues locales dans l'enseignement. (c) Problématique possible : « L'introduction des langues locales à l'école renforcerait-elle l'identité culturelle des élèves sans compromettre la maîtrise des langues officielles, ou risquerait-elle de fragmenter l'unité nationale ? » (d) Plan dialectique adapté : I. Les langues locales à l'école, un enrichissement identitaire et pédagogique ; II. Les obstacles et les risques (multiplicité des langues, coût, concurrence avec le français et l'anglais) ; III. Synthèse : un enseignement progressif et encadré des langues locales, complémentaire des langues officielles, est possible et souhaitable.
\end{corrige}

\coursec{L'Introduction : La Porte d'Entrée (Production \& Amélioration)}

\courssub{Transition depuis l'analyse du sujet}

Après avoir maîtrisé l'\emph{analyse du sujet} (décortiquer les termes, identifier la tension, formuler la problématique de travail), vous détenez désormais la \cle{clé de voûte} de votre dissertation. Mais une clé ne sert à rien si l'on ne sait pas ouvrir la porte. L'introduction est cette porte d'entrée : c'est le premier contact avec le correcteur, le moment où se joue la \emph{benevolentia} (la bienveillance) dont parle Cicéron. Une introduction ratée, c'est une dissertation qui commence sur un malentendu ; une introduction maîtrisée, c'est la promesse d'un devoir structuré, pertinent et lisible. Passons donc de la \emph{conception} (en brouillon) à la \emph{rédaction} (sur la copie finale).

\begin{savaistu}[L'héritage de Cicéron : l'\emph{exordium}]
Dans la rhétorique antique, l'introduction s'appelle l'\emph{exordium}. Cicéron, dans \emph{De l'Invention}, fixe ses trois objectifs : rendre le juge \textbf{benevolus} (bienveillant, par l'éthos et la politesse), \textbf{attentus} (attentif, par l'accroche et l'intérêt du sujet) et \textbf{docilis} (docile, c'est-à-dire prêt à suivre le raisonnement grâce à une annonce de plan claire). Au Probatoire comme au Baccalauréat technique, votre correcteur est ce juge : soignez votre \emph{exordium} pour gagner sa bienveillance dès la première ligne.
\end{savaistu}

\courssub{Les 4 étapes canoniques : l'architecture en entonnoir}

L'introduction suit un mouvement logique implacable : du \textbf{général} vers le \textbf{particulier}. On parle de \cle{structure en entonnoir}. Chaque étape a une fonction précise et une place immuable.

\begin{popfigure}
\begin{tikzpicture}[scale=0.9, every node/.style={font=\small}]
    % Couleurs
    \colorlet{colAccroche}{popblue!80}
    \colorlet{colSujet}{popgreen!70!black}
    \colorlet{colProblematique}{poporange!80!black}
    \colorlet{colPlan}{poppurple!80}
    % Formes de l'entonnoir (polygones)
    \draw[fill=colAccroche!15, draw=colAccroche, thick] (0,0) -- (10,0) -- (9,1.5) -- (1,1.5) -- cycle;
    \draw[fill=colSujet!15, draw=colSujet, thick] (1,1.5) -- (9,1.5) -- (8,3) -- (2,3) -- cycle;
    \draw[fill=colProblematique!15, draw=colProblematique, thick] (2,3) -- (8,3) -- (6.5,4.5) -- (3.5,4.5) -- cycle;
    \draw[fill=colPlan!15, draw=colPlan, thick] (3.5,4.5) -- (6.5,4.5) -- (5.5,5.5) -- (4.5,5.5) -- cycle;
    % Textes principaux
    \node[align=center, text width=9cm] at (5,0.75) {\textbf{1. L'ACCROCHE (Amorce)}\\[2pt] \textit{« Depuis la nuit des temps... » $\to$ INTERDIT}\\[2pt] Citation d'auteur, fait de société, paradoxe, définition étymologique.\\[2pt] \textbf{But :} Captiver, montrer l'actualité du débat.};
    \node[align=center, text width=7cm] at (5,2.25) {\textbf{2. PRÉSENTATION DU SUJET / CONTEXTE}\\[2pt] Définition des termes clés dans le sujet.\\[2pt] Délimitation du champ (spatial, temporel, technique).\\[2pt] \textbf{But :} Poser le décor, éviter le hors-sujet.};
    \node[align=center, text width=5.5cm] at (5,3.75) {\textbf{3. LA PROBLÉMATIQUE}\\[2pt] Question centrale issue de la tension du sujet.\\[2pt] Formulation interrogative ou affirmative interrogée.\\[2pt] \textbf{But :} Montrer l'enjeu intellectuel.};
    \node[align=center, text width=4cm] at (5,5.0) {\textbf{4. ANNONCE DU PLAN}\\[2pt] \textit{En premier lieu... Ensuite... Enfin...}\\[2pt] Grandes parties (pas les arguments détaillés).\\[2pt] \textbf{But :} Guider le lecteur (docilis).};
    % Flèches de convergence
    \draw[<->, thick, popmuted, dashed] (0.5,0.75) -- node[left, xshift=-2mm] {\textbf{Général}} (0.5,2.25);
    \draw[<->, thick, popmuted, dashed] (0.5,2.25) -- node[left, xshift=-2mm] {\textbf{Particulier}} (0.5,3.75);
    \draw[<->, thick, popmuted, dashed] (0.5,3.75) -- node[left, xshift=-2mm] {\textbf{Focalisé}} (0.5,5.0);
\end{tikzpicture}
\legende{Structure en entonnoir de l'introduction : la largeur représente le degré de généralité, la hauteur l'avancée dans le raisonnement.}
\end{popfigure}

\begin{definition}[Les 4 temps obligatoires]
Une introduction de dissertation technique au Cameroun comprend \textbf{exactement} quatre mouvements séquencés :
\begin{enumerate}
    \item \textbf{L'Accroche (Amorce) :} Une phrase ou deux (3 lignes max) qui ouvre le débat sans citer le sujet mot pour mot.
    \item \textbf{La Présentation du sujet / Contexte :} Définition des concepts clés, cadrage (pays, époque, domaine technique), rappel de l'actualité si nécessaire.
    \item \textbf{La Problématique :} La question unique à laquelle la dissertation va répondre. Elle se place \textbf{à la fin de l'amorce}, juste avant l'annonce du plan.
    \item \textbf{L'Annonce du plan :} Les 2 ou 3 grandes parties annoncées avec des connecteurs de structure (\textit{En premier lieu, Ensuite, Enfin / Pour terminer}).
\end{enumerate}
\end{definition}

\courssub{Étape 1 : L'Accroche — Typologie et Pièges}

L'accroche n'est pas un ornement : elle prouve que le sujet vous intéresse et que vous avez une culture (littéraire, technique, sociologique). Au lycée technique, on valorise les accroches \textbf{ancrées dans le réel camerounais ou technique}.

\begin{propriete}[Typologie des accroches recommandées en Terminale Tech]
\begin{itemize}
    \item \textbf{Citation d'auteur africain/francophone :} \textit{« L'écriture est une arme », Mongo Beti} (pour un sujet sur la littérature engagée) ; \textit{« Le blanc est un homme comme les autres », Ferdinand Oyono} (sujet sur le racisme/altérité).
    \item \textbf{Fait de société / Statistique officielle (MINJEC, INS, MINPROFF) :} \textit{« Selon l'INS, 68\% des jeunes camerounais de 15-35 ans sont en sous-emploi »} (sujet : « L'entrepreneuriat jeune au Cameroun »).
    \item \textbf{Paradoxe / Tension :} \textit{« Le Cameroun, « grenier de la CEMAC », importe pourtant 80\% de son blé »} (sujet : « Agriculture et sécurité alimentaire »).
    \item \textbf{Définition étymologique / technique :} \textit{« Le mot "technique" vient du grec \emph{tekhnê} : art, savoir-faire. Or, aujourd'hui... »} (sujet : « Technique et humanité »).
\end{itemize}
\end{propriete}

\begin{attention}[Les 4 accroches « tueuses » de notes (À BANNIR)]
\begin{enumerate}
    \item \textbf{« Depuis la nuit des temps... / Depuis l'aube de l'humanité... »} : Vide, faux historiquement, marque l'absence de réflexion.
    \item \textbf{« Le sujet que nous allons traiter est... »} : Vous n'êtes pas un speaker radio ; le correcteur a le sujet devant les yeux.
    \item \textbf{La citation apprise par cœur mais hors-sujet :} \textit{« L'enfer, c'est les autres » (Sartre)} pour un sujet sur \textit{« L'intelligence artificielle dans l'industrie »}.
    \item \textbf{L'accroche personnelle / « Je » :} \textit{« Moi, quand j'étais petit... »} $\to$ La dissertation est un exercice \textbf{impersonnel} (on utilise « On », « Nous », ou la voix passive).
\end{enumerate}
\end{attention}

\begin{exemplebox}[Accroches types sur un sujet technique]
\textbf{Sujet :} « L'intelligence artificielle dans l'industrie camerounaise : opportunités et défis. »

\begin{itemize}
    \item \textbf{[Statistique MINJEC/INS]} : « D'après le rapport 2023 du Ministère de l'Emploi et de la Formation Professionnelle, moins de 12\% des PME camerounaises ont intégré des solutions numériques avancées dans leur chaîne de production. »
    \item \textbf{[Paradoxe]} : « Alors que le Cameroun forme chaque année des centaines d'ingénieurs en data science, nos usines de transformation du bois ou du cacao restent largement artisanales. »
    \item \textbf{[Citation technique]} : « L'automatisation ne remplace pas l'homme, elle le libère de la pénibilité », affirmait un rapport de l'UNIDO 2022. Qu'en est-il sur le terrain camerounais ?
\end{itemize}
\end{exemplebox}

\courssub{Étape 2 : Présentation du sujet — Le mouvement entonnoir}

C'est le moment de la \cle{définition contextuelle}. Ne donnez pas la définition du dictionnaire (Larousse/Robert) mais la définition \textbf{opératoire} : ce que le terme signifie \emph{dans ce sujet précis}.

\begin{methode}[Technique du « Post-it » pour ordonner l'amorce]
Pour ne pas mélanger les étapes sur votre brouillon :
\begin{enumerate}
    \item Prenez 4 petits papiers (Post-it).
    \item \textbf{Post-it 1 (Accroche) :} Écrivez votre phrase d'ouverture (citation, chiffre, paradoxe).
    \item \textbf{Post-it 2 (Sujet) :} Écrivez les définitions des 2-3 mots-clés + le cadrage (Cameroun, secteur industriel, années 2020).
    \item \textbf{Post-it 3 (Problématique) :} Écrivez votre question unique, formulée au conditionnel ou à l'indicatif interrogatif.
    \item \textbf{Post-it 4 (Plan) :} Écrivez : « 1. [Thèse] 2. [Antithèse] 3. [Synthèse] » avec les connecteurs.
    \item \textbf{Ordre :} Alignez-les : 1 $\to$ 2 $\to$ 3 $\to$ 4. Vérifiez les liens logiques (connecteurs).
    \item \textbf{Rédaction finale :} Recopiez en un seul paragraphe fluide sur la copie.
\end{enumerate}
\end{methode}

\begin{exemplebox}[Définition contextuelle vs Définition dictionnaire]
\textbf{Sujet :} « La maintenance préventive, gage de productivité dans l'industrie agroalimentaire. »

\begin{itemize}
    \item \textbf{Mauvais (Dictionnaire) :} « La maintenance est l'action de maintenir en bon état. La productivité est le rapport production/moyens. » $\to$ Trop général, hors tension du sujet.
    \item \textbf{Bon (Contextuelle) :} « Dans le secteur agroalimentaire camerounais, soumis à des normes strictes (ISO 22000, Codex Alimentarius), la \textbf{maintenance préventive} se définit comme l'ensemble des opérations planifiées (lubrification, calibration, changement pièces d'usure) visant à éviter l'arrêt inopiné des lignes de conditionnement, garantissant ainsi la \textbf{productivité} (tonnes/jour conformes). »
\end{itemize}
\end{exemplebox}

\courssub{Étape 3 : La Problématique — Le cœur nucléaire}

C'est la \cle{seule question} à laquelle votre développement répondra. Elle doit naître \textbf{naturellement} de la présentation du sujet (d'où sa place à la fin de l'amorce).

\begin{propriete}[Critères d'une bonne problématique]
\begin{enumerate}
    \item \textbf{Unique :} Une seule phrase interrogative (ou affirmative interrogée : \textit{« Il convient de se demander dans quelle mesure... »}).
    \item \textbf{Ouverte :} Pas de réponse par « Oui » ou « Non ». Elle appelle une analyse nuancée (\textit{Dans quelle mesure / Comment / Pourquoi / Quels sont les enjeux de...}).
    \item \textbf{Issue de la tension du sujet :} Elle oppose deux termes du sujet (ex: \textit{Opportunités} vs \textit{Défis} ; \textit{Tradition} vs \textit{Modernité} ; \textit{Coût} vs \textit{Qualité}).
    \item \textbf{Recouvre tout le développement :} Les 2 ou 3 parties du plan doivent être les étapes de la réponse.
\end{enumerate}
\end{propriete}

\begin{exoresolu}[Formulation de la problématique]
\textbf{Sujet :} « L'intelligence artificielle dans l'industrie camerounaise : opportunités et défis. »

\begin{itemize}
    \item \textbf{Faible :} « L'IA est-elle une opportunité pour l'industrie ? » (Fermée, binaire).
    \item \textbf{Moyenne :} « Quels sont les avantages et les inconvénients de l'IA ? » (Catalogue, pas de tension).
    \item \textbf{Excellente :} \textbf{« Dans quelle mesure l'intégration de l'intelligence artificielle peut-elle devenir un levier de compétitivité durable pour l'industrie camerounaise, au-delà des obstacles structurels (coût, formation, énergie) qu'elle soulève ? »}
    \begin{itemize}
        \item \textit{Analyse :} « Levier de compétitivité » = Opportunités (Partie 1). « Obstacles structurels » = Défis (Partie 2). « Au-delà / Durable » = Synthèse/Prospective (Partie 3).
    \end{itemize}
\end{itemize}
\tcblower
\textbf{Solution détaillée :} La problématique excellente reprend les mots du sujet (\textit{intelligence artificielle, industrie camerounaise, opportunités, défis}) mais les articulate dans une tension dialectique (\textit{levier} vs \textit{obstacles}) et annonce implicitement la structure en 3 parties (Thèse/Antithèse/Synthèse). Elle utilise « Dans quelle mesure » (standard académique).
\end{exoresolu}

\courssub{Étape 4 : L'Annonce du Plan — La carte routière}

L'annonce ne liste pas les arguments (ex: \textit{« D'abord nous verrons que l'IA réduit les coûts, ensuite qu'elle améliore la qualité... »}$\to$ \textbf{Interdit}). Elle annonce les \textbf{grandes articulations logiques}.

\begin{methode}[Formules types pour l'annonce du plan (à adapter)]
\begin{itemize}
    \item \textbf{Plan Dialectique (Thèse / Antithèse / Synthèse) :}
    \textit{« En premier lieu, nous analyserons les potentialités offertes par l'IA pour moderniser l'outil productif. Ensuite, nous examinerons les freins structurels qui entravent son déploiement massif. Enfin, nous montrerons qu'une stratégie nationale d'accompagnement est la condition sine qua non d'une appropriation durable. »}
    \item \textbf{Plan Thématique (Causes / Conséquences / Solutions) :}
    \textit{« Nous identifierons d'abord les causes du retard technologique. Nous en mesurerons ensuite les impacts sur la compétitivité. Nous proposerons enfin les axes d'une politique industrielle volontariste. »}
    \item \textbf{Connecteurs obligatoires :} \textbf{En premier lieu / D'abord / Premièrement} $\to$ \textbf{Ensuite / Dans un second temps / Par la suite} $\to$ \textbf{Enfin / Pour terminer / Ultérieurement}.
\end{itemize}
\end{methode}

\begin{attention}[Formules proscrites (Zéro point style/structure)]
\begin{itemize}
    \item \textbf{« Je vais dire que... / Je vais expliquer... / Mon plan est le suivant : »} $\to$ Subjectif, oral, enfantin.
    \item \textbf{« Première partie, Deuxième partie, Troisième partie »} $\to$ Administratif, pas littéraire.
    \item \textbf{Dévoiler les arguments :} \textit{« Nous verrons que l'IA baisse les coûts, puis qu'elle crée du chômage... »} $\to$ Tue la curiosité du correcteur (\emph{docilis} perdu).
    \item \textbf{Oublier l'annonce :} Introduction qui s'arrête à la problématique $\to$ Note plafonnée (structure invisible).
\end{itemize}
\end{attention}

\courssub{Atelier d'amélioration : Du brouillon à la copie finale}

C'est ici que se joue la note de \emph{style} et de \emph{cohérence}. Le brouillon est le laboratoire ; la copie est l'œuvre finie.

\begin{popfigure}
\begin{tikzpicture}[scale=0.85]
    % Styles
    \tikzset{
        cell/.style={rectangle, draw=popline, thick, minimum height=1.8cm, text width=5.2cm, align=left, inner sep=4pt, font=\small},
        header/.style={rectangle, fill=popblue, text=white, font=\bfseries\small, minimum height=0.8cm, text width=5.2cm, align=center, inner sep=2pt},
        tag/.style={rectangle, fill=poporange, text=white, font=\bfseries\tiny, minimum height=0.5cm, text width=2.5cm, align=center, inner sep=1pt},
        taggood/.style={rectangle, fill=popgreen, text=white, font=\bfseries\tiny, minimum height=0.5cm, text width=2.5cm, align=center, inner sep=1pt},
        tagbad/.style={rectangle, fill=popink, text=white, font=\bfseries\tiny, minimum height=0.5cm, text width=2.5cm, align=center, inner sep=1pt},
        arrow/.style={->, thick, popmuted, shorten >=2pt, shorten <=2pt}
    }
    % En-têtes
    \node[header] (h1) at (0,0) {Version FAIBLE (Note ~0,5/4)};
    \node[header] (h2) at (6,0) {Version MOYENNE (Note ~2/4)};
    \node[header] (h3) at (12,0) {Version EXCELLENTE (Note 3,5-4/4)};
    % Contenu Faible
    \node[cell, anchor=north] (c1) at (0,-0.8) {
        \textbf{Depuis la nuit des temps}, l'homme a cherché à faciliter son travail. Aujourd'hui, l'intelligence artificielle arrive au Cameroun. \textbf{Je vais vous dire} ce que c'est. L'IA c'est des ordinateurs qui pensent. L'industrie camerounaise a des problèmes. \textbf{Est-ce que l'IA est bien ?} \textbf{Mon plan :} 1. Les avantages. 2. Les inconvénients. 3. Mon avis.
    };
    \node[tagbad, anchor=north west] at (c1.north west) {ACCROCHE : Lieux communs};
    \node[tagbad, anchor=north east] at (c1.north east) {SUJET : Pas de définition};
    \node[tagbad, anchor=south west] at (c1.south west) {PROBLÉMATIQUE : Fermée};
    \node[tagbad, anchor=south east] at (c1.south east) {PLAN : Arguments dévoilés + « Je »};
    % Contenu Moyen
    \node[cell, anchor=north] (c2) at (6,-0.8) {
        Selon le MINPOSTEL, le numérique représente 4\% du PIB. L'intelligence artificielle (IA) désigne des algorithmes d'apprentissage automatique. Dans l'industrie, elle permet la maintenance prédictive. Cependant, le coût des infrastructures freine son adoption. \textbf{Quels sont les apports et limites de l'IA pour l'industrie locale ?} \textbf{En premier lieu,} nous verrons les gains de productivité. \textbf{Ensuite,} les difficultés techniques. \textbf{Enfin,} les perspectives.
    };
    \node[tag, anchor=north west] at (c2.north west) {ACCROCHE : Statistique OK};
    \node[tag, anchor=north east] at (c2.north east) {SUJET : Définition technique};
    \node[tag, anchor=south west] at (c2.south west) {PROBLÉMATIQUE : Ouverte mais liste};
    \node[tag, anchor=south east] at (c2.south east) {PLAN : Thématique simple};
    % Contenu Excellent
    \node[cell, anchor=north] (c3) at (12,-0.8) {
        \textit{« La technique n'est pas un destin, c'est un choix »,} rappelait Jacques Ellul. Pourtant, au Cameroun, l'usine de transformation de cacao de Nkoemvone tourne encore à 60\% de sa capacité, faute de maintenance prédictive. L'intelligence artificielle, entendue ici comme l'automatisation cognitive des processus industriels, promet une rupture. \textbf{Dans quelle mesure l'IA peut-elle résorber le déficit de compétitivité structurel de notre industrie, sans aggraver la fracture numérique territoriale ?} \textbf{En premier lieu,} il convient d'analyser le potentiel de modernisation de l'outil productif. \textbf{Ensuite,} nous questionnerons la viabilité socio-économique de ce transfert technologique. \textbf{Enfin,} nous esquisserons les contours d'une souveraineté numérique industrielle.
    };
    \node[taggood, anchor=north west] at (c3.north west) {ACCROCHE : Citation + Fait local};
    \node[taggood, anchor=north east] at (c3.north east) {SUJET : Définition opératoire + Cadrage};
    \node[taggood, anchor=south west] at (c3.south west) {PROBLÉMATIQUE : Tension dialectique};
    \node[taggood, anchor=south east] at (c3.south east) {PLAN : Dialectique + Verbes d'action};
\end{tikzpicture}
\legende{Analyse comparative de trois introductions sur le même sujet technique. La version excellente respecte le formalisme (impersonnel, connecteurs, tension) et ancre le sujet dans le réel camerounais.}
\end{popfigure}

\courssub{Analyse d'introductions types (Probatoire 2022-2024)}

\begin{exemplebox}[Sujet Probatoire 2023 (Séries F/G) : « Le jeune face à l'emploi au Cameroun »]
\textbf{Introduction élève (Note 3/4) :}
\begin{quote}
« Le chômage des jeunes est une bombe à retardement », alertait le Premier Ministre en 2022. Avec un taux de sous-emploi des 15-35 ans atteignant 68\% selon l'INS (2022), l'insertion professionnelle demeure le premier défi de la politique nationale de l'emploi (PNE). L'\emph{emploi} désigne ici l'activité rémunérée stable ; la \emph{jeunesse} la tranche 15-35 ans. \textbf{Comment transformer ce dividende démographique en moteur de croissance, plutôt qu'en facteur d'instabilité ?} \textbf{En premier lieu,} nous analyserons l'inadéquation formation-emploi. \textbf{Ensuite,} nous étudierons les initiatives d'auto-emploi et leurs limites. \textbf{Enfin,} nous proposerons les leviers d'une réforme structurelle du marché du travail.
\end{quote}
\textbf{Points forts :} Accroche ancrée (PM + INS), définitions contextuelles, problématique dialectique (dividende vs instabilité), plan dialectique clair, verbes d'action (\textit{analyser, étudier, proposer}). \textbf{Axes d'amélioration :} L'accroche pourrait être plus littéraire ; la transition vers la problématique est un peu brutale.
\end{exemplebox}

\courssub{Exercice d'écriture chronométrée (15 minutes)}

\begin{exercice}[Entraînement Bac : Introduction technique]
\textbf{Consigne :} Rédigez une introduction complète (15-20 lignes) sur le sujet suivant. Respectez les 4 étapes. Chronométrez-vous (15 min max).

\textbf{Sujet :} \textit{« La maintenance préventive, gage de productivité dans l'industrie agroalimentaire camerounaise. »}

\textbf{Aide-mémoire (Brouillon) :}
\begin{itemize}
    \item \textbf{Accroche :} Chiffre MINCOMMERCE / INS sur pertes post-récolte ou arrêt usine (ex: 30\% pertes cacao/banane) OU Citation technique (AFNOR).
    \item \textbf{Sujet :} Définir \textit{maintenance préventive} (planifiée, vs curative) + \textit{productivité} (rendement/qualité/coûts) + Cadrage (Agroalimentaire CM, normes ISO 22000).
    \item \textbf{Problématique :} Tension \textit{Coût immédiat de la maintenance} vs \textit{Gain long terme / conformité export}.
    \item \textbf{Plan :} 1. Apport technique/économique (Thèse). 2. Freins financiers/organisationnels (Antithèse). 3. Stratégie intégrée GMAO / Formation (Synthèse).
\end{itemize}
\end{exercice}

\begin{corrige}[Exercice n°1 - Proposition de correction]
\textbf{Introduction type (Niveau Excellent) :}

D'après l'Organisation des Nations Unies pour l'Alimentation (FAO), les pertes post-récolte en Afrique subsaharienne atteignent 37\% de la production vivrière, soit un manque à gagner de 4 milliards de dollars annuels ; au Cameroun, l'arrêt imprévu d'une ligne de conditionnement de jus de fruit peut coûter jusqu'à 5 millions de FCFA par jour. La \textbf{maintenance préventive}, définie comme l'ensemble des interventions planifiées (inspection, remplacement programmé, analyse vibratoire) sur les équipements de transformation agroalimentaire, s'oppose à la maintenance curative, réactionnelle et coûteuse. La \textbf{productivité} ne se mesure pas seulement en tonnes/jour, mais en conformité aux normes HACCP/ISO 22000, condition \emph{sine qua non} de l'export vers la zone CEMAC et l'Union Européenne. \textbf{Dans quelle mesure la mise en place d'une politique de maintenance préventive rigoureuse constitue-t-elle le levier principal pour sécuriser la productivité et la compétitivité de l'industrie agroalimentaire camerounaise, face à la contrainte budgétaire des PME ?} \textbf{En premier lieu,} nous démontrerons que la prévention technique réduit drastiquement les arrêts non planifiés et garantit la qualité sanitaire. \textbf{Ensuite,} nous analyserons les obstacles financiers, humains et organisationnels qui expliquent la réticence des industriels locaux. \textbf{Enfin,} nous montrerons que l'adoption d'une Gestion de Maintenance Assistée par Ordinateur (GMAO) couplée à la formation continue des techniciens permet de transformer ce coût en investissement rentable.
\end{corrige}

\begin{aretenir}[Barème de notation de l'introduction (Probatoire / Baccalauréat)]
Sur une dissertation notée sur \textbf{20 points}, l'introduction pèse généralement \textbf{3 à 4 points} (selon les académies) :
\begin{center}
\begin{tabularx}{\linewidth}{|l|X|c|}\hline
\textbf{Critère} & \textbf{Indicateurs de réussite} & \textbf{Points} \\ \hline
\textbf{Accroche} & Pertinente, variée, ancrée (citation, chiffre, paradoxe), pas de lieu commun. & 1 pt \\ \hline
\textbf{Présentation Sujet / Problématique} & Termes définis \emph{in contextu}, cadrage clair, problématique unique, ouverte, dialectique, bien placée (fin amorce). & 1 pt \\ \hline
\textbf{Annonce du Plan} & Structure claire (2 ou 3 parties), connecteurs de structure, verbes d'action (analyser, examiner, montrer), pas d'arguments dévoilés, ton impersonnel. & 1 pt \\ \hline
\textbf{Qualité de la langue / Fluidité} & Zéro faute majeure (accords, syntaxe), vocabulaire précis (technique), transitions fluides entre les 4 étapes, paragraphe unique aéré. & 0,5 à 1 pt \\ \hline
\end{tabularx}
\end{center}
\textbf{Conseil :} Une introduction à 3,5/4 met le correcteur en confiance (\emph{benevolus, attentus, docilis}) pour la lecture du développement.
\end{aretenir}

\courssub{Synthèse : Check-list « Prêt à copier »}

Avant de valider votre introduction sur la copie finale, vérifiez ces 10 points (technique du \emph{self-check}) :

\begin{center}
\begin{tabularx}{\linewidth}{|c|l|X|}\hline
\# & \textbf{Point de contrôle} & \textbf{Action si NON} \\ \hline
1 & L'accroche est-elle \textbf{personnalisée} (pas « Depuis la nuit des temps ») ? & Remplacez par citation/chiffre/paradoxe. \\ \hline
2 & Les termes clés sont-ils définis \textbf{dans le contexte du sujet} (pas Larousse) ? & Réécrivez la définition opératoire. \\ \hline
3 & Le cadrage (Cameroun, secteur, période) est-il explicite ? & Ajoutez une phrase de cadrage. \\ \hline
4 & La problématique est-elle \textbf{unique, ouverte, dialectique} ? & Reformulez avec « Dans quelle mesure / Comment ». \\ \hline
5 & La problématique est-elle \textbf{juste avant l'annonce du plan} ? & Déplacez-la (c'est la charnière). \\ \hline
6 & Le plan annonce-t-il \textbf{2 ou 3 grandes parties logiques} (pas les arguments) ? & Reformulez les titres des parties. \\ \hline
7 & Les connecteurs (\textit{En premier lieu, Ensuite, Enfin}) sont-ils présents ? & Insérez-les impérativement. \\ \hline
8 & Le ton est-il \textbf{impersonnel} (pas de « Je », « Mon », « Nous allons voir ») ? & Passivisez / Utilisez « Il convient de... / Nous analyserons... ». \\ \hline
9 & L'introduction tient-elle en \textbf{un seul paragraphe} (ou deux max, sans saut de ligne intempestif) ? & Fusionnez les blocs. \\ \hline
10 & Relu à voix haute (mentalement) : \textbf{fluidité, vocabulaire précis, zéro faute} ? & Corrigez les accords, précisions lexic. \\ \hline
\end{tabularx}
\end{center}

\begin{methode}[Gestion du temps en examen (3h-4h dissertation)]
\begin{itemize}
    \item \textbf{0h00 - 0h45 :} Analyse sujet + Plan détaillé (brouillon).
    \item \textbf{0h45 - 1h00 :} \textbf{Rédaction soignée de l'introduction sur la COPIE} (utilisez la check-list). Ne la laissez pas au brouillon pour la recopier à la hâte à la fin.
    \item \textbf{1h00 - 2h30 :} Rédaction du développement (sur la copie, en sautant des lignes).
    \item \textbf{2h30 - 2h45 :} Rédaction de la conclusion (sur la copie).
    \item \textbf{2h45 - 3h00 :} Relecture globale (cohérence, fautes, lisibilité).
\end{itemize}
\textbf{Règle d'or :} L'introduction et la conclusion se rédigent \textbf{sur la copie définitive} en premier et en dernier, proprement. Le développement se construit par blocs.
\end{methode}

\begin{aretenir}[L'introduction, miroir de votre rigueur]
L'introduction n'est pas une formalité administrative : c'est la \textbf{vitrine de votre intelligence}. Elle prouve au correcteur que vous avez \textbf{compris la tension du sujet}, que vous \textbf{maîtrisez le vocabulaire technique}, que vous \textbf{structurez votre pensée} et que vous \textbf{respectez les codes de l'écrit académique}. Une introduction « propre » (structure, langue, ancrage) vaut souvent la différence entre la mention \emph{Assez Bien} et \emph{Bien}. Travaillez-la comme un artisan : mesurez, taillez, assemblez, polissez.
\end{aretenir}

\coursec{Le Développement : Architecture, Paragraphe, Citations \& Transitions}

\courssub{De l'introduction au développement : la promesse tenue}

Après avoir franchi le seuil de l'\emph{introduction} — où vous avez accroché le lecteur, posé le sujet, défini les termes, énoncé votre thèse et annoncé votre plan —, vous entrez maintenant dans le cœur de la dissertation : le \textbf{développement}. C'est ici que se joue la crédibilité de votre argumentation. L'introduction a fait une \emph{promesse} (le plan) ; le développement doit la \emph{tenir} par une démonstration rigoureuse, progressive et cohérente.

Rappelons le principe cardinal : \textbf{une dissertation n'est pas une suite d'idées jetées au fil de la plume, c'est une architecture}. Chaque étage (partie) repose sur des murs porteurs (paragraphes) eux-mêmes faits de briques solidement liées (arguments, exemples, citations). Si un mur fléchit, l'édifice s'effondre.

\begin{aretenir}[L'enjeu du développement]
Le développement doit \textbf{prouver} la thèse annoncée en introduction. Il s'organise en \textbf{deux ou trois grandes parties} (selon le plan dialectique : Thèse/Antithèse/Synthèse ou Thématique), chacune subdivisée en \textbf{paragraphes} (2 à 3 par partie). Chaque paragraphe défend \textbf{une seule idée directrice} (thèse partielle) qui s'inscrit dans la logique de la partie.
\end{aretenir}

\courssub{L'anatomie du paragraphe parfait : la règle « 1 paragraphe = 1 idée »}

\textbf{Intuition visuelle.}\par
Ouvrez n'importe quel ouvrage de référence, manuel scientifique ou article de presse de qualité : le texte est découpé en blocs séparés par une ligne blanche (l'alinéa). Ce n'est pas une fantaisie typographique, c'est le signal visuel d'une \textbf{unité de pensée}. En dissertation, cette règle est \textbf{inviolable} : \cle{1 paragraphe = 1 thèse partielle = 1 alinéa}.

\begin{propriete}[Structure interne canonique du paragraphe argumentatif]
Un paragraphe de dissertation technique suit un schéma reproductible en \textbf{5 temps} :
\begin{enumerate}
    \item \textbf{Phrase-topic (Thèse partielle) :} Première phrase, elle énonce l'idée force du paragraphe. Elle répond à la question : « Que prouve ce paragraphe ? »
    \item \textbf{Argument(s) explicatif(s) :} Déploiement logique (définition, analyse, mécanisme, cause/effet). Pourquoi cette thèse partielle est-elle recevable ?
    \item \textbf{Exemple(s) / Preuve(s) concrets :} Ancrage dans le réel (fait social, donnée technique, référence littéraire, statistique). L'exemple \emph{illustre}, il ne \emph{remplace pas} l'argument.
    \item \textbf{Citation(s) insérée(s) :} Appui d'autorité (auteur, œuvre, texte étudié). Doit être \emph{intégrée syntaxiquement} et \emph{commentée}.
    \item \textbf{Mini-synthèse / Transition :} Dernière phrase. Elle résume l'apport du paragraphe \textbf{et} ouvre logiquement sur le suivant.
\end{enumerate}
\end{propriete}

\begin{popfigure}
\begin{tikzpicture}[
    box/.style={draw=popblue, thick, rounded corners=6pt, fill=popblue!5, minimum height=1.8cm, text width=3.2cm, align=center, font=\small},
    arrow/.style={-Stealth, thick, popdark},
    link/.style={draw=poporange, thick, dashed, rounded corners=4pt, fill=poporange!5, minimum height=1.2cm, text width=2.8cm, align=center, font=\small\itshape}
]
% Main flow
\node[box, align=center] (topic) at (0,0){\textbf{1. Phrase-topic}\\(Thèse partielle)\\\emph{« L'école technique forme... »}};
\node[box, align=center] (arg1) at (0,-2.5){\textbf{2. Argument 1}\\(Explication)\\\emph{Acquisition de compétences...}};
\node[link, align=center] (ex1) at (4.5,-2.5){\textbf{Exemple 1}\\\emph{Stage en entreprise, projet tuteuré}};
\node[box, align=center] (arg2) at (0,-5){\textbf{3. Argument 2}\\(Analyse approfondie)\\\emph{Adaptation au marché...}};
\node[link, align=center] (ex2) at (4.5,-5){\textbf{Exemple 2}\\\emph{Insertion pro. élevée (données MINESEC récentes)}};
\node[box, align=center] (cit) at (0,-7.5){\textbf{4. Citation intégrée}\\\emph{« L'outil sans le savoir... »}\\(Auteur, Œuvre)};
\node[box, align=center] (trans) at (0,-10){\textbf{5. Mini-synthèse + Transition}\\\emph{Ainsi, l'école outille ; mais faut-il\\encore qu'elle émancipe...}};

% Arrows
\draw[arrow] (topic.south) -- (arg1.north);
\draw[arrow] (arg1.east) -- node[above,font=\tiny,pos=0.5] {illustre} (ex1.west);
\draw[arrow] (arg1.south) -- (arg2.north);
\draw[arrow] (arg2.east) -- node[above,font=\tiny,pos=0.5] {confirme} (ex2.west);
\draw[arrow] (arg2.south) -- (cit.north);
\draw[arrow] (cit.south) -- (trans.north);

% Braces for grouping
\draw[popdark, decorate, decoration={brace, amplitude=8pt}] (topic.north east) -- (cit.south east) node[midway, right=12pt, align=left, font=\small] {Corps du paragraphe\\ (Preuve \& Analyse)};
\draw[popdark, decorate, decoration={brace, amplitude=8pt}] (trans.north west) -- (trans.south west) node[midway, left=12pt, font=\small] {Charnière};
\end{tikzpicture}
\legende{Schéma : Anatomie d'un paragraphe parfait. Les boîtes bleues sont les étapes obligatoires (argumentation) ; les boîtes orange pointillées sont les ancrages concrets (exemples). La flèche descendante montre l'enchaînement logique ; la transition finale (charnière) relie au paragraphe suivant.}
\end{popfigure}

\begin{attention}[Piège mortel : le « catalogue » d'exemples]
Ne confondez \textbf{jamais} \emph{argument} et \emph{exemple}.
\begin{itemize}
    \item \textbf{Argument :} « La formation technique développe l'autonomie par la résolution de problèmes réels. » (Raisonnement abstrait)
    \item \textbf{Exemple :} « Lors du projet de fin d'année, l'élève conçoit un système d'irrigation solaire pour son village. » (Cas concret)
    \item \textbf{Erreur :} Enchaîner « Exemple 1, Exemple 2, Exemple 3 » sans phrase-topic ni argument explicatif. Le correcteur voit une \emph{liste de courses}, pas une démonstration.
\end{itemize}
\emph{Règle d'or :} L'exemple est la \emph{vitrine} de l'argument ; l'argument est la \emph{structure} du magasin.
\end{attention}

\courssub{La technique du « Sandwich » : maîtriser l'insertion des citations}

Une citation n'est pas un ornement : c'est une \textbf{preuve d'autorité}. Au Bac technique, elle doit provenir des textes au programme (\emph{Le Vieux Nègre et la Médaille}, textes d'anthologie, documents techniques) ou d'auteurs de référence. La laisser « nue » (sans introduction ni commentaire) est une faute grave.

\begin{methode}[La technique du « Sandwich » pour la citation — 3 étapes obligatoires]
\begin{enumerate}
    \item \textbf{Le pain du bas — INTRODUIRE (Contexte + Annonce) :} Situez la citation (Auteur, Œuvre, circonstance) et annoncez son rôle : \emph{« C'est ce que souligne Ferdinand Oyono dans \emph{Le Vieux Nègre et la Médaille} lorsqu'il écrit : »}
    \item \textbf{La garniture — CITER (Respect du texte) :} Guillemets français \og ... \fg{}, italique pour le titre de l'œuvre, respect scrupuleux de la ponctuation et de l'orthographe originales. Si citation $> 3$ lignes : \textbf{bloc indenté}, corps plus petit, \textbf{sans guillemets}.
    \item \textbf{Le pain du haut — COMMENTER (Analyse + Portée) :} Expliquez \emph{pourquoi} cette citation prouve votre argument. \emph{« Cette réplique révèle l'ironie amère du colonisé face à la médaille : l'objet symbolique devient chaîne. »}
\end{enumerate}
\emph{Jamais de citation orpheline : Introduire $\rightarrow$ Citer $\rightarrow$ Commenter.}
\end{methode}

\begin{exemplebox}[Insertion syntaxique vs bloc]
\textbf{Intégration syntaxique (courte) :} Meka « \og ne comprenait pas pourquoi on lui donnait une médaille pour avoir servi les Blancs \fg{} » (Oyono, \emph{Le Vieux Nègre et la Médaille}, p.~45), ce qui illustre l'incompréhension culturelle.

\textbf{Introduction par deux-points (moyenne) :} L'ironie transparaît dans cette réplique : \og Meka, tu as bien servi, voici ta médaille \fg{} (p.~45) : la familiarité du tutoiement masque la violence symbolique.

\textbf{Citation en bloc (longue $> 3$ lignes) :}
\begin{quote}
\itshape
Il regarda la médaille. Elle brillait au soleil. Il la tourna et la retourna dans ses mains calleuses. Elle était belle, très belle. Mais elle ne valait pas la peine qu'on l'eût tuée pour elle. (Oyono, \emph{Le Vieux Nègre et la Médaille}, p.~89)
\end{quote}
Cette description sensorielle (brillance, calleuses) oppose la valeur marchande de l'objet à la valeur humaine de la vie perdue.
\end{exemplebox}

\begin{propriete}[Critères de choix d'une citation pertinente]
\begin{itemize}
    \item \textbf{Pertinence :} Elle doit \emph{étayer} l'argument en cours, pas un autre.
    \item \textbf{Brièveté :} Préférez un extrait ciblé (1 phrase, 1 verset) à un long passage dont vous n'exploiterez que 10\%.
    \item \textbf{Identification :} Auteur + Œuvre (+ Acte/Scène/Page si théâtre/roman) \textbf{systématiques} entre parenthèses ou en note.
    \item \textbf{Variété :} Alternez citations littéraires, techniques, statistiques, proverbes.
\end{itemize}
\end{propriete}

\courssub{Les transitions : les charnières invisibles de la démonstration}

Si les paragraphes sont les wagons, les transitions sont les \textbf{attelages}. Sans eux, le train se disloque. La transition se place \textbf{à la fin du paragraphe} (dernière phrase) ou \textbf{au début du suivant} (première phrase), idéalement les deux (écho).

\begin{propriete}[Anatomie d'une transition réussie = 3 composantes]
Une transition efficace opère une \textbf{synthèse-relais} :
\begin{enumerate}
    \item \textbf{Rappel (Rétrospection) :} Reprend le mot-clé ou l'idée force du paragraphe qui s'achève.
    \item \textbf{Lien logique (Connecteur) :} Marque la relation intellectuelle (opposition, approfondissement, cause, concession...).
    \item \textbf{Annonce (Prospection) :} Présente la thèse partielle du paragraphe qui commence.
\end{enumerate}
\emph{Exemple :} « \textbf{Ainsi}, l'école technique \textbf{outille} l'élève (rappel) ; \textbf{cependant} (opposition), \textbf{la compétence technique sans conscience citoyenne reste incomplète} (annonce). »
\end{propriete}

\begin{popfigure}
\begin{tikzpicture}[
    decision/.style={draw=popblue, thick, rounded corners=8pt, fill=popblue!10, text width=3.5cm, align=center, font=\small\bfseries, minimum height=1.5cm},
    answer/.style={draw=popgreen, thick, rounded corners=6pt, fill=popgreen!10, text width=3.2cm, align=center, font=\small, minimum height=1.2cm},
    arrow/.style={-Stealth, thick, popdark}
]
\node[decision, align=center] (start) at (0,0){\textbf{QUESTION :}\\\emph{Quelle relation logique\\entre l'idée A et l'idée B ?}};
\node[answer, align=center] (opp) at (-5,-3){\textbf{OPPOSITION / CONCESSION}\\\emph{Idée B contredit ou nuance A}\\\par\vspace{2pt}\textit{Cependant, Toutefois, Pourtant, Certes... mais, En revanche}};
\node[answer, align=center] (appr) at (0,-3){\textbf{APPROFONDISSEMENT / AJOUT}\\\emph{Idée B renforce ou complète A}\\\par\vspace{2pt}\textit{Par ailleurs, De surcroît, En outre, D'ailleurs, Qui plus est}};
\node[answer, align=center] (cause) at (5,-3){\textbf{CAUSE / CONSÉQUENCE}\\\emph{Idée B explique ou résulte de A}\\\par\vspace{2pt}\textit{C'est pourquoi, Par conséquent, Dès lors, En effet, Par suite}};
\node[answer, align=center] (explic) at (-5,-6.5){\textbf{EXPLICATION / ILLUSTRATION}\\\emph{Idée B précise ou donne un ex. de A}\\\par\vspace{2pt}\textit{Autrement dit, En d'autres termes, Par exemple, Notamment}};
\node[answer, align=center] (concl) at (5,-6.5){\textbf{CONCLUSION PARTIELLE}\\\emph{Idée B tire le bilan de A}\\\par\vspace{2pt}\textit{En définitive, Au total, En somme, Bilan :}};

\draw[arrow] (start.south west) -- node[left,font=\tiny,pos=0.5] {Non / Nuance} (opp.north);
\draw[arrow] (start.south) -- node[above,font=\tiny,pos=0.5] {Oui / Plus loin} (appr.north);
\draw[arrow] (start.south east) -- node[right,font=\tiny,pos=0.5] {Pourquoi / Résultat} (cause.north);
\draw[arrow] (opp.south) -- node[left,font=\tiny,pos=0.5] {Préciser} (explic.north);
\draw[arrow] (cause.south) -- node[right,font=\tiny,pos=0.5] {Bilan} (concl.north);
\end{tikzpicture}
\legende{Arbre de décision : Comment choisir son connecteur de transition ? Partez de la relation logique réelle entre vos deux idées, pas du connecteur qui « sonne bien ».}
\end{popfigure}

\begin{exemplebox}[Transitions types selon le plan dialectique]
\textbf{Plan Thèse / Antithèse / Synthèse (sujet : « La technique libère-t-elle l'homme ? »)}

\textit{Fin partie Thèse (dernier §) :} « \textbf{En somme}, la technique libère l'homme des tâches pénibles et repousse les limites physiques ; \textbf{cependant}, cette émancipation matérielle s'accompagne de nouvelles servitudes. »

\textit{Début partie Antithèse (premier §) :} « \textbf{En effet}, la dépendance technologique crée de nouvelles aliénations : l'homme devient l'appendice de la machine. »

\textit{Fin partie Antithèse :} « \textbf{Néanmoins}, refuser la technique serait un retour en arrière illusoire ; \textbf{c'est pourquoi} il faut repenser le rapport homme/technique. »

\textit{Début partie Synthèse :} « \textbf{Dès lors}, la libération véritable réside dans la \textbf{maîtrise critique} de l'outil technique. »
\end{exemplebox}

\courssub{Gestion de l'équilibre : l'architecture symétrique}

Un développement déséquilibré (ex. : Partie I = 3 pages, Partie II = 10 lignes) trahit une pensée déséquilibrée. Au Probatoire/Bac technique, la \textbf{symétrie quantitative} reflète la \textbf{rigueur qualitative}.

\begin{exemplebox}[Repères chiffrés pour la copie Bac (double feuille, écriture standard)]
\begin{center}
\begin{tabularx}{\linewidth}{|l|X|X|}\hline
\textbf{Élément} & \textbf{Volume cible} & \textbf{Signal d'alerte} \\ \hline
Paragraphe unique & 15 -- 25 lignes & $< 8$ lignes (sous-développé) ou $> 35$ lignes (fourre-tout) \\ \hline
Grande partie (I, II, III) & 2 à 3 paragraphes $\approx$ 40--70 lignes & Déséquilibre $> 1:2$ entre parties \\ \hline
Développement total & $\approx$ 2 à 3 pages (sur 4--5 totales) & Introduction + Conclusion $\approx$ 1 page \\ \hline
\end{tabularx}
\end{center}
\emph{Astuce :} Comptez vos lignes à l'entraînement. Visez \textbf{20 lignes/paragraphe} = rythme de croisière idéal.
\end{exemplebox}

\begin{aretenir}[Check-list d'équilibre avant de rendre la copie]
\begin{itemize}
    \item [$\square$] Même nombre de paragraphes par grande partie (ex: 2/2/2 ou 3/3/3).
    \item [$\square$] Longueur comparable (pas de paragraphe « nain » ni « géant »).
    \item [$\square$] Chaque partie a sa \textbf{phrase d'appel} (début) et sa \textbf{transition de sortie} (fin).
    \item [$\square$] Les connecteurs de transition varient (pas 5 « Par ailleurs » de suite).
\end{itemize}
\end{aretenir}

\courssub{Atelier rédaction guidée : deux paragraphes types (Thèse \& Antithèse)}

Appliquons la méthode sur un sujet technique transversal, en mobilisant \emph{Le Vieux Nègre et la Médaille} d'Oyono.

\textbf{Sujet :} « La reconnaissance institutionnelle (médaille, diplôme, certification) garantit-elle la dignité du travailleur technique ? »

\textbf{Plan choisi :} I. La reconnaissance formelle : nécessaire mais insuffisante (Thèse partielle). II. La dignité véritable naît de l'autonomie et de la maîtrise (Antithèse). III. Vers une reconnaissance par les pairs et l'utilité sociale (Synthèse).

\begin{exoresolu}[Rédaction guidée : Paragraphe 1 (Partie I - Thèse) \& Paragraphe 2 (Partie II - Antithèse)]
\textbf{Énoncé :} Sujet : « La reconnaissance institutionnelle (médaille, diplôme, certification) garantit-elle la dignité du travailleur technique ? » — Plan : I. La reconnaissance formelle : nécessaire mais insuffisante. II. La dignité véritable naît de l'autonomie et de la maîtrise. III. Vers une reconnaissance par les pairs et l'utilité sociale.
\tcblower
\textbf{Paragraphe 1 — Idée directrice :} \emph{La médaille/diplôme valide une compétence technique mais ne saurait remplacer la considération humaine.}

\textit{[Phrase-topic]} La certification officielle constitue un \textbf{repère objectif} de qualification professionnelle. \textit{[Argument 1]} Elle standardise les acquis, facilite l'insertion et protège contre l'arbitraire : le diplôme est une \emph{monnaie d'échange} sur le marché du travail. \textit{[Exemple 1]} Ainsi, le Brevet de Technicien (BT) ou le Baccalauréat Technique attestent, au Cameroun, d'un niveau de compétences vérifié par le MINESEC, ouvrant droit à la fonction publique ou au secteur privé. \textit{[Citation intégrée - Sandwich]} C'est ce que suggère ironiquement Oyono lorsque le commandant remet la médaille à Meka : \og Meka, tu as bien servi, voici ta médaille \fg{} (\emph{Le Vieux Nègre et la Médaille}, p.~45) : l'institution \emph{reconnaît} le service rendu par un symbole. \textit{[Mini-synthèse + Transition]} \textbf{Toutefois}, cette reconnaissance \emph{formelle} masque souvent une \emph{invisibilité} humaine : la médaille brille sur la veste, mais l'homme reste transparent aux yeux du pouvoir.

\textbf{Paragraphe 2 — Idée directrice :} \emph{La dignité du technicien réside dans l'autonomie de jugement, non dans le parchemin.}

\textit{[Phrase-topic]} Au-delà du titre, la \textbf{dignité professionnelle} s'enracine dans la capacité à décider, innover et assumer ses choix techniques. \textit{[Argument 1]} Un technicien qui n'exécute que des consignes sans comprendre le \emph{pourquoi} ni le \emph{comment} reste un exécutant, fut-il décoré. \textit{[Exemple 1]} Dans l'atelier de maintenance, le technicien senior qui diagnostique une panne inédite par raisonnement scientifique (et non par notice) gagne le respect de ses pairs — respect que nul diplôme ne saurait décréter. \textit{[Citation intégrée]} Meka, au crépuscule de sa vie, le pressent : \og Il regarda la médaille... Elle ne valait pas la peine qu'on l'eût tuée pour elle \fg{} (Oyono, p.~89) : l'objet symbolique ne rachète pas la vie aliénée. \textit{[Mini-synthèse + Transition]} \textbf{Par conséquent}, la véritable reconnaissance ne vient pas d'en haut (hiérarchie), mais de \emph{l'intérieur} (maîtrise) et de \emph{l'extérieur} (utilité sociale perçue par la communauté).
\end{exoresolu}

\begin{methode}[Auto-évaluation de votre paragraphe (grille 5 points)]
Après rédaction, cochez :
\begin{enumerate}
    \item \textbf{Topic :} La 1ère phrase annonce-t-elle clairement l'idée du § ?
    \item \textbf{Argument :} Y a-t-il une explication logique (pas seulement un exemple) ?
    \item \textbf{Preuve :} Exemple concret + Citation \emph{sandwichée} (intégrée + commentée) ?
    \item \textbf{Transition :} La dernière phrase relie-t-elle au § suivant (rappel + lien + annonce) ?
    \item \textbf{Équilibre :} 15--25 lignes ? Connecteurs variés ? Zéro répétition ?
\end{enumerate}
\end{methode}

\courssub{Focus langue : Progression thématique, cohésion \& cohérence (Rappel d'unité)}

Le programme insiste sur la \textbf{morphosyntaxe de la progression thématique} : comment faire « couler » le sens d'une phrase à l'autre à l'intérieur du paragraphe.

\begin{propriete}[Les trois piliers de la fluidité textuelle]
\begin{itemize}
    \item \textbf{Progression thématique (Thème/Rhème) :} Le \emph{Thème} (point de départ, connu) d'une phrase devient le \emph{Rhème} (information nouvelle) de la précédente. \emph{Ex :} « \textbf{Le technicien} (Thème) diagnostique la panne. \textbf{Ce diagnostic} (Thème = Rhème antérieur) révèle une défaillance structurelle. »
    \item \textbf{Cohésion (Marqueurs de surface) :} Connecteurs logiques (\emph{donc, cependant, en effet}), reprises pronominales (\emph{il, ce qui, celui-ci}), chaines lexicales (synonymes, hyperonymes : \emph{médaille $\rightarrow$ distinction $\rightarrow$ récompense $\rightarrow$ symbole}).
    \item \textbf{Cohérence (Architecture profonde) :} Enchaînement logique des idées (définition $\rightarrow$ illustration $\rightarrow$ analyse $\rightarrow$ conséquence). Un texte cohésif mais incohérent est du « charabia bien lié ».
\end{itemize}
\end{propriete}

\begin{savaistu}[Étymologie : « Faire une transition »]
Le mot \emph{transition} vient du latin \emph{transire} (\emph{trans-} « au-delà, par-dessus » + \emph{ire} « aller »). Littéralement : \textbf{« passer par-dessus »}. En dissertation, faire une transition, c'est \emph{franchir une étape sans la brûler} : on pose le pied sur la rive d'en face en gardant l'appui sur la rive d'ici. Brûler l'étape (pas de transition) = sauter le fossé = risque de chute pour le lecteur.
\end{savaistu}

\courssub{Exercice d'intégration \& Remédiation ciblée}

\begin{exercice}[Entraînement : « Chirurgie du paragraphe »]
Voici un paragraphe \textbf{défectueux} (sujet : « Technique et responsabilité »). Identifiez les 5 erreurs majeures, puis réécrivez-le en appliquant la structure canonique + technique du Sandwich + transition.

\begin{quote}
\itshape
La technique est dangereuse. Par exemple Tchernobyl. Aussi les voitures polluent. Comme dit Heidegger « La technique n'est pas un outil ». Donc il faut faire attention. Par ailleurs le téléphone portable rend accro. En conclusion la technique c'est pas bon.
\end{quote}

\textbf{Consignes :}
\begin{enumerate}
    \item Repérez : absence de phrase-topic, catalogue d'exemples sans arguments, citations orphelines, connecteurs inadaptés, conclusion abusive.
    \item Réécrivez en 20 lignes : Thèse partielle claire $\rightarrow$ Argument $\rightarrow$ Exemple (Tchernobyl \emph{ou} portable) $\rightarrow$ Citation (Sandwich) $\rightarrow$ Transition vers § suivant (ex: « Cependant, rejeter la technique... »).
\end{enumerate}
\end{exercice}

\begin{corrige}[Exercice : Chirurgie du paragraphe]
\textbf{Erreurs diagnostiquées :}
\begin{enumerate}
    \item Pas de thèse partielle (phrase 1 trop vague).
    \item « Par exemple Tchernobyl » : exemple nu, sans argument explicatif (catalogue).
    \item « Aussi les voitures polluent » : connecteur \emph{aussi} (conséquence) inadapté pour un 2ème exemple ; pas de lien logique.
    \item Citation « nue », mal citée (pas de référence), non commentée.
    \item « En conclusion » au milieu du développement ; thèse finale (« c'est pas bon ») contredit le sujet (nuance absente).
\end{enumerate}

\textbf{Proposition de réécriture (structure canonique) :}
\emph{« Si la technique démultiplie la puissance d'agir, elle engage proportionnellement la responsabilité de ses concepteurs et utilisateurs. [Topic] En effet, tout système technique porte en lui un risque systémique que la seule compétence technique ne saurait maîtriser sans éthique. [Argument] La catastrophe de Tchernobyl (1986) l'illustre tragiquement : non pas la défaillance d'un réacteur, mais l'enchaînement de décisions bureaucratiques privilégiant le rendement sur la sûreté. [Exemple commenté] Rabelais avertit à juste titre : \og Science sans conscience n'est que ruine de l'âme \fg{} (\emph{Pantagruel}, Chap. 8) : la puissance technique impose sa propre logique que l'homme doit penser, pas seulement subir. [Citation sandwichée] \textbf{Cependant}, cette responsabilité ne saurait mener au refus technophobe ; elle appelle au contraire une \textbf{maîtrise critique} que seule l'éducation technique humaniste peut fonder. [Transition vers § Antithèse : l'éducation technique comme rempart] »}
\end{corrige}

\begin{attention}[Pièges de dernière minute (Check-list « Nuit avant le Bac »)]
\begin{itemize}
    \item [$\square$] \textbf{Saut de ligne visible} entre chaque paragraphe (alinea = 1 cm ou ligne blanche).
    \item [$\square$] \textbf{Pas de paragraphe unique} de 2 pages (découpez !).
    \item [$\square$] \textbf{Citations :} Guillemets \og \fg{} + italique œuvre + (Auteur, p.) + \textbf{commentaire}.
    \item [$\square$] \textbf{Transitions :} Au moins 1 mot de liaison logique \emph{par passage} d'un § à l'autre.
    \item [$\square$] \textbf{Variété lexicale :} Évitez les « aussi / donc / par ailleurs » en boucle. Piochez dans l'arbre de décision.
\end{itemize}
\end{attention}

\begin{aretenir}[Synthèse visuelle : Le paragraphe « Bac Ready »]
\begin{center}
\begin{tabularx}{\linewidth}{|l|X|}\hline
\textbf{Zone} & \textbf{Contenu type (exemple : Partie I, §1)} \\ \hline
L1--L2 & \textbf{Phrase-topic :} « La certification technique valide les acquis mais ne garantit pas l'autonomie. » \\ \hline
L3--L8 & \textbf{Argument + Exemple :} Explication du rôle du diplôme (MINESEC) + Exemple concret (BT Électrotechnique, stage). \\ \hline
L9--L14 & \textbf{Citation Sandwich :} Intro $\rightarrow$ \og Meka, tu as bien servi...\fg{} (Oyono, p.45) $\rightarrow$ Analyse : ironie institutionnelle. \\ \hline
L15--L18 & \textbf{Approfondissement :} Nuance (diplôme nécessaire $\neq$ suffisant), ouverture sur compétences transverses. \\ \hline
L19--L20 & \textbf{Transition :} « \textbf{Néanmoins}, c'est dans l'exercice réel du métier que se forge la légitimité... » (Annonce §2). \\ \hline
\end{tabularx}
\end{center}
\emph{Objectif :} 20 lignes, 5 mouvements, 1 thèse partielle, 0 citation nue, 1 transition sortante.
\end{aretenir}

\textbf{Transition vers la Section 4.}\par
Nous avons désormais les \textbf{outils de construction} (paragraphe, citation, transition, équilibre) pour bâtir le corps de la dissertation. Mais une dissertation technique au Cameroun ne saurait être complète sans une \textbf{étude intégrée des textes au programme}. La Section 4 mobilisera \emph{Le Vieux Nègre et la Médaille} (Lectures méthodiques 2 \& 3) et le \textbf{texte explicatif} pour nourrir vos arguments d'une culture littéraire et technique ancrée dans le contexte africain. C'est là que vos paragraphes gagneront en \emph{profondeur} et en \emph{authenticité}.

\coursec{Étude Intégrée : Le Vieux Nègre et la Médaille (Lectures 2 \& 3) \& Texte Explicatif}

Dans la section précédente, nous avons appris à construire le développement d'une dissertation : un paragraphe bien articulé, des citations insérées avec soin, des transitions qui guident le lecteur. Mais pour écrire un bon paragraphe, il faut disposer d'un \emph{exemple littéraire solide}, analysé en profondeur. C'est précisément l'objet de cette section : faire du roman de Ferdinand Oyono, \emph{Le Vieux Nègre et la Médaille} (1956), un \cle{réservoir d'arguments} et d'illustrations directement mobilisables le jour du Baccalauréat technique. Nous y croiserons également un savoir de langue essentiel : le \cle{texte explicatif} et la \cle{progression thématique}.

\courssub{Rappel contextuel : Oyono et le Cameroun colonial}

\begin{definition}[L'œuvre et son contexte]
\emph{Le Vieux Nègre et la Médaille} est un roman camerounais publié en 1956 par Ferdinand Oyono. L'action se déroule dans le Cameroun colonial des années 1950, sous administration française. Le héros, \cle{Meka}, vieillard du village de Doum, reçoit une médaille pour avoir « donné » ses terres à la mission catholique et ses deux fils à la guerre. Le roman dénonce, par l'\cle{ironie}, l'idéologie de l'\emph{assimilation} et l'hypocrisie de l'administration coloniale.
\end{definition}

L'intuition à retenir est simple : Oyono montre un homme qui croit être honoré par le colonisateur, et qui découvre, au prix de sa vie, que cet honneur est un \emph{leurre}. La médaille, symbole de reconnaissance, devient le symbole de la \cle{trahison} et de la \cle{vanité} des promesses coloniales.

\begin{savaistu}[Un écrivain engagé devenu ministre]
Ferdinand Oyono (1929-2010) n'a pas été seulement romancier : il a été diplomate (ambassadeur du Cameroun) et ministre. Publié en 1956, soit quatre ans avant l'indépendance du Cameroun (1960), \emph{Le Vieux Nègre et la Médaille} a préfiguré la décolonisation en montrant, de l'intérieur d'un village africain, l'effondrement du mythe colonial de la « mission civilisatrice ».
\end{savaistu}

\courssub{Lecture méthodique n°2 : la cérémonie et le discours du Commandant}

\textbf{Situation de l'extrait.} Le 14 juillet, à Doum, le Commandant de cercle décore Meka devant la foule assemblée. Son discours est un modèle de \cle{texte explicatif et argumentatif} de propagande : il explique et justifie « l'œuvre de la France ».

\textbf{Observation guidée.} Lisons le discours comme des enquêteurs : qui parle ? à qui ? dans quel but ? Nous constatons que le Commandant célèbre Meka comme un « ami de la France », qu'il accumule les éloges, mais que le narrateur, en filigrane, laisse apparaître le grotesque de la scène : un vieillard épuisé, debout sous le soleil, décoré pour s'être dépossédé lui-même. C'est le mécanisme de l'\cle{ironie dramatique} : le lecteur comprend ce que le personnage ne comprend pas encore.

\begin{definition}[Le texte explicatif]
Le \cle{texte explicatif} est un texte à dominante argumentative qui vise à faire comprendre une notion, un phénomène ou une thèse. Ses procédés typiques sont : la \textbf{définition}, la \textbf{classification}, l'\textbf{illustration} (exemples), la \textbf{reformulation} (marqueurs : \emph{autrement dit}, \emph{c'est-à-dire}, \emph{en effet}) et la \textbf{comparaison}.
\end{definition}

\begin{center}
\begin{tabularx}{\linewidth}{|l|X|}\hline
\rowcolor{popblueL} \textbf{Critère} & \textbf{Analyse de l'extrait n°2 (discours du Commandant)} \\ \hline
Thèse & L'ordre colonial est légitime et généreux : la France « récompense » ses fidèles. \\ \hline
Arguments & Arguments d'\cle{autorité} (l'État, la médaille, la cérémonie officielle) ; argument de la « gratitude » (Meka a donné ses terres et ses fils). \\ \hline
Procédés explicatifs & Éloge amplifié, généralisations (« les évolués », « l'œuvre civilisatrice »), reformulations flatteuses. \\ \hline
Connecteurs & Connecteurs d'addition et de justification (\emph{car}, \emph{en effet}, \emph{c'est pourquoi}). \\ \hline
Ton & Solennel et paternaliste en surface ; \cle{ironique} en profondeur (registre satirique). \\ \hline
\end{tabularx}
\end{center}

\begin{attention}[Ne pas confondre les points de vue]
Piège classique au Bac : attribuer à Oyono les propos du Commandant. Il faut distinguer trois voix : celle de l'\textbf{auteur} (Oyono, critique envers la colonisation), celle du \textbf{personnage} (Meka, naïf puis lucide), et celle du \textbf{Commandant} (porte-parole de la propagande coloniale). Dans une dissertation, écrivez : « le Commandant déclare... », puis « le narrateur ironise... », jamais « Oyono dit que la France est généreuse ».
\end{attention}

\courssub{Lecture méthodique n°3 : la mort de Meka et la progression thématique}

\textbf{Situation de l'extrait.} Après la cérémonie, Meka, ivre, est battu et emprisonné par les gendarmes qu'il croyait ses amis. Libéré, humilié, il meurt peu après, pendant la tornade. Le récit bascule du comique au \cle{tragique}.

\textbf{Démarche d'investigation.} Pour comprendre comment le texte « avance », soulignons les \emph{débuts de phrases} (les thèmes) et les \emph{informations nouvelles} (les rhèmes). Nous observons deux mouvements simultanés : le nom « Meka » revient sans cesse en tête de phrase (le texte reste centré sur lui), tandis que l'information circule : la médaille apporte l'honneur, l'honneur tourne à la honte, la honte conduit à la mort.

\begin{definition}[Les trois types de progression thématique]
\begin{itemize}
\item \cle{Progression thématique constante} : le même thème (souvent le sujet grammatical) est répété en début de phrases successives. Exemple : « \textbf{Meka} se leva. \textbf{Meka} avança. \textbf{Il} salua la foule. »
\item \cle{Progression en chaîne} : le rhème (l'information nouvelle) d'une phrase devient le thème de la phrase suivante. Exemple : « Meka reçut une \textbf{médaille}. Cette \textbf{médaille} lui attira des envieux. »
\item \cle{Progression dérivée} : des thèmes différents dérivent d'un même hyper-thème. Exemple : hyper-thème « la cérémonie » $\rightarrow$ « le discours », « la foule », « les gendarmes ».
\end{itemize}
\end{definition}

\begin{popfigure}
\begin{center}
\begin{tikzpicture}[font=\small, >=Stealth]
% Ligne du thème constant
\draw[very thick, popblue] (0,0) -- (12,0);
\foreach \x in {0.5,2.5,4.5,6.5,8.5,10.5}{
  \node[circle, fill=popblue, inner sep=1.6pt] at (\x,0) {};
}
\node[above, popblue, font=\small\bfseries] at (0.5,0.15) {Meka};
\node[above, popblue] at (2.5,0.15) {Il};
\node[above, popblue] at (4.5,0.15) {Le vieillard};
\node[above, popblue] at (6.5,0.15) {Il};
\node[above, popblue] at (8.5,0.15) {Meka};
\node[above, popblue] at (10.5,0.15) {Son corps};
\node[popblue, font=\small\itshape] at (11.6,0.35) {thème constant};
% Chaîne thématique
\node[draw=poporange, rounded corners, fill=poporangeL, inner sep=4pt] (h) at (2,-1.8) {l'honneur};
\node[draw=poporange, rounded corners, fill=poporangeL, inner sep=4pt] (d) at (6,-1.8) {la honte};
\node[draw=poporange, rounded corners, fill=poporangeL, inner sep=4pt] (m) at (10,-1.8) {la mort};
\draw[->, very thick, poporange] (h) -- (d);
\draw[->, very thick, poporange] (d) -- (m);
\node[poporange, font=\small\itshape] at (6,-2.6) {progression en chaîne : le rhème devient thème};
% Liens
\draw[dashed, popmuted] (2.5,0) -- (h);
\draw[dashed, popmuted] (6.5,0) -- (d);
\draw[dashed, popmuted] (10.5,0) -- (m);
\end{tikzpicture}
\end{center}
\legende{Progression thématique de l'extrait n°3 : en haut, le thème constant « Meka » (reprises nominales et pronominales) ; en bas, la chaîne énonciative honneur $\rightarrow$ honte $\rightarrow$ mort qui donne sa dynamique tragique au passage.}
\end{popfigure}

Ce double mouvement est remarquable : la \emph{stabilité} du thème (Meka ne change pas, il subit) contraste avec la \emph{dégradation} des rhèmes (de l'honneur à la mort). C'est exactement le genre d'observation fine qui valorise une copie de dissertation.

\courssub{Mobilisation pour la dissertation : trois figures de l'argument}

\begin{aretenir}[Trois lectures de Meka utilisables au Bac]
\begin{itemize}
\item \textbf{Meka, figure de l'aliéné volontaire} : il a intériorisé les valeurs du colonisateur au point de donner ses terres et ses fils ; il illustre l'\emph{aliénation culturelle}.
\item \textbf{La médaille, symbole vide} : objet de « toc », elle promet l'intégration et n'apporte que l'humiliation ; elle illustre l'\emph{illusion}.
\item \textbf{Le vieux nègre, lucidité tardive} : sa révolte finale et sa mort signifient que la prise de conscience, même tardive, révèle la vérité du système colonial.
\end{itemize}
\end{aretenir}

\begin{exemplebox}[Un exemple chiffré pour réviser]
L'œuvre compte environ 120 pages : sa lecture intégrale est rapide et faisable en une semaine. Les trois citations les plus rentables au Bac portent sur : (1) la médaille, objet de « toc » ; (2) le discours du Commandant lors de la cérémonie du 14 juillet ; (3) la mort de Meka pendant la tornade. Apprenez une phrase courte pour chacune.
\end{exemplebox}

\textbf{Exploitation argumentative.} Supposons deux sujets de dissertation :
\begin{itemize}
\item \emph{« Le progrès n'est-il qu'une illusion ? »} : Meka croit progresser socialement grâce à la médaille ; l'épisode de la prison prouve que ce « progrès » était un mirage. La médaille devient l'illustration parfaite de la thèse.
\item \emph{« Le choc des cultures »} : la cérémonie du 14 juillet met face à face deux mondes (la fête républicaine française et le village de Doum) ; l'incompréhension réciproque débouche sur la violence. Meka incarne l'homme écrasé entre deux cultures.
\end{itemize}

\begin{methode}[Insérer une citation d'Oyono en trois étapes]
\begin{enumerate}
\item \textbf{Contextualiser} : qui parle ? à qui ? dans quelle situation ? (« Lors de la cérémonie du 14 juillet, le Commandant déclare devant la foule... »)
\item \textbf{Citer exactement} : phrase courte, entre guillemets, fidèle au texte.
\item \textbf{Analyser} : dégager l'ironie ou la portée critique (« Cet éloge, que le narrateur laisse entendre comme grotesque, révèle l'hypocrisie de... »).
\end{enumerate}
\end{methode}

\begin{exoresolu}[Transformer un passage narratif en paragraphe explicatif argumentatif]
Énoncé : à partir de l'épisode de la médaille, rédigez un paragraphe argumentatif (thèse, exemple, analyse) sur le sujet : \emph{« Les honneurs peuvent-ils tromper ? »}
\tcblower
\textbf{Solution détaillée.} \\
\emph{Thèse} : Les honneurs officiels peuvent masquer une domination et tromper celui qui les reçoit. \emph{Exemple contextualisé} : dans \emph{Le Vieux Nègre et la Médaille} de Ferdinand Oyono, le vieux Meka est décoré par le Commandant, le 14 juillet, pour avoir donné ses terres à la mission et ses fils à la guerre. \emph{Citation} : la médaille, objet de « toc », brille un instant avant de perdre tout éclat. \emph{Analyse} : cet honneur n'est qu'un instrument de propagande ; dès le soir même, Meka est battu et emprisonné par ceux-là mêmes qui l'ont décoré, autrement dit la reconnaissance coloniale était un leurre. \emph{Conclusion du paragraphe} : l'exemple d'Oyono démontre qu'un honneur accordé par un pouvoir injuste n'honore pas le récipiendaire, il l'enchaîne. \\
Remarquez la mécanique : le passage narratif (la scène de la cérémonie) est devenu un \emph{argument} grâce à la contextualisation, à la citation exacte et à l'analyse ironique.
\end{exoresolu}

\begin{exercice}[Entraînement Bac]
\begin{enumerate}
\item Relevez dans l'extrait du discours du Commandant deux procédés du texte explicatif (définition, reformulation, illustration) et justifiez leur fonction persuasive.
\item Identifiez, dans un paragraphe de l'extrait n°3, le type de progression thématique dominante et expliquez son effet.
\item Rédigez un paragraphe argumentatif (8 à 10 lignes) utilisant la mort de Meka pour illustrer le sujet : \emph{« Faut-il parfois souffrir pour comprendre ? »}
\end{enumerate}
\end{exercice}

\begin{corrige}[Exercice : éléments de correction]
\begin{enumerate}
\item La reformulation (\emph{c'est-à-dire}, \emph{autrement dit}) permet au Commandant d'habiller la dépossession en générosité ; l'illustration (l'exemple de Meka présenté comme modèle) transforme un cas particulier en argument d'autorité.
\item Progression en chaîne dominante : le rhème « médaille » devient thème, puis « honneur » engendre « honte » ; effet : le lecteur suit mécaniquement la dégradation du destin de Meka, ce qui produit l'émotion tragique.
\item Attendu : thèse annoncée, contextualisation (après la prison, humilié, Meka meurt pendant la tornade), citation courte, analyse (la souffrance ouvre les yeux : sa révolte finale est une lucidité), phrase de synthèse.
\end{enumerate}
\end{corrige}

\courssub{Compte rendu écrit : fiche de révision synthétique}

Pour conclure cette étude intégrée, voici la fiche de révision à recopier et à compléter ; elle constitue votre \cle{compte rendu} des deux lectures méthodiques.

\begin{center}
\begin{tabularx}{\linewidth}{|l|X|}\hline
\rowcolor{popgreenL} \textbf{Rubrique} & \textbf{Contenu à retenir} \\ \hline
Thèmes majeurs & Colonisation et assimilation ; illusion et trahison ; religion et hypocrisie ; vieillesse et lucidité. \\ \hline
Personnages & Meka (héros tragique, naïf puis lucide) ; le Commandant (propagande) ; Kelara et Engamba (les siens, révolte) ; le père Vandermayer (mission). \\ \hline
Citations-clés & La médaille de « toc » ; le discours du 14 juillet ; la mort de Meka sous la tornade. \\ \hline
Portée sociologique & Le roman dénonce l'aliénation d'une génération africaine séduite puis brisée par le mythe colonial ; il annonce l'indépendance de 1960. \\ \hline
Usage en dissertation & Exemple mobilisable pour les sujets sur l'illusion, le progrès, le choc des cultures, la critique sociale du roman. \\ \hline
\end{tabularx}
\end{center}

\begin{attention}[Dernière mise en garde avant la rédaction]
Deux erreurs coûtent cher : (1) \emph{raconter} le roman au lieu d'\emph{argumenter} — la dissertation n'est pas un résumé ; (2) citer de mémoire approximative — une citation déformée discrédite toute la copie. Préférez une citation courte et exacte à une longue citation incertaine.
\end{attention}

Nous disposons désormais d'un corpus analysé, de citations contrôlées et d'une méthode d'insertion. Il reste à apprendre, dans la section suivante, à refermer la dissertation par une conclusion qui synthétise sans répéter et qui ouvre sans divaguer.

\coursec{La Conclusion : Synthèse, Bilan \& Ouverture (Production \& Amélioration)}

Dans la section précédente, notre lecture méthodique du \emph{Vieux Nègre et la médaille} de Ferdinand Oyono nous a montré comment un texte explicatif construit sa démonstration jusqu'à son terme : Meka, après avoir reçu sa médaille, comprend l'illusion de la reconnaissance coloniale, et le roman se referme sur une prise de conscience qui éclaire rétroactivement tout le récit. Il en va de même de la dissertation : la \cle{conclusion} n'est pas une simple formalité, c'est le moment où tout le parcours argumentatif trouve son sens. Un correcteur du Probatoire ou du Baccalauréat technique lit la conclusion avec une attention particulière : c'est la dernière impression qu'il garde de votre copie, et c'est elle qui confirme — ou ruine — la cohérence de l'ensemble.

\courssub{La structure en trois temps : le miroir de l'introduction}

\begin{definition}[La conclusion]
La \cle{conclusion} est la dernière partie de la dissertation. Elle remplit trois fonctions inséparables : \textbf{répondre} explicitement à la problématique posée dans l'introduction, \textbf{bilanter} synthétiquement le chemin parcouru dans le développement, et \textbf{ouvrir} la réflexion vers un horizon plus large. On dit qu'elle est le \emph{miroir} de l'introduction : l'introduction allait du général vers le sujet (entonnement), la conclusion va du sujet vers le général (élargissement).
\end{definition}

Observez le schéma suivant : la dissertation entière a une forme de \emph{sablier}. L'introduction part du large (amorce, contexte) et se resserre sur la problématique ; le développement est la partie étroite et structurée ; la conclusion s'élargit à nouveau jusqu'à l'ouverture.

\begin{popfigure}
\begin{center}
\begin{tikzpicture}[scale=1]
% Sablier : contour
\fill[popblueL] (0,4) -- (4,4) -- (2.3,2) -- (4,0) -- (0,0) -- (1.7,2) -- cycle;
\draw[popdark,thick] (0,4) -- (4,4) -- (2.3,2) -- (4,0) -- (0,0) -- (1.7,2) -- cycle;
% Étiquettes
\node[popdark,font=\bfseries] at (2,3.4) {INTRODUCTION};
\node[popdark,font=\small] at (2,2.85) {(large $\rightarrow$ problématique)};
\node[popdark,font=\bfseries] at (2,1.55) {DÉVELOPPEMENT};
\node[popdark,font=\small] at (2,1.1) {(étroit, structuré : I, II, III)};
\node[popdark,font=\bfseries] at (2,0.45) {CONCLUSION};
% Flèches de mouvement
\draw[-{Stealth},poporange,very thick] (4.5,3.8) -- (4.5,2.6) node[midway,right,font=\small\itshape]{resserrement};
\draw[-{Stealth},popgreen,very thick] (-0.5,0.2) -- (-0.5,1.4) node[midway,left,font=\small\itshape]{élargissement};
% Ouverture
\draw[-{Stealth},poppurple,very thick] (2,-0.15) -- (2,-0.9);
\node[poppurple,font=\bfseries] at (2,-1.2) {OUVERTURE (vers le large)};
\end{tikzpicture}
\end{center}
\legende{Structure en sablier de la dissertation : l'introduction resserre le propos vers la problématique, le développement démontre, la conclusion ré-élargit jusqu'à l'ouverture.}
\end{popfigure}

Détaillons les trois temps obligatoires.

\textbf{Premier temps : reprise de la problématique et réponse directe.} On reformule la question directrice du sujet (sans la recopier mot pour mot) et on y apporte une réponse claire, qui est votre \cle{thèse finale}. Exemple : « Nous nous sommes demandé si la médaille offerte à Meka constituait une véritable reconnaissance. Il apparaît, au terme de notre analyse, que cette distinction n'est qu'un instrument de propagande coloniale. »

\textbf{Deuxième temps : bilan synthétique des grandes parties.} On résume en deux ou trois phrases les idées fortes de chaque partie du développement, sans introduire \emph{aucun} argument nouveau.

\textbf{Troisième temps : l'ouverture.} On déplace le regard vers une autre œuvre, une autre époque, une autre discipline ou une perspective d'avenir, en gardant un lien logique avec le sujet.

\courssub{Rédiger la thèse finale : nuancer et modaliser}

La thèse finale ne doit être « ni tout blanc, ni tout noir ». La pensée d'un candidat de Terminale technique se reconnaît à sa capacité de \cle{nuance}. Pour cela, on utilise des \cle{modalisateurs}, c'est-à-dire des expressions qui introduisent la thèse avec mesure ou avec force contrôlée :

\begin{itemize}
\item « \emph{Il apparaît que...} » (constat raisonné) ;
\item « \emph{Force est de constater que...} » (évidence admise) ;
\item « \emph{En définitive...} », « \emph{Au terme de cette réflexion...} » (bilan) ;
\item « \emph{Si... , il n'en demeure pas moins que...} » (concession puis affirmation).
\end{itemize}

\begin{exemplebox}[Thèse radicale contre thèse nuancée]
Sujet : « Le progrès technique détruit-il les traditions ? »
\begin{itemize}
\item Thèse radicale (faible) : « Le progrès détruit totalement les traditions. »
\item Thèse nuancée (forte) : « Si le progrès technique érode certaines pratiques traditionnelles, il apparaît qu'il peut aussi les préserver et les diffuser, à condition d'être mis au service des communautés. »
\end{itemize}
\end{exemplebox}

\begin{attention}[Bannir le « Je pense que »]
En conclusion, le « \emph{Je pense que} » est interdit : il affaiblit la démonstration en la ramenant à une opinion personnelle. L'argumentation impose l'évidence par des tournures impersonnelles : « \emph{Il ressort de cette étude que...} », « \emph{Il est donc établi que...} ». Votre pensée s'exprime à travers la démonstration, non à travers le « je ».
\end{attention}

\courssub{Les techniques de synthèse : le bilan sans répétition}

Le bilan n'est pas un résumé mot pour mot de vos parties : c'est un \emph{regroupement par mots-clés}. Chaque partie du développement est condensée en une idée-force, et ces idées sont reliées par des \cle{connecteurs de conclusion} : « \emph{Au total...} », « \emph{En somme...} », « \emph{Pour conclure...} », « \emph{Ainsi...} », « \emph{Dès lors...} ».

\begin{methode}[La technique du « Brouillon inversé »]
Pour construire un bilan fidèle sans recopier son développement :
\begin{enumerate}
\item Relire son développement partie par partie.
\item Sur le brouillon, noter en une ligne l'\textbf{idée forte} de chaque partie (un mot-clé + un verbe).
\item Transformer chaque note en une phrase complète du bilan.
\item Relier ces phrases par un connecteur de conclusion.
\item Vérifier qu'aucun argument nouveau ne s'est glissé dans le bilan.
\end{enumerate}
\end{methode}

\begin{exemplebox}[Application du brouillon inversé]
Développement : I. La médaille comme récompense apparente ; II. La médaille comme instrument de domination ; III. La prise de conscience de Meka.\\
Notes du brouillon : « récompense = façade » / « domination = propagande » / « Meka = lucidité ».\\
Bilan rédigé : « En somme, la médaille n'apparaît d'abord que comme une façade honorifique ; elle se révèle ensuite un instrument de propagande coloniale ; elle devient enfin, grâce à la lucidité de Meka, le symbole d'une dignité retrouvée. »
\end{exemplebox}

\courssub{L'ouverture : déplacer le regard sans l'éteindre}

\begin{definition}[L'ouverture]
L'\cle{ouverture} est la dernière phrase (ou les deux dernières phrases) de la conclusion. Elle n'est pas une « conclusion bis » : elle \emph{déplace le regard} — dans l'espace (un autre pays), dans le temps (une autre époque, l'avenir) ou dans la discipline (littérature, histoire, sociologie, technique) — sans éteindre la réflexion. Elle doit garder un \textbf{lien logique} avec le sujet.
\end{definition}

Les critères de pertinence d'une ouverture sont au nombre de trois : un \textbf{lien logique} explicite avec la thèse défendue, un \textbf{élargissement} réel (on apporte un éclairage nouveau, on ne répète pas), et l'absence de \textbf{hors-sujet} (on ne part pas dans une direction sans rapport avec la problématique).

\begin{popfigure}
\begin{center}
\begin{tikzpicture}[mindmap, grow cyclic,
  every node/.style={concept, font=\small},
  root concept/.append style={concept color=popblue, font=\bfseries},
  level 1/.append style={level distance=3.4cm, sibling angle=72},
  level 2/.append style={level distance=2.2cm, sibling angle=45, font=\scriptsize}]
\node[root concept]{Types d'ouvertures}
  child[concept color=poporange]{ node {Littéraire}
    child { node {\emph{Une vie de boy} (Oyono)} }
    child { node {\emph{Le Pauvre Christ de Bomba} (Beti)} } }
  child[concept color=popgreen]{ node {Historique}
    child { node {Indépendances africaines (1960)} }
    child { node {Décolonisation et néocolonialisme} } }
  child[concept color=poppurple]{ node {Sociologique}
    child { node {Urbanisation de Yaoundé et Douala} }
    child { node {Exode rural des jeunes} } }
  child[concept color=poppink]{ node {Technique / Professionnel}
    child { node {Impact de l'IA sur les métiers} }
    child { node {Innovation durable et éthique} } }
  child[concept color=popgold]{ node {Philosophique}
    child { node {Qu'est-ce que la dignité ?} }
    child { node {Liberté et apparences} } };
\end{tikzpicture}
\end{center}
\legende{Carte mentale des types d'ouvertures classés par domaine, avec des exemples ancrés dans le contexte camerounais et africain.}
\end{popfigure}

\begin{savaistu}[L'ouverture attendue au Baccalauréat technique]
Au Bac technique, les sujets liés aux filières professionnelles (Maintenance, Gestion, Industries) attendent souvent une ouverture sur l'\textbf{éthique professionnelle} ou l'\textbf{innovation durable}. Par exemple, pour un sujet sur le travail, une ouverture sur la place de la technologie dans l'atelier de demain, ou sur la responsabilité du technicien face à l'environnement, est toujours bien accueillie par le correcteur.
\end{savaistu}

\courssub{Atelier d'amélioration : d'une conclusion bâclée à une conclusion maîtrisée}

\begin{exoresolu}[Transformer une conclusion faible]
\textbf{Sujet :} « La médaille de Meka est-elle une véritable récompense ? »\\
\textbf{Conclusion bâclée :} « Voilà. Je pense que la médaille n'est pas une vraie récompense. L'avenir nous le dira. »\\
\textbf{Travail demandé :} réécrire cette conclusion en trois temps.
\tcblower
\textbf{Diagnostic :} la conclusion tient en une ligne, contient un « je pense », aucun bilan, et une phrase creuse (« L'avenir nous le dira »).\\
\textbf{Conclusion améliorée :} « Nous nous sommes demandé si la médaille remise à Meka constituait une véritable récompense. Il ressort de notre analyse que cette distinction, loin d'honorer le vieil homme, sert les intérêts de l'administration coloniale. En somme, la médaille n'est d'abord qu'une façade honorifique ; elle se révèle ensuite un instrument de propagande ; elle devient enfin, par la lucidité de Meka, le signe d'une dignité reconquise. Au-delà du roman d'Oyono, on peut s'interroger sur les décorations et distinctions distribuées aujourd'hui : toutes expriment-elles une reconnaissance sincère, ou certaines servent-elles encore à masquer des rapports de domination ? »\\
Cette version reprend la problématique, formule une thèse modalisée, bilan fidèle en trois idées-forces, et ouverture pertinente (lien logique : les distinctions ; élargissement : l'époque actuelle).
\end{exoresolu}

\begin{attention}[Les trois erreurs qui coûtent cher]
\begin{itemize}
\item \textbf{Introduire un nouvel argument} : la conclusion ne démontre plus ; toute idée nouvelle prouve que le plan était incomplet.
\item \textbf{Répéter mot pour mot l'introduction} : le miroir n'est pas une photocopie ; reformulez.
\item \textbf{Finir par une phrase creuse} : « L'avenir nous le dira », « C'est ainsi que le monde tourne » sont des ouvertures vides, sans lien avec le sujet.
\end{itemize}
\textbf{Sanction chiffrée :} une conclusion absente ou hors-sujet coûte au minimum \textbf{2 points} sur la note globale, au titre de la cohérence générale de la copie. C'est le prix d'une année de travail perdue en trois lignes négligées.
\end{attention}

\courssub{Auto-évaluation et entraînement}

Avant de rendre votre copie, soumettez votre conclusion à la grille suivante :

\begin{center}
\begin{tabularx}{\linewidth}{|l|X|c|}
\hline
\rowcolor{popblueL} \textbf{Critère} & \textbf{Question de contrôle} & \textbf{Oui/Non} \\
\hline
Réponse explicite & Ma conclusion répond-elle clairement à la problématique ? & \\
\hline
Thèse modalisée & Ai-je utilisé un modalisateur et banni le « je pense » ? & \\
\hline
Synthèse fidèle & Mon bilan reprend-il les idées-forces sans argument nouveau ? & \\
\hline
Ouverture pertinente & Mon ouverture a-t-elle un lien logique avec le sujet ? & \\
\hline
Correction de la langue & Orthographe, accords et ponctuation sont-ils relus ? & \\
\hline
\end{tabularx}
\end{center}

\begin{exercice}[Trois conclusions pour un même sujet]
Sujet : « Le progrès technique menace-t-il les traditions camerounaises ? » Rédigez trois conclusions complètes (trois temps chacune) correspondant à trois positions différentes :
\begin{enumerate}
\item une \textbf{thèse radicale} : le progrès détruit les traditions ;
\item une \textbf{thèse nuancée} : le progrès transforme mais peut préserver ;
\item une \textbf{thèse technique} : le progrès est un outil au service de la transmission (numérisation des contes, applications en langues locales), avec une ouverture sur l'éthique professionnelle du technicien.
\end{enumerate}
\end{exercice}

\begin{corrige}[Exercice — éléments de réponse]
\textbf{1. Thèse radicale :} « ... Force est de constater que le progrès technique, en imposant ses rythmes et ses valeurs, balaye les pratiques ancestrales. En somme, il uniformise les modes de vie, dévalorise les savoir-faire locaux et coupe la jeunesse de ses racines. Reste à savoir si une société sans traditions conserve encore une âme. »\\
\textbf{2. Thèse nuancée :} « ... Il apparaît que le progrès transforme les traditions sans nécessairement les anéantir. Au total, s'il menace certains rites, il offre aussi des moyens inédits de les conserver et de les diffuser. Tout dépend de l'usage que la communauté en fait : la tradition ne meurt que lorsqu'on cesse de la transmettre. »\\
\textbf{3. Thèse technique :} « ... En définitive, le progrès technique est un outil neutre que le technicien peut mettre au service de la transmission culturelle : numérisation des contes, applications en langues nationales, archives sonores des griots. C'est pourquoi la formation du technicien camerounais devrait intégrer une dimension éthique : innover, oui, mais durablement et au service de la mémoire collective. »
\end{corrige}

\begin{aretenir}[L'essentiel de la conclusion]
Une conclusion réussie tient en trois temps : \textbf{réponse} explicite et modalisée à la problématique, \textbf{bilan} synthétique des idées-forces (technique du brouillon inversé), \textbf{ouverture} pertinente qui déplace le regard. Jamais de « je pense », jamais d'argument nouveau, jamais de phrase creuse. La conclusion est la dernière image laissée au correcteur : soignez-la comme la porte de sortie d'un beau voyage.
\end{aretenir}

\coursec{Cohésion, Cohérence, Remédiation \& Intégration Finale}

Vous savez désormais construire une introduction solide, rédiger un développement structuré avec citations et transitions, et conclure avec bilan et ouverture. Il reste une étape décisive, celle qui sépare une copie moyenne d'une excellente copie : la \cle{relecture} et la \cle{remédiation}. Au Baccalauréat technique, une dissertation n'est jamais « finie » à la fin de la conclusion : elle est finie quand elle a été \textbf{relue, corrigée et vérifiée}. Cette dernière section vous donne les outils de cette phase finale et vous prépare à l'épreuve complète.

\courssub{Cohérence et cohésion : les deux piliers du texte réussi}

\begin{definition}[Cohérence et cohésion]
La \cle{cohérence} est la qualité d'un texte dont les idées forment un tout logique et compréhensible : c'est le \textbf{sens global} (niveau \emph{macro} : le plan, la progression de l'argumentation, la réponse au sujet). La \cle{cohésion} est l'ensemble des \textbf{liens linguistiques explicites} qui relient les phrases entre elles (niveau \emph{micro} : connecteurs, reprises pronominales, synonymes, temps verbaux).
\end{definition}

L'intuition est simple : la cohérence est le \emph{plan de la maison}, la cohésion est le \emph{ciment entre les briques}. On peut avoir de belles briques (des phrases correctes) sans maison (pas de logique globale) ; on peut avoir un beau plan qui s'écroule faute de ciment (des idées justes mais des phrases juxtaposées sans liens).

\begin{exemplebox}[Deux textes à comparer]
\textbf{Texte A} (cohésion sans cohérence) : « Le vieux nègre reçoit une médaille. Cependant, la pluie tombe sur le village. En effet, les colons sont présents. » Les connecteurs sont là, mais les idées ne se suivent pas : le texte est incohérent.

\textbf{Texte B} (cohérence + cohésion) : « Le vieux nègre reçoit une médaille. \emph{Cette récompense}, censée honorer sa fidélité, révèle en réalité l'hypocrisie coloniale. \emph{En effet}, \emph{il} est décoré non pour sa dignité, mais pour sa soumission. » Mêmes liens, mais chaque phrase prolonge logiquement la précédente.
\end{exemplebox}

\courssub{La progression thématique : le ciment du texte}

\begin{definition}[Progression thématique]
La \cle{progression thématique} est la manière dont l'information circule de phrase en phrase. On distingue :
\begin{itemize}
\item la \cle{progression constante} (ou à thème unique) : chaque phrase reprend le même thème et ajoute une information nouvelle ($T_1 \to T_1 + I_1 \to T_1 + I_2$) ;
\item la \cle{progression en chaîne} (ou linéaire) : l'information nouvelle d'une phrase devient le thème de la suivante ($T_1 \to I_1 \to I_1 + I_2 \to I_2 + I_3$).
\end{itemize}
Les procédés qui assurent cette progression sont : la \textbf{reprise pronominale} (il, celle-ci, ils), la \textbf{reprise nominale} (ce vieillard, cet homme), la \textbf{synonymie} (médaille / récompense / distinction), les \textbf{champs lexicaux} et les \textbf{connecteurs}.
\end{definition}

\begin{propriete}[Choisir sa progression selon la partie]
La \cle{progression en chaîne} assure la \textbf{fluidité narrative} : elle convient à la présentation de l'œuvre dans l'introduction et aux exemples développés. La \cle{progression constante} assure la \textbf{focalisation argumentative} : elle convient au cœur des paragraphes du développement, où l'on examine un même point sous plusieurs angles. Le bon élève \textbf{alterne} les deux selon les parties de sa dissertation.
\end{propriete}

\begin{exemplebox}[Progression constante dans un paragraphe argumentatif]
« \textbf{La médaille} est d'abord un symbole d'aliénation. \textbf{Elle} récompense la soumission et non le mérite. \textbf{Cette distinction} transforme le vieillard en objet de propagande. \textbf{La médaille} devient ainsi l'instrument de sa propre humiliation. » Le thème reste constant ; chaque phrase approfondit l'analyse.
\end{exemplebox}

\begin{methode}[La relecture « à l'envers »]
Pour traquer les fautes d'accord et de cohésion sans être happé par le sens :
\begin{enumerate}
\item Terminez entièrement la rédaction (conclusion comprise).
\item Repartez de la \textbf{dernière phrase} de la copie et remontez vers la première, phrase par phrase.
\item Pour chaque phrase isolée, vérifiez : accord sujet-verbe, accords dans le groupe nominal, temps verbal, ponctuation.
\item Vérifiez ensuite chaque \textbf{charnière} (fin de paragraphe / début du suivant) : y a-t-il un connecteur ou une transition ?
\end{enumerate}
Lire à l'envers casse le flux du sens : votre œil voit enfin les fautes que la lecture normale « corrige » automatiquement.
\end{methode}

\begin{attention}[Piège du Bac technique]
Certains sujets du Baccalauréat technique exigent d'ancrer la réflexion dans la \textbf{vie professionnelle} (ex. : « L'homme et le travail », « Le progrès technique est-il un bien ? »). Ne faites pas de la littérature pure : illustrez avec des réalités de votre filière (normes de sécurité, atelier, chantier, entreprise camerounaise) \emph{en plus} des références littéraires. Une copie qui ignore l'ancrage technique perd des points de pertinence.
\end{attention}

\courssub{Le processus de révision en cinq étapes}

\begin{popfigure}
\begin{center}
\begin{tikzpicture}[node distance=0.55cm,
  etape/.style={rectangle, rounded corners=3pt, draw=popblue, fill=popblueL, text width=4.6cm, align=center, minimum height=0.9cm, font=\small\bfseries},
  detail/.style={rectangle, draw=popmuted, fill=popcream, text width=6.3cm, align=left, font=\small, minimum height=0.9cm},
  fleche/.style={-{Stealth[length=3mm]}, thick, poporange}]
\node[etape] (e1) {1. Macro-structure};
\node[detail, right=0.5cm of e1] (d1) {Sujet respecté ? Problématique résolue ? Plan apparent en 2 ou 3 parties ?};
\node[etape, below=of e1] (e2) {2. Micro-structure};
\node[detail, right=0.5cm of e2] (d2) {Chaque paragraphe : idée + exemple + citation intégrée + analyse ?};
\node[etape, below=of e2] (e3) {3. Cohésion};
\node[detail, right=0.5cm of e3] (d3) {Connecteurs, reprises pronominales, transitions entre les parties ?};
\node[etape, below=of e3] (e4) {4. Langue};
\node[detail, right=0.5cm of e4] (d4) {Orthographe, accords, conjugaison, ponctuation (relecture à l'envers) ?};
\node[etape, below=of e4] (e5) {5. Propreté};
\node[detail, right=0.5cm of e5] (d5) {Ratures propres, paragraphes visibles (alinéas), marge, présentation ?};
\draw[fleche] (e1) -- (e2);
\draw[fleche] (e2) -- (e3);
\draw[fleche] (e3) -- (e4);
\draw[fleche] (e4) -- (e5);
\end{tikzpicture}
\end{center}
\legende{Processus de révision en cinq étapes : on descend du général (le plan) vers le particulier (la présentation). Chaque étape se fait en une passe de relecture distincte.}
\end{popfigure}

\begin{exemplebox}[Chiffre à retenir]
Selon les statistiques d'harmonisation MINESEC, les copies \textbf{relues} (auto-correction effective dans les 30 dernières minutes) gagnent en moyenne \textbf{1,5 à 2 points} sur 20 par rapport aux copies rendues brutes de décoffrage. Trente minutes de relecture valent donc plus que trente minutes de rédaction supplémentaire.
\end{exemplebox}

\courssub{Atelier de remédiation : corriger une copie}

\begin{exoresolu}[Repérer et corriger les défauts d'un extrait de copie]
Énoncé. Voici un extrait de copie d'élève (anonymisée) sur \emph{Le Vieux Nègre et la médaille} de Ferdinand Oyono. Relevez quatre défauts et proposez une correction.

« Meka reçoit une médaille. "Vous êtes un bon nègre" dit le commandant. Les Blancs se moquent de lui. Donc la médaille est une humiliation. Ensuite parlons de la révolte de Meka. »
\tcblower
Solution détaillée.
\begin{enumerate}
\item \textbf{Citation orpheline} : « "Vous êtes un bon nègre" dit le commandant » est jetée sans introduction ni analyse. Correction : « Le commandant résume le mépris colonial en déclarant : \emph{« Vous êtes un bon nègre »} : l'adjectif "bon" réduit Meka à un animal docile. »
\item \textbf{Absence de cohésion} : « Les Blancs se moquent de lui » arrive sans lien. Correction : ajouter une reprise et un connecteur : « \emph{Cette scène} montre que les Blancs, \emph{en réalité}, se moquent de lui. »
\item \textbf{Conclusion hâtive} : « Donc la médaille est une humiliation » affirme sans analyser. Correction : développer le raisonnement avant de conclure le paragraphe.
\item \textbf{Transition manquante} : « Ensuite parlons de la révolte de Meka » est une annonce de plan orale, interdite à l'écrit. Correction : « Si la médaille révèle l'aliénation de Meka, elle provoque aussi, par un retournement décisif, sa prise de conscience. »
\end{enumerate}
\end{exoresolu}

\begin{attention}[Les quatre fautes qui coûtent le plus cher]
\begin{enumerate}
\item La \textbf{citation orpheline} (non introduite, non analysée) ;
\item la \textbf{rupture de cohérence} (un paragraphe qui ne répond plus à la problématique) ;
\item la \textbf{transition annoncée} (« nous allons maintenant voir... », « dans un premier temps je vais parler de... ») ;
\item le \textbf{hors-sujet technique} (oublier d'ancrer les exemples dans la filière quand le sujet l'exige).
\end{enumerate}
\end{attention}

\courssub{La check-list des cinq dernières minutes}

\begin{aretenir}[Check-list de révision finale]
Avant de rendre la copie, cochez mentalement :
\begin{enumerate}
\item \textbf{Sujet respecté ?} (ai-je répondu à la question posée, pas à une autre ?)
\item \textbf{Problématique résolue ?} (ma conclusion répond-elle à la question de l'introduction ?)
\item \textbf{Plan apparent ?} (2 ou 3 grandes parties visibles, alinéas nets ?)
\item \textbf{Citations intégrées ?} (chacune introduite, entre guillemets, analysée ?)
\item \textbf{Transitions fluides ?} (aucune annonce de plan scolaire ?)
\item \textbf{Conclusion + ouverture ?} (bilan + élargissement ?)
\item \textbf{Relecture orthographe / accords ?} (passe « à l'envers » effectuée ?)
\end{enumerate}
\end{aretenir}

\begin{exemplebox}[Gestion du temps au Baccalauréat (4 heures)]
\begin{center}
\begin{tabularx}{\linewidth}{|l|X|}\hline
\rowcolor{popblueL} \textbf{Durée} & \textbf{Étape} \\ \hline
45 min & Analyse du sujet, mobilisation des idées, plan détaillé au brouillon \\ \hline
30 min & Rédaction de l'introduction (au brouillon puis au propre) \\ \hline
1 h 45 & Rédaction du développement (environ 35 min par partie) \\ \hline
30 min & Rédaction de la conclusion (bilan + ouverture) \\ \hline
30 min & Relecture et correction (processus en 5 étapes) \\ \hline
\end{tabularx}
\end{center}
Règle d'or : \textbf{jamais} de développement au brouillon intégral ; seuls le plan et l'introduction y passent.
\end{exemplebox}

\begin{savaistu}[Le capital sympathie du correcteur]
Un correcteur du Baccalauréat lit en moyenne \textbf{80 copies en 4 heures}, soit environ 3 minutes par copie en première approche. Une copie \textbf{aérée} (alinéas visibles), \textbf{structurée} (parties nettes), \textbf{sans ratures}, avec une introduction et une conclusion soignées, part avec un énorme capital sympathie : le correcteur la lit avec bienveillance. À contenu égal, la présentation fait la différence.
\end{savaistu}

\courssub{Grille d'auto-évaluation et de co-évaluation}

\begin{popfigure}
\begin{center}
\begin{tikzpicture}
\draw[fill=popblueL, draw=popblue] (0,0) rectangle (15,0.8);
\node[font=\small\bfseries, anchor=west] at (0.15,0.4) {Critère};
\node[font=\small\bfseries, anchor=west] at (4.2,0.4) {Indicateur observable};
\node[font=\small\bfseries, anchor=west] at (10.2,0.4) {Oui / Non};
\node[font=\small\bfseries, anchor=west] at (12.4,0.4) {Commentaire};
\foreach \y/\crit/\ind in {
  -1.0/Introduction/{Accroche, sujet, problématique, plan annoncés},
  -2.0/Développement/{2 ou 3 parties ; idée + exemple + citation + analyse},
  -3.0/Cohésion/{Connecteurs, reprises, transitions rédigées},
  -4.0/Citations/{Introduites, exactes, analysées},
  -5.0/Conclusion/{Bilan, réponse à la problématique, ouverture},
  -6.0/Langue et présentation/{Orthographe, accords, alinéas, propreté}} {
  \draw[draw=popmuted] (0,\y) rectangle (15,\y+0.8);
  \node[font=\small, anchor=west] at (0.15,\y+0.4) {\crit};
  \node[font=\small, anchor=west] at (4.2,\y+0.4) {\ind};
  \node[font=\small, anchor=west] at (10.2,\y+0.4) {$\square$ \quad $\square$};
  \draw[draw=popmuted] (4.0,\y) -- (4.0,\y+0.8);
  \draw[draw=popmuted] (10.0,\y) -- (10.0,\y+0.8);
  \draw[draw=popmuted] (12.2,\y) -- (12.2,\y+0.8);
}
\draw[draw=popmuted] (4.0,0) -- (4.0,0.8);
\draw[draw=popmuted] (10.0,0) -- (10.0,0.8);
\draw[draw=popmuted] (12.2,0) -- (12.2,0.8);
\end{tikzpicture}
\end{center}
\legende{Grille d'auto-évaluation et de co-évaluation : à utiliser en binôme lors des ateliers de remédiation, puis seul avant de rendre toute dissertation.}
\end{popfigure}

\courssub{Préparation à l'oral et compte rendu}

L'évaluation ne s'arrête pas à la copie : l'élève doit savoir \textbf{justifier sa démarche} à l'oral, comme devant un jury.

\begin{methode}[Défendre sa dissertation à l'oral (simulation jury)]
\begin{enumerate}
\item \textbf{Justifier son plan} : « J'ai choisi un plan dialectique parce que le sujet opposait deux thèses... » (une phrase par partie, en expliquant la \emph{logique} d'enchaînement).
\item \textbf{Expliquer le choix d'une citation} : dire \emph{où} elle se situe dans l'œuvre, \emph{pourquoi} elle illustre l'argument, et ce qu'elle apporte de plus qu'une paraphrase.
\item \textbf{Répondre à un contre-argument} : le jury joue l'avocat du diable (« La médaille n'est-elle pas malgré tout un honneur ? »). Répondre en trois temps : reconnaître la pertinence (« Il est vrai que... »), nuancer (« cependant... »), réaffirmer sa thèse avec une preuve.
\end{enumerate}
\end{methode}

\courssub{Le dossier personnel : la « boîte à outils » de l'élève}

Chaque élève constitue au fil de l'année un \cle{dossier personnel} de dissertation, réutilisable le jour de l'examen (dans la mémoire, non sur la table !) :

\begin{itemize}
\item \textbf{Fiches connecteurs} : addition (de plus, en outre), opposition (cependant, néanmoins), cause (car, en effet, puisque), conséquence (ainsi, par conséquent), concession (certes... mais), illustration (ainsi, c'est le cas de).
\item \textbf{Plans types} : thématique, dialectique (thèse / antithèse / synthèse), analytique (constat / causes / conséquences), avec pour chacun un exemple de sujet traité.
\item \textbf{Citations clés} : un répertoire de 10 à 15 citations apprises d'Oyono (\emph{Le Vieux Nègre et la médaille}, \emph{Une vie de boy}) et de Mongo Beti, classées par thème (colonialisme, aliénation, résistance, dignité).
\item \textbf{Modèles d'introduction et de conclusion} : deux ou trois gabarits réussis et corrigés en classe, à adapter.
\item \textbf{Grille d'auto-évaluation} (voir figure ci-dessus) et check-list finale.
\end{itemize}

\courssub{Évaluation formative : la dissertation complète}

\begin{exercice}[Dissertation d'intégration (notée sur 20)]
Sujet : « La médaille, dans le roman d'Oyono, est-elle une récompense ou un instrument de domination ? » Vous rédigerez une dissertation complète (introduction, développement en deux ou trois parties, conclusion avec ouverture) en vous appuyant sur votre boîte à outils. Barème indicatif inspiré de la grille officielle MINESEC : respect du sujet et problématique (4 pts) ; cohérence du plan et de l'argumentation (6 pts) ; cohésion, citations et transitions (4 pts) ; langue et présentation (6 pts).
\end{exercice}

\begin{corrige}[Exercice — modalités du retour individualisé]
La correction se fait en deux temps. \textbf{1. Correction collective} : lecture d'une copie-type réussie, repérage des attendus du barème, comparaison avec la grille d'auto-évaluation. \textbf{2. Retour individualisé} : chaque copie reçoit une fiche en deux colonnes — \emph{Points forts} (au moins deux, nommés précisément : « problématique claire », « citation bien analysée p. 2 ») et \emph{Axes de progression} (au plus deux, formulés comme des consignes actionnables : « introduire chaque citation par un verbe de parole », « rédiger la transition I/II »). L'élève réécrit ensuite \textbf{un seul paragraphe} en appliquant ses axes de progression : c'est cette réécriture ciblée, et non la note, qui produit le progrès.
\end{corrige}

\begin{aretenir}[L'essentiel de la section]
\begin{itemize}
\item \textbf{Cohérence} = logique globale des idées ; \textbf{cohésion} = liens linguistiques entre les phrases. Les deux sont indispensables.
\item La \textbf{progression thématique} (constante ou en chaîne) est le ciment du texte ; on l'alterne selon les parties.
\item La révision suit \textbf{cinq étapes} : macro-structure, micro-structure, cohésion, langue, propreté — avec la technique de la \textbf{relecture à l'envers}.
\item Au Bac : \textbf{45 / 30 / 1 h 45 / 30 / 30} ; les 30 dernières minutes de relecture valent 1,5 à 2 points.
\item La \textbf{boîte à outils} personnelle (connecteurs, plans types, citations, modèles) est le trésor de guerre du candidat.
\end{itemize}
\end{aretenir}

\newpage

\coursec{Activité d'intégration}

\coursec{II : Rédiger une dissertation}

\courssub{Objectifs de l'activité}
\noindent Cette activité d'intégration, placée en fin d'unité, vise à \textbf{mobiliser l'ensemble des savoirs et savoir-faire} travaillés lors des séquences précédentes (structure de la dissertation, argumentation, insertion de citations, progression thématique, rédaction de l'introduction et de la conclusion, cohésion/cohérence, expression orale) dans une \textbf{situation authentique d'examen} (Probatoire / Baccalauréat technique). Concrètement, l'élève doit être capable de :
\begin{itemize}
    \item \textbf{Analyser} un sujet de dissertation à double entrée (héritage / modernité) pour en dégager la problématique et construire un plan dialectique pertinent.
    \item \textbf{Rédiger} une introduction complète (amorce, sujet posé, sujet divisé, problématique, annonce du plan) et une conclusion (bilan synthétique, réponse à la problématique, ouverture).
    \item \textbf{Construire} un développement structuré en trois parties, composé de paragraphes types (argument, explication, illustration par citation/document, transition).
    \item \textbf{Exploiter} le corpus fourni (extraits de \emph{Le Vieux Nègre et la Médaille}, données statistiques MINJEC, article de presse) pour étayer l'argumentation.
    \item \textbf{Maîtriser} les outils de la cohérence textuelle (progression thématique : reprise, chaînage, progression ; connecteurs logiques ; champs lexicaux).
    \item \textbf{Défendre} sa production à l'oral : justifier ses choix de plan, sa thèse, son ouverture, face à un jury simulé.
\end{itemize}

\courssub{Situation}
\noindent \textbf{Contexte :} Le Ministère de la Jeunesse et de l'Éducation Civique (MINJEC) lance, dans le cadre de la \emph{Semaine Nationale de la Jeunesse 2025}, un concours national d'essais destiné aux élèves de Terminale des lycées techniques du Cameroun. Le thème retenu est :
\begin{center}
    \textbf{« Héritage culturel et modernité technique : le jeune camerounais face au défi de l'identité »}
\end{center}

\noindent \textbf{Corpus documentaire fourni aux candidats :}
\begin{enumerate}
    \item \textbf{Sujet de dissertation} (cf. thème ci-dessus).
    \item \textbf{Extraits de \emph{Le Vieux Nègre et la Médaille} (Ferdinand Oyono) :}
        \begin{itemize}
            \item \textit{Extrait A (Discours du Commandant) :} « …Cette médaille, mes amis, c'est la preuve que la France n'oublie pas ses enfants… Elle marque l'entrée de Meka dans la grande famille de la civilisation… » (Symbolique de la reconnaissance coloniale / assimilation).
            \item \textit{Extrait B (Monologue intérieur de Meka) :} « …Il la sentait lourde sur sa poitrine, cette médaille. Lourde comme le fardeau des ans… Et si c'était une chaîne dorée ? … Moi, Meka, le premier médaillé de la région… » (Ambivalence : fierté identitaire vs aliénation symbolique).
            \item \textit{Extrait C (Réaction de Kelara, la femme de Meka) :} « …Elle regarda la médaille, puis son mari… "C'est beau", dit-elle simplement. Mais dans ses yeux, il y avait cette lueur ironique qui disait : "Tu as vendu ton âme pour un bout de fer-blanc"… » (Regard critique, lucidité féminine, résistance silencieuse).
        \end{itemize}
    \item \textbf{Infographie MINJEC 2023 — « Pratiques langagières des 15-25 ans au Cameroun » :}
        \begin{center}
            \begin{tabularx}{\linewidth}{|l|X|X|}\hline
            \rowcolor{popblueL} \textbf{Indicateur} & \textbf{Zone Urbaine (Yaoundé/Douala)} & \textbf{Zone Rurale} \\ \hline
            Usage principal \textbf{Français/Anglais} (maison/école) & 68\% & 42\% \\ \hline
            Usage principal \textbf{Langues nationales} (maison) & 22\% & 55\% \\ \hline
            \textbf{Code-switching} (mélange Fr/An + Langue nationale) & 85\% & 60\% \\ \hline
            Sentiment d'appartenance \textbf{"Camerounais d'abord"} & 72\% & 88\% \\ \hline
            \end{tabularx}
        \end{center}
        \legende{Source : Enquête MINJEC 2023, échantillon 5\,000 jeunes.}
    \item \textbf{Article de presse : « L'IA générative : menace pour les langues africaines ? » (\emph{Le Jour}, mars 2024).}
        \textit{Extrait clé :} « …Les modèles de langage (LLM) sont entraînés à 95\% sur des corpus anglais/français. Les langues à faible ressource (bassa, bamileke, fulfulde, etc.) risquent l'extinction numérique. Pourtant, des start-ups camerounaises (ex: \emph{KamerNLP}) développent des outils de traduction automatique pour "nourrir" l'IA en langues locales. La technique peut devenir le gardien de l'héritage, à condition d'une volonté politique et citoyenne… »
\end{enumerate}

\noindent \textbf{Déroulement de l'atelier (3 séances / 5h au total) :}
\begin{itemize}
    \item \textbf{Séance 1 (2h) :} Analyse du sujet $\rightarrow$ Plan détaillé $\rightarrow$ Rédaction de l'Introduction $\rightarrow$ Rédaction du 1er paragraphe de développement (Thèse : L'héritage comme socle) avec insertion d'une citation d'Oyono + Transition vers l'antithèse.
    \item \textbf{Séance 2 (2h) :} Rédaction de la 2ème partie (Antithèse : Modernité technique, risque d'effacement) et 3ème partie (Synthèse : Réappropriation créative) $\rightarrow$ Rédaction de la Conclusion $\rightarrow$ Relecture croisée en binôme avec \textbf{Grille Cohésion/ Cohérence} (voir Ressources).
    \item \textbf{Séance 3 (1h) :} \textbf{Simulation orale} (5 min par élève) : « Défendez votre thèse face au jury : Pourquoi votre plan est-il le plus pertinent ? Justifiez votre ouverture. » $\rightarrow$ Évaluation par les pairs (grille simplifiée) + Feedback professeur.
\end{itemize}

\courssub{Consignes}
\noindent Les consignes ci-dessous guident le travail de l'élève tout au long des trois séances. Elles sont conçues pour activer progressivement les compétences ciblées.

\textbf{Séance 1 — Analyse, Planification, Introduction, Amorce du développement}
\begin{enumerate}
    \item \textbf{Analyse du sujet :} Identifiez les \cle{mots-clés} (Héritage culturel, Modernité technique, Jeune camerounais, Défi, Identité). Dégagez les \cle{termes du débat} : opposition / complémentarité ? Définissez le \cle{champ notionnel} (culture, technique, identité, aliénation, réappropriation).
    \item \textbf{Problématisation :} Formulez \textbf{une seule problématique} ouverte, issue de la tension entre les termes. \textit{Exemple attendu :} « Dans quelle mesure la modernité technique, loin de dissoudre l'héritage culturel, peut-elle devenir le levier d'une réinvention identitaire pour la jeunesse camerounaise ? »
    \item \textbf{Construction du plan détaillé :} Adoptez une structure \cle{dialectique} (Thèse / Antithèse / Synthèse) ou \cle{analytique} (Causes / Conséquences / Solutions) justifiée. Rédigez les \cle{titres des trois parties} et les \cle{sous-arguments} (idées directrices) pour chaque partie. Précisez \textbf{quel document du corpus} illustre chaque sous-argument.
    \item \textbf{Rédaction de l'Introduction :} Appliquez la méthode en \textbf{5 étapes} (Amorce contextuelle $\rightarrow$ Sujet posé $\rightarrow$ Sujet divisé / Débat $\rightarrow$ Problématique $\rightarrow$ Annonce du plan). Soignez l'accroche (citation, chiffre clé de l'infographie, actualité).
    \item \textbf{Premier paragraphe de développement (Thèse) :} Rédigez le paragraphe type : \textbf{Idée directrice} (L'héritage culturel est le socle identitaire indispensable) $\rightarrow$ \textbf{Explication} $\rightarrow$ \textbf{Illustration} par \textbf{citation intégrée} (Extrait B : Meka « lourd comme le fardeau des ans ») $\rightarrow$ \textbf{Transition} vers l'antithèse (Mais cet héritage est mis à l'épreuve…).
\end{enumerate}

\textbf{Séance 2 — Achèvement du développement, Conclusion, Révision}
\begin{enumerate}
    \setcounter{enumi}{5}
    \item \textbf{Deuxième partie (Antithèse) :} « La modernité technique : facteur de rupture et de risque d'effacement identitaire ». Deux paragraphes minimum. Utilisez l'\textbf{Infographie MINJEC} (chiffres sur perte langues nationales) et l'\textbf{Article \emph{Le Jour}} (IA, extinction numérique). Insérez au moins une citation courte de l'article. Assurez la \cle{transition} vers la synthèse.
    \item \textbf{Troisième partie (Synthèse) :} « Vers une hybridation féconde : le jeune camerounais, acteur de la réappropriation ». Deux paragraphes. Mobilisez l'\textbf{Extrait C (Kelara)} (lucidité/résistance) et l'exemple des \textbf{start-ups camerounaises} (article \emph{Le Jour}). Montrez comment la technique \cle{sert} l'héritage.
    \item \textbf{Rédaction de la Conclusion :} Appliquez la méthode en \textbf{3 temps} : \textbf{Bilan synthétique} (reprise des 3 parties sans copier-coller) $\rightarrow$ \textbf{Réponse directe à la problématique} (thèse finale) $\rightarrow$ \textbf{Ouverture} (perspective concrète : ex. réforme curriculaire, rôle de l'école technique, engagement citoyen).
    \item \textbf{Relecture croisée (Binôme) :} Échangez les copies. Utilisez la \textbf{Grille Cohésion/Cohérence} (voir Ressources) pour annoter : progression thématique (reprise/chaînage), connecteurs, ponctuation, accord participes, pertinence citations. Rendez la copie annotée à l'auteur pour correction finale.
\end{enumerate}

\textbf{Séance 3 — Défense orale (Simulation Jury)}
\begin{enumerate}
    \setcounter{enumi}{9}
    \item \textbf{Préparation (10 min) :} Relisez votre copie. Préparez 3 arguments-clés pour défendre : 1) La pertinence de votre plan (pourquoi dialectique ?). 2) La solidité de votre thèse (pourquoi cette synthèse ?). 3) La pertinence de votre ouverture (pourquoi celle-ci ?).
    \item \textbf{Passage oral (5 min) :} Face au jury (2 élèves + 1 prof), présentez votre défense sans lire votre dissertation. Répondez aux questions du jury (ex: « Pourquoi avoir mis Kelara en synthèse et pas en antithèse ? », « Votre ouverture n'est-elle pas trop vague ? »).
    \item \textbf{Évaluation :} Remplissage de la \textbf{Grille Orale} (Contenu / Argumentation / Expression orale / Attitude) par les pairs et le prof. Débriefing collectif (10 min).
\end{enumerate}

\courssub{Ressources à mobiliser}
\noindent \textbf{1. Structure type de la dissertation (Rappel mémento) :}
\begin{itemize}
    \item \textbf{Introduction :} \cle{Amorce} (Général $\rightarrow$ Spécifique) $\rightarrow$ \cle{Sujet posé} (Reformulation) $\rightarrow$ \cle{Sujet divisé} (Débat / Tension) $\rightarrow$ \cle{Problématique} (Question unique, ouverte) $\rightarrow$ \cle{Annonce du plan} (Mots-clés des parties).
    \item \textbf{Développement :} 3 Parties (I, II, III). Chaque partie = 2 à 3 paragraphes. \textbf{Paragraphe type :} \cle{Idée directrice} (en tête) $\rightarrow$ \cle{Explication/Analyse} $\rightarrow$ \cle{Preuve/Illustration} (Citation intégrée + Référence précise + Commentaire) $\rightarrow$ \cle{Transition} (Lien logique vers idée suivante).
    \item \textbf{Conclusion :} \cle{Bilan} (Synthèse I + II + III) $\rightarrow$ \cle{Réponse/Thèse} (Affirmation tranchée) $\rightarrow$ \cle{Ouverture} (Nouvel horizon, question liée, perspective concrète).
\end{itemize}

\noindent \textbf{2. Progression thématique (Cohérence) — Outils linguistiques :}
\begin{center}
    \begin{tabularx}{\linewidth}{|l|X|X|}\hline
    \rowcolor{popblueL} \textbf{Type} & \textbf{Mécanisme} & \textbf{Exemple / Marqueurs} \\ \hline
    \textbf{Reprise (Anaphore)} & Reprise du \textbf{même thème} (sujet) en début de phrase/paragraphe. & \textit{« L'héritage culturel… Il constitue… Cette mémoire… Elle permet… »} \\ \hline
    \textbf{Chaînage} & Thème de la phrase N = Rhème de la phrase N-1. & \textit{« La technique menace les langues. \textbf{Cette menace} pousse à l'innovation. \textbf{L'innovation} nécessite des ressources… »} \\ \hline
    \textbf{Progression} & Thème constant, Rhèmes qui s'enchaînent logiquement (cause $\rightarrow$ conséquence $\rightarrow$ solution). & \textit{« Le jeune camerounais \textbf{hérite}. Il \textbf{reçoit}. Il \textbf{transforme}. Il \textbf{transmet}. »} \\ \hline
    \textbf{Connecteurs logiques} & Articulent la relation inter-phrastique / inter-paragraphique. & \textit{Addition : de plus, en outre. Opposition : cependant, pourtant. Conséquence : par conséquent, dès lors. Explication : en effet, c'est que…} \\ \hline
    \end{tabularx}
\end{center}

\noindent \textbf{3. Insertion de citation (Norme Bac) :}
\begin{itemize}
    \item \textbf{Courte ($<$ 3 lignes) :} Intégrée à la phrase, entre guillemets français « … », commentée immédiatement. \textit{Ex: Meka ressent ce poids comme une « chaîne dorée » (Extrait B), symbole de l'aliénation consentie.}
    \item \textbf{Longue ($>$ 3 lignes) :} Alinéa, retrait, corps plus petit, sans guillemets, introduite par deux-points. Commentaire obligatoire après.
    \item \textbf{Référence :} (Auteur, Œuvre, Extrait/Acte/Scène) ou (Source, Date) pour documents non littéraires.
\end{itemize}

\noindent \textbf{4. Grille de relecture croisée — Cohésion / Cohérence (Séance 2) :}
\begin{center}
    \begin{tabularx}{\linewidth}{|l|X|c|}\hline
    \rowcolor{popblueL} \textbf{Critère} & \textbf{Indicateurs observables} & \textbf{Validé ?} \\ \hline
    \textbf{Progression thématique} & Identification du thème de départ + types de progression (Reprise / Chaînage / Progression) repérés dans le développement. & $\square$ \\ \hline
    \textbf{Connecteurs} & Présence, variété, justesse (pas de « et » systématique). Liens logiques explicites entre parties (Transitions). & $\square$ \\ \hline
    \textbf{Citations} & Intégration syntaxique correcte / Commentaire explicite / Référence précise (Extrait A, B, C, Infographie, Article). & $\square$ \\ \hline
    \textbf{Structure paragraphaire} & 1 Idée directrice claire par paragraphe / Alinéa visible / Ponctuation forte. & $\square$ \\ \hline
    \textbf{Langue (Code)} & Accords (sujet/verbe, participe passé) / Temps (présent de vérité, passé narration) / Vocabulaire précis (pas de familier). & $\square$ \\ \hline
    \end{tabularx}
\end{center}

\noindent \textbf{5. Grille d'évaluation orale (Séance 3) — Simplifiée (Sur 10 pts) :}
\begin{itemize}
    \item \textbf{Contenu / Argumentation (4 pts) :} Pertinence de la défense du plan (2 pts) + Justification de la thèse/ouverture (2 pts).
    \item \textbf{Expression orale (3 pts) :} Clarté, débit, vocabulaire précis, syntaxe orale correcte, regard/contact jury.
    \item \textbf{Attitude / Interaction (3 pts) :} Écoute questions, reformulation, aisance, respect temps (5 min).
\end{itemize}

\courssub{Démarche et corrigé commenté}
\begin{exoresolu}{Corrigé type guidé — « La Médaille de la Modernité »}
\textbf{ÉTAPE 1 — ANALYSE DU SUJET ET PROBLÉMATIQUE}
\begin{itemize}
    \item \textbf{Mots-clés :} \cle{Héritage culturel} (traditions, langues, mémoire, valeurs, médaille Meka) ; \cle{Modernité technique} (école, langues officielles, IA, numérique, start-ups) ; \cle{Jeune camerounais} (sujet acteur, 15-25 ans, cible stats MINJEC) ; \cle{Défi de l'identité} (tension, construction, risque perte, opportunité).
    \item \textbf{Débat (Sujet divisé) :} L'héritage est-il un frein ou un tremplin face à la modernité ? La technique détruit-elle l'identité ou offre-t-elle de nouveaux outils pour la dire ?
    \item \textbf{Problématique retenue :} \textbf{« Comment le jeune camerounais peut-il transformer la tension entre l'héritage culturel et la modernité technique en une dynamique de réinvention identitaire, sans subir ni l'aliénation (médaille dorée), ni l'effacement (extinction numérique) ? »}
    \item \textit{Justification :} Cette problématique intègre les 4 documents (Médaille = aliénation ; Stats = effacement ; IA = menace/outil ; Kelara/Start-ups = réappropriation) et appelle une structure dialectique.
\end{itemize}

\textbf{ÉTAPE 2 — PLAN DÉTAILLÉ (Structure Dialectique)}

\textbf{INTRODUCTION} (À rédiger in extenso — voir ci-dessous)

\textbf{DÉVELOPPEMENT}
\begin{itemize}
    \item \textbf{I. L'HÉRITAGE CULTUREL : SOCLE INDISPENSABLE ET ENJEU DE RECONNAISSANCE (Thèse)}
    \begin{itemize}
        \item A. L'héritage comme mémoire vivante et ancrage identitaire (Extrait B : Meka « fardeau des ans », fierté d'être « premier médaillé »).
        \item B. La quête de reconnaissance légitime par l'Autre (Extrait A : Discours Commandant « entrée dans la civilisation » — analyse critique : reconnaissance asymétrique).
        \item C. La résistance silencieuse et lucide face à l'assimilation (Extrait C : Kelara, « lueur ironique », « vendu ton âme » — transition : mais cet héritage est aujourd'hui fragilisé par une modernité technique invasive).
    \end{itemize}
    \item \textbf{II. LA MODERNITÉ TECHNIQUE : FACTEUR DE RUPTURE ET RISQUE D'EFFACEMENT IDENTITAIRE (Antithèse)}
    \begin{itemize}
        \item A. L'école et les langues officielles comme vecteurs de déracinement progressif (Infographie MINJEC : 68\% usage Fr/An en zone urbaine vs 22\% langues nationales ; corrélation avec sentiment d'appartenance).
        \item B. La révolution numérique et l'IA : l'extinction numérique des langues à « faible ressource » (Article \emph{Le Jour} : 95\% corpus anglais/français, « menace pour les langues africaines »).
        \item C. La médaille moderne : le statut social conditionné par la maîtrise technique exclusive (Transition : pourtant, la technique n'est pas fatale ; elle peut être détournée, réappropriée).
    \end{itemize}
    \item \textbf{III. VERS UNE HYBRIDATION FÉCONDE : LE JEUNE CAMEROUNAIS, ACTEUR DE LA RÉAPPROPRIATION (Synthèse)}
    \begin{itemize}
        \item A. La lucidité critique comme point de départ (Extrait C : Kelara, modèle de conscience claire ; le jeune ne doit pas être dupe de la « médaille »).
        \item B. L'innovation endogène : la technique au service de l'héritage (Article \emph{Le Jour} : \emph{KamerNLP}, traduction automatique, « nourrir l'IA en langues locales »).
        \item C. L'école technique comme lieu privilégié de cette synthèse : former des « ingénieurs de culture » (Proposition concrète : modules langues nationales en codage, projets fin d'études sur patrimoine numérisé).
    \end{itemize}
\end{itemize}

\textbf{CONCLUSION} (À rédiger in extenso — voir ci-dessous)

\tcblower

\textbf{ÉTAPE 3 — RÉDACTION DE L'INTRODUCTION (Modèle commenté)}

\begin{quote}
    \textit{[Amorce contextuelle — Accroche par le corpus]} « Tu as vendu ton âme pour un bout de fer-blanc », lance Kelara à son mari Meka, fraîchement décoré de la médaille coloniale dans \emph{Le Vieux Nègre et la Médaille} de Ferdinand Oyono. Cette réplique cinglante résume le drame séculaire de l'Africain face à la modernité importée : l'échange inégal d'une identité profonde contre une reconnaissance symbolique.
    
    \textit{[Sujet posé — Reformulation]} Aujourd'hui, au Cameroun, ce « bout de fer-blanc » a changé de nature : ce n'est plus seulement la médaille coloniale, mais le smartphone, le code informatique, l'intelligence artificielle, le français ou l'anglais imposés comme seules langues de la réussite. Le jeune camerounais, qui représente plus de 60\% de la population, se trouve au cœur de ce face-à-face entre un \textbf{héritage culturel} multiséculaire (langues, traditions, valeurs communautaires) et une \textbf{modernité technique} planétaire, vectrice de progrès mais aussi d'uniformisation.
    
    \textit{[Sujet divisé — Débat]} Si l'héritage apparaît comme le socle sans lequel l'individu se perd (Meka sans ses ancêtres n'est rien), la modernité technique s'impose comme un horizon incontournable pour exister dans le monde contemporain (statistiques MINJEC : 68\% des urbains parlent français/anglais à la maison). Faut-il dès lors choisir : tradition ou modernité ? Identité ou compétence technique ?
    
    \textit{[Problématique]} \textbf{Comment le jeune camerounais peut-il transformer la tension entre l'héritage culturel et la modernité technique en une dynamique de réinvention identitaire, sans subir ni l'aliénation, ni l'effacement ?}
    
    \textit{[Annonce du plan]} Pour répondre, nous montrerons d'abord que l'héritage culturel constitue le socle identitaire indispensable, mais qu'il est menacé par une modernité technique porteuse de rupture (I). Nous analyserons ensuite comment cette modernité, par l'école et le numérique, agit comme un facteur d'effacement linguistique et mémoriel (II). Enfin, nous démontrerons que la réappropriation critique des outils techniques par la jeunesse elle-même ouvre la voie à une hybridation féconde, seule garante d'une identité vivante (III).
\end{quote}

\noindent \textbf{Commentaire méta :} L'amorce ancre directement dans l'œuvre au programme (Oyono). Le sujet posé définit les termes techniques. Le sujet divisé oppose les deux pôles avec des preuves chiffrées. La problématique est unique, ouverte, issue de la tension. L'annonce du plan utilise les mots-clés des parties (Socle / Menace / Réappropriation).

\textbf{ÉTAPE 4 — PREMIER PARAGRAPHE DÉVELOPPEMENT (I-A) AVEC CITATION ET TRANSITION}

\begin{quote}
    \textbf{[Idée directrice]} L'héritage culturel n'est pas un folklore figé ; il est la \cle{mémoire vivante} qui ancre le sujet dans une temporalité longue et lui donne une épaisseur existentielle.
    
    \textbf{[Explication]} Pour Meka, le vieil homme d'Oyono, cet héritage se ressent physiquement : le poids des ans, le poids des ancêtres. La médaille qu'il reçoit n'est pas un objet neutre ; elle vient se greffer sur cette histoire charnelle.
    
    \textbf{[Illustration — Citation intégrée]} C'est ce que révèle son monologue intérieur lorsqu'il confesse : \og Il la sentait lourde sur sa poitrine, cette médaille. Lourde comme le fardeau des ans \fg{} (Extrait B). La comparaison \og comme le fardeau des ans \fg{} établit une \cle{progression thématique par chaînage} : le poids de la médaille (modernité/colonisateur) se confond avec le poids de l'héritage (ancêtres/tradition). Meka ne peut porter la première sans assumer le second.
    
    \textbf{[Commentaire de la citation]} Cette « lourdeur » est ambivalente : elle signifie à la fois la charge de la mémoire (fierté d'être « le premier médaillé ») et la contrainte de l'histoire. L'héritage est donc le \textbf{prérequis} à toute réception de la modernité.
    
    \textbf{[Transition vers I-B]} Toutefois, cette reconnaissance de l'héritage par la modernité (la médaille remise par le Commandant) repose sur un malentendu fondamental : l'assimilation.
\end{quote}

\noindent \textbf{Commentaire méta :} Respect du schéma paragraphe type. Citation courte intégrée à la syntaxe (guillemets français). Commentaire linguistique (comparaison, chaînage). Transition explicite vers le sous-argument suivant.

\textbf{ÉTAPE 5 — DEUXIÈME PARTIE (II) — EXEMPLE DE PARAGRAPHE (II-A) AVEC DONNÉES STATISTIQUES}

\begin{quote}
    \textbf{[Idée directrice]} L'école technique et l'administration, par l'imposition des langues officielles, opèrent un \cle{déracinement linguistique} progressif, premier pas vers l'effacement identitaire.
    
    \textbf{[Explication + Données]} L'infographie MINJEC 2023 est sans appel : en zone urbaine, \textbf{68\%} des 15-25 ans utilisent principalement le français ou l'anglais à la maison et à l'école, contre seulement \textbf{22\%} pour les langues nationales. Ce basculement domestique signe la relégation de l'héritage linguistique au second plan.
    
    \textbf{[Analyse]} Ce n'est pas un hasard si le \textbf{code-switching} (mélange) atteint 85\% en ville : il marque une transition forcée, non un bilinguisme équilibré. La maîtrise technique (codage, bureautique, communication professionnelle) s'acquiert quasi exclusivement en langue officielle, créant une \cle{corrélation forte} entre « réussite technique » et « abandon de la langue maternelle ».
    
    \textbf{[Transition vers II-B]} Ce phénomène s'accélère dramatiquement avec l'avènement de l'intelligence artificielle générative.
\end{quote}

\textbf{ÉTAPE 6 — TROISIÈME PARTIE (III) — SYNTHÈSE ET OUVERTURE CONCRÈTE (III-B)}

\begin{quote}
    \textbf{[Idée directrice]} La réponse au défi ne vient pas d'un retour au passé, mais de la \cle{réappropriation technique} : faire de l'outil numérique le gardien de l'héritage.
    
    \textbf{[Illustration — Article Le Jour]} L'article de \emph{Le Jour} (2024) cite l'exemple de \emph{KamerNLP}, start-up camerounaise qui développe des modèles de traduction pour « nourrir l'IA en langues locales ». C'est l'inversion totale du rapport de force : la technique (IA), menace identifiée (extinction numérique), devient \textbf{condition de survie} des langues (bassa, bamileke, fulfulde).
    
    \textbf{[Lien avec Oyono]} Là où Meka subissait la médaille (Extrait A) et où Kelara la déconstruisait par l'ironie (Extrait C), le jeune ingénieur de \emph{KamerNLP} \textbf{agit}. Il ne demande pas la reconnaissance de l'Autre (Commandant), il construit l'outil pour la sienne (Communauté).
    
    \textbf{[Transition vers Conclusion]} Cette dynamique prouve que l'identité camerounaise de demain ne se jouera pas \og contre \fg{} la technique, mais \og par \fg{} elle.
\end{quote}

\textbf{ÉTAPE 7 — RÉDACTION DE LA CONCLUSION (Modèle commenté)}

\begin{quote}
    \textit{[Bilan synthétique]} Au terme de cette analyse, il apparaît que la médaille de Meka — symbole de l'héritage reconnu par l'Autre — a muté. Elle n'est plus un morceau de métal offert par le Commandant, mais l'ensemble des compétences techniques (numérique, IA, langues officielles) que la modernité propose au jeune camerounais. Nous avons vu que cet héritage est un \textbf{socle vital} (I), que la modernité technique porte en elle un \textbf{risque majeur d'effacement} linguistique et mémoriel via l'école et le numérique (II), mais qu'une \textbf{réappropriation citoyenne et innovante} (start-ups, code-switching conscient, école technique rénovée) permet d'inverser la vapeur (III).
    
    \textit{[Réponse à la problématique — Thèse finale]} \textbf{Le défi de l'identité ne se résout pas par le refus de la modernité, ni par l'oubli de l'héritage, mais par leur hybridation créative : le jeune camerounais doit devenir le « forgeron » de sa propre médaille, utilisant le feu de la technique pour tremper l'acier de sa culture.}
    
    \textit{[Ouverture — Perspective concrète]} Cette forge passe nécessairement par l'\textbf{École Technique}. Intégrer l'apprentissage des langues nationales et du patrimoine culturel \emph{dans} les modules de programmation, de CAO/DAO, de maintenance industrielle — et non plus en marge — serait le gage d'une modernité \og camerounaisée \fg{}. Comme le suggère l'ironie lucide de Kelara, il ne s'agit pas de refuser la médaille, mais de ne plus en faire le prix de son âme.
\end{quote}

\noindent \textbf{Commentaire méta :} Le bilan reprend les 3 mots-clés des parties (Socle / Risque / Réappropriation) sans répétition littérale. La thèse finale est une phrase forte, imagée (« forgeron », « tremper l'acier »), répondant directement à la problématique. L'ouverture est ancrée dans le réel de la série (École Technique, modules pro) et fait écho à l'œuvre étudiée (Kelara).
\end{exoresolu}

\courssub{Bilan de l'activité}
\noindent Cette activité d'intégration \textbf{« La Médaille de la Modernité »} a permis de \textbf{valider l'atteinte des objectifs terminaux} du programme de Français en Terminale Technique :

\begin{itemize}
    \item \textbf{Maîtrise de la forme dissertation :} Les élèves ont produit une copie complète (Intro, Développement 3 parties, Conclusion) respectant les canons de l'examen (probatoire/bac), dans le temps imparti (simulation 4h écrit + 1h oral).
    \item \textbf{Argumentation documentée :} L'exploitation croisée des 4 documents (Littérature/Oyono, Statistiques/INS-MINJEC, Presse/Le Jour) a entraîné la compétence \cle{« mobiliser des connaissances »} et \cle{« exploiter un corpus hétérogène »}, exigée en réception et production.
    \item \textbf{Qualité langagière (Cohésion/Cohérence) :} La relecture croisée guidée par la grille a rendu visible l'usage effectif de la \cle{progression thématique} (reprise, chaînage, progression) et des \cle{connecteurs logiques}, transformant des connaissances déclaratives (cours Langue) en compétences procédurales (rédaction).
    \item \textbf{Insertion des citations :} Le passage obligé par la citation intégrée (Extrait B, Article) avec commentaire a ancré la norme « Bac » (guillemets, référence, analyse).
    \item \textbf{Transfert oral / Écrit :} La défense orale (Séance 3) a révélé l'écart fréquent entre \og savoir rédiger \fg{} et \og savoir justifier sa rédaction \fg{}. L'explicitation du \textbf{choix du plan} et de l'\textbf{ouverture} a consolidé la conscience métacognitive de la structure argumentative.
    \item \textbf{Ancrage camerounais :} Le sujet, le corpus (Oyono, MINJEC, Le Jour, KamerNLP) et l'ouverture (École Technique) ont ancré l'exercice dans la réalité sociolinguistique et technique vécue par les élèves, donnant du sens à l'apprentissage scolaire.
\end{itemize}

\noindent \textbf{Pistes de remédiation post-activité (pour le professeur) :}
\begin{itemize}
    \item \textit{Problème récurrent :} Transition I $\rightarrow$ II souvent faible (simple « Cependant »). \textit{Action :} Exercice ciblé sur les \cle{transitions dialectiques} (reprise du thème + marqueur d'opposition + annonce nouveau thème).
    \item \textit{Problème récurrent :} Ouverture trop générale (« Il faut sensibiliser »). \textit{Action :} Atelier \og Ouvertures concrètes \fg{} : lier obligatoirement l'ouverture à la spécialité technique de l'élève (Génie civil $\rightarrow$ patrimoine architectural ; Informatique $\rightarrow$ corpus linguistiques ; Électrotechnique $\rightarrow$ éclairage sites patrimoniaux).
    \item \textit{Problème récurrent :} Citation « parachutée » sans commentaire. \textit{Action :} Systématiser l'exercice \og 1 citation = 1 commentaire linguistique/sémantique \fg{} en début de séance.
\end{itemize}

\noindent En somme, cet atelier constitue une \textbf{évaluation sommative formatrice} : il certifie les acquis de l'unité « Dissertation » tout en diagnostiquant les points de vigilance pour la préparation finale aux épreuves certificatives.

\newpage

\coursec{Méthodes et exercices résolus}

\coursec{Leçon 14 : Rédiger une dissertation}

\courssub{Méthodes et Exercices Résolus}

\begin{methode}[Analyser un sujet de dissertation : du mot-clé à la problématique]
    \textbf{Objectif :} Maîtriser l'étape préliminaire indispensable avant toute rédaction : l'analyse rigoureuse du sujet pour en extraire la tension dialectique.
    \begin{enumerate}
        \item \textbf{Cercler les mots-clés} (termes forts, verbes, connecteurs logiques) et les \textbf{délimiter} (définitions précises, nuances, oppositions).
        \item \textbf{Identifier le type de sujet} : 
            \begin{itemize}
                \item \cle{Sujet affirmatif} (ex: « La tradition freine le progrès ») $\rightarrow$ Transformer en question.
                \item \cle{Sujet interrogatif} (ex: « La tradition freine-t-elle le progrès ? ») $\rightarrow$ Réponse non binaire.
                \item \cle{Sujet citation} $\rightarrow$ Analyser l'auteur, le contexte, expliciter la pensée.
            \end{itemize}
        \item \textbf{Problématiser} : Formuler la \cle{problématique} (question centrale porteuse d'un débat). Elle ne doit pas être une simple reformulation, mais révéler une \textbf{tension} (ex: « Dans quelle mesure la tradition, garante de l'identité, peut-elle devenir un obstacle à l'émancipation moderne ? »).
        \item \textbf{Construire le plan dialectique} (Thèse / Antithèse / Synthèse) ou \cle{thématique} (Causes / Conséquences / Solutions) en veillant à l'équilibre des parties. Chaque grande partie = 2 à 3 sous-parties (arguments).
    \end{enumerate}
    \textbf{Reflexe Bac :} Toujours écrire le \textbf{brouillon structuré} (mots-clés, définitions, problématique, plan détaillé) avant de rédiger l'introduction. C'est noté au Probatoire.
\end{methode}

\begin{exoresolu}[Analyse et planification : « L'école peut-elle former le citoyen ? »]
    \textbf{Énoncé :} Vous traitez le sujet : \emph{« L'école peut-elle former le citoyen ? »}. Produisez l'analyse complète du sujet (mots-clés, définitions, type, problématique) et le plan détaillé (grandes parties et arguments principaux avec exemples littéraires/historiques).
    \tcblower
    \textbf{Solution détaillée :}
    
    \textbf{1. Analyse des mots-clés et délimitations :}
    \begin{itemize}
        \item \cle{L'école} : Institution d'enseignement formel (primaire, secondaire, supérieur). Lieu de transmission des savoirs, de socialisation secondaire, de certification. \textbf{Nuance} : Distinguer l'institution (programmes, structure) des acteurs (enseignants, pairs).
        \item \cle{Former} : Façonner, instruire, éduquer, transmettre des valeurs, des compétences, un esprit critique. Implique une intentionnalité, une durée, un résultat visible.
        \item \cle{Le citoyen} : Membre d'une communauté politique (État), titulaire de droits et de devoirs, capable de participer à la vie publique (vote, débat, respect des lois), doté d'une conscience civique et morale. \textbf{Opposition} : Sujet (obéissance) vs Citoyen (participation critique).
    \end{itemize}
    
    \textbf{2. Type de sujet :} Sujet \textbf{interrogatif direct} (question fermée en apparence « Oui/Non », mais exigeant une réponse nuancée « Dans quelle mesure »).
    
    \textbf{3. Problématique :}
    \emph{« Si l'école a pour mission officielle de former le citoyen par l'instruction civique et la socialisation, les contradictions entre ses finalités (sélection, obéissance) et l'idéal démocratique (esprit critique, autonomie) ne risquent-elles pas de faire de l'élève un sujet docile plutôt qu'un citoyen acteur ? »}
    
    \textbf{4. Plan détaillé (Dialectique) :}
    
    \textbf{I. L'école, lieu privilégié de la formation citoyenne (Thèse)}
    \begin{itemize}
        \item A. La transmission explicite des valeurs républicaines/démocratiques (programmes d'éducation civique, histoire, hymne, drapeau). \textit{Ex: Cours de "Vie citoyenne" au Cameroun, "Éducation morale et civique" en France.}
        \item B. La socialisation par la vie collective : respect des règles, élection des délégués, apprentissage du débat contradictoire. \textit{Ex: Conseil de vie lycéenne (CVL).}
        \item C. L'accès au savoir comme condition de l'autonomie : lire, écrire, compter, esprit scientifique $\rightarrow$ immunisation contre la manipulation.
    \end{itemize}
    
    \textbf{II. Les limites structurelles et idéologiques de l'institution scolaire (Antithèse)}
    \begin{itemize}
        \item A. La forme scolaire favorise la soumission à l'autorité (maître unique, notation, discipline) $\rightarrow$ formation du \cle{sujet} plus que du citoyen. \textit{Ref: Foucault, "Surveiller et punir" ; Bourdieu, "La reproduction".}
        \item B. L'inégalité sociale scolaire (reproduction des élites) contredit l'égalité citoyenne. L'école légitime les dominants. \textit{Ex: Écarts de réussite selon l'origine sociale (PISA).}
        \item C. Le décalage entre discours officiel (esprit critique) et pratiques pédagogiques (par cœur, conformisme). \textit{Ex: "Le Vieux Nègre et la Médaille" : l'école coloniale formait des auxiliaires, pas des citoyens.}
    \end{itemize}
    
    \textbf{III. Vers une refondation de l'école pour une citoyenneté effective (Synthèse)}
    \begin{itemize}
        \item A. Pédagogies actives et émancipatrices : coopération, projets, classe inversée, philosophie pour enfants (Lipman). L'élève devient acteur.
        \item B. Ouverture de l'école sur la cité : stages, associations, conseils d'élèves réels, éducation aux médias (lutte contre fake news).
        \item C. Formation continue des enseignants à la posture de "facilitateur de citoyenneté" plutôt que de "dépositaire du savoir".
    \end{itemize}
    
    \textbf{Conclusion de l'analyse :} Ce plan assure la progression logique : Constat institutionnel $\rightarrow$ Critique sociologique $\rightarrow$ Perspectives réformistes. Il permet d'insérer des références précises (auteurs, textes du programme, actualité).
\end{exoresolu}

\begin{methode}[Rédiger l'introduction : la méthode de l'entonnoir (Production \& Amélioration)]
    \textbf{Objectif :} Produire une introduction type Bac en 4 étapes logiques (environ 15-20 lignes), puis savoir l'améliorer (vocabulaire, fluidité, accroche).
    \begin{enumerate}
        \item \textbf{L'accroche (Amorce) :} 1 à 2 phrases. \cle{Accroche contextuelle} (actualité, histoire), \cle{littéraire} (citation, référence au programme), ou \cle{généralisante} (définition large, lieu commun à dépasser). \textit{Éviter : "Depuis la nuit des temps..."}
        \item \textbf{La présentation du sujet :} Reprise des termes du sujet, définitions rapides des concepts clés, mise en perspective (enjeux sociaux, philosophiques, littéraires).
        \item \textbf{La problématique :} \cle{La question unique} qui guide toute la dissertation. Elle doit contenir la tension (mots : \emph{paradoxe, tension, ambivalence, dans quelle mesure, jusqu'à quel point}).
        \item \textbf{L'annonce du plan :} \cle{Formule impersonnelle} : « Nous analyserons d'abord... / Nous examinerons ensuite... / Nous montrerons enfin... » ou « Première partie... Deuxième partie... Troisième partie... ». \textbf{Interdiction :} « Je vais dire que... », « Mon plan est... ».
    \end{enumerate}
    \textbf{Amélioration (Réécriture) :} Vérifier la \cle{cohésion} (connecteurs : \emph{or, toutefois, en effet, par conséquent}), la \cle{cohérence} (lien logique accroche $\rightarrow$ sujet $\rightarrow$ problème), le \cle{registre de langue} (soutenu, impersonnel), la \cle{ponctuation}.
\end{methode}

\begin{exoresolu}[Rédaction et amélioration d'une introduction]
    \textbf{Énoncé :} À partir de l'analyse précédente (Sujet : « L'école peut-elle former le citoyen ? »), rédigez une introduction complète. Puis, proposez une version améliorée en justifiant les corrections (vocabulaire, connecteurs, style).
    \tcblower
    \textbf{Solution détaillée :}
    
    \textbf{Version 1 (Brouillon / Production brute) :}
    L'école est une institution importante dans la société. Depuis toujours, on dit que l'école forme le citoyen. Le sujet demande si l'école peut former le citoyen. Former le citoyen, c'est lui apprendre ses droits et devoirs. Mais l'école a aussi des défauts comme les inégalités. La problématique est : L'école forme-t-elle vraiment le citoyen ou pas ? Nous verrons d'abord que l'école forme le citoyen. Ensuite que l'école a des limites. Enfin comment améliorer l'école.
    
    \textbf{Critique de la V1 (Erreurs types) :} 
    \begin{itemize}
        \item Accroche plate ("Depuis toujours...", "institution importante").
        \item Répétitions ("forme le citoyen" 3 fois).
        \item Problématique binaire ("ou pas"), pas de tension.
        \item Annonce du plan au "Je", style oral, pas de connecteurs logiques entre les parties.
        \item Registre familier ("défauts", "améliorer l'école").
    \end{itemize}
    
    \textbf{Version 2 (Améliorée - Niveau Bac/Probatoire) :}
    \emph{« L'instruction publique, selon le mot de Jules Ferry, a pour but de faire des citoyens. »} Cette affirmation fondatrice de l'école républicaine pose d'emblée l'ambition citoyenne de l'institution scolaire. Pourtant, au-delà de la transmission des savoirs académiques, \cle{former le citoyen} suppose de construire un sujet autonome, conscient de ses droits et capable de participation critique à la vie de la Cité. \cle{Or}, entre cette finalité affichée et la réalité des pratiques pédagogiques, une tension persiste. \textbf{Dans quelle mesure l'école, par sa structure même et ses contenus, parvient-elle à transformer l'élève en citoyen acteur, sans le réduire à un sujet docile ?} Pour répondre, nous analyserons \textbf{en premier lieu} comment l'école assure, par ses programmes et sa vie collective, une socialisation civique fondamentale. \textbf{Dans un second temps}, nous mettrons en lumière les contradictions structurelles — reproduction des inégalités, forme scolaire autoritaire — qui entravent cette émancipation. \textbf{Enfin}, nous envisagerons les pistes de refondation pédagogique et institutionnelle pour une citoyenneté véritablement effective.
    
    \textbf{Justifications des améliorations :}
    \begin{itemize}
        \item \textbf{Accroche} : Citation de Ferry (référence culturelle, programme Histoire/EMC), verbe "poser d'emblée".
        \item \textbf{Présentation} : Définition implicite du citoyen ("sujet autonome..."), mot-charnière \cle{Or} pour basculer vers le problème.
        \item \textbf{Problématique} : Formulation "Dans quelle mesure...", tension "citoyen acteur / sujet docile".
        \item \textbf{Plan} : Connecteurs formels (\emph{en premier lieu, dans un second temps, enfin}), tournures impersonnelles, vocabulaire précis (\emph{socialisation civique, reproduction, émancipation, refondation}).
    \end{itemize}
\end{exoresolu}

\begin{methode}[Construire un paragraphe argumenté : Citation, Transition, Progression thématique]
    \textbf{Objectif :} Rédiger un paragraphe unique (un argument = un paragraphe) respectant la structure standard et assurant la cohésion interne (progression thématique) et externe (transition).
    \begin{enumerate}
        \item \textbf{Phrase d'introduction (Thèse partielle) :} Énonce l'argument principal de la partie (ex: « D'abord, l'école assure une socialisation civique par l'enseignement explicite des règles. »).
        \item \textbf{Explication / Analyse :} Développe le mécanisme, la cause, la logique. Utiliser des connecteurs logiques (\emph{car, en effet, ainsi, par conséquent}).
        \item \textbf{Exemple / Illustration / Citation :} \cle{Preuve concrète}. 
            \begin{itemize}
                \item \textbf{Citation courte} (intégrée dans la phrase) : \emph{Comme l'écrit Durkheim, « l'éducation est l'action exercée par les générations adultes sur celles qui ne sont pas encore mûres pour la vie sociale ».}
                \item \textbf{Citation longue} (isolée, retrait, guillemets français \og ... \fg{}) : Introduite par deux-points, commentée \textbf{après}.
            \end{itemize}
        \item \textbf{Commentaire de l'exemple :} \textbf{Obligatoire.} Expliquer \emph{pourquoi} cet exemple prouve l'argument. Ne jamais laisser une citation "parler d'elle-même".
        \item \textbf{Transition (Mini-conclusion + Ouverture) :} Résume l'argument \textbf{ET} annonce le suivant par opposition ou approfondissement. \textit{Modèle :} « Si l'enseignement explicite pose les bases du civisme, \textbf{toutefois} la vie collective à l'école en offre une mise à l'épreuve plus concrète. »
    \end{enumerate}
    \textbf{Progression thématique (Cohésion) :} Choisir un thème fil conducteur (ex: \emph{L'école / Elle / Cette institution / Le système éducatif}) et le reprendre (pronom, synonyme) en début de phrases pour lier les énoncés (progression \cle{continue} ou \cle{linéaire}).
\end{methode}

\begin{exoresolu}[Rédaction d'un paragraphe avec citation et transition]
    \textbf{Énoncé :} Rédigez le \textbf{premier paragraphe de la partie I} (Thèse : L'école forme le citoyen par la socialisation civique explicite) pour le sujet « L'école peut-elle former le citoyen ? ». Intégrez une citation de \emph{Durkheim} ou \emph{Jules Ferry}. Soignez la progression thématique et la transition vers l'argument suivant (vie collective/délégués).
    \tcblower
    \textbf{Solution détaillée :}
    
    \emph{En premier lieu, l'école remplit sa mission citoyenne par une \textbf{transmission explicite et ritualisée des valeurs républicaines}. \textbf{Elle} inscrit au cœur de ses programmes l'apprentissage des droits, des devoirs et des symboles de la Nation, transformant l'ignorance civique en conscience partagée. \textbf{En effet}, l'instruction civique et morale, loin d'être une simple matière, structure le rapport de l'élève à la loi commune. \textbf{Cette institution} trouve sa justification chez Émile Durkheim lorsqu'il affirme : \og L'éducation est l'action exercée par les générations adultes sur celles qui ne sont pas encore mûres pour la vie sociale \fg{}. \textbf{Cette citation} révèle que l'école agit comme le creuset où la société se reproduit elle-même en inculquant à l'individu les normes collectives indispensables à la cohésion. \textbf{Ainsi}, par la connaissance de la Constitution, le respect du drapeau ou l'apprentissage de la laïcité, l'élève intériorise le cadre légal de sa future citoyenneté. \textbf{Toutefois}, si le savoir civique pose le socle théorique, \textbf{il ne saurait suffire sans l'expérience vécue de la démocratie au quotidien au sein même de l'établissement.} }
    
    \textbf{Analyse technique (Pourquoi ça marche) :}
    \begin{itemize}
        \item \textbf{Progression thématique (Continue) :} Thème = \cle{L'école / Elle / Cette institution / Elle}. Reprise pronominale/synonymique en tête de 4 phrases sur 6 $\rightarrow$ Fluidité maximale.
        \item \textbf{Structure Argumentaire} : Thèse partielle $\rightarrow$ Explication (\emph{En effet}) $\rightarrow$ Citation intégrée/introduite $\rightarrow$ \textbf{Commentaire explicite} (\emph{Cette citation révèle que...}) $\rightarrow$ Exemples concrets (Constitution, drapeau, laïcité) $\rightarrow$ Transition.
        \item \textbf{Transition} : \emph{Toutefois} (marque l'opposition/limite) + Résumé (\emph{socle théorique}) + Ouverture (\emph{expérience vécue / démocratie au quotidien}) $\rightarrow$ Appelle logiquement l'argument B (Vie collective/Délégués).
        \item \textbf{Registre} : Impersonnel, vocabulaire précis (\emph{intériorise, creuset, cohésion, socle}).
    \end{itemize}
\end{exoresolu}

\begin{methode}[Maîtriser la progression thématique, la cohésion et la cohérence (Langue Française)]
    \textbf{Objectif :} Transformer un texte "déchiré" en texte fluide en manipulant les chaînes de référence (progression thématique) et les connecteurs (cohésion), pour servir la logique du propos (cohérence).
    \begin{enumerate}
        \item \textbf{Identifier le thème de chaque phrase} (Sujet grammatical ou thème énonciatif). Repérer les \cle{chaînes de référence} (nom $\rightarrow$ pronom $\rightarrow$ synonyme $\rightarrow$ hyperonyme).
        \item \textbf{Choisir le type de progression} adapté à l'argumentation :
            \begin{itemize}
                \item \cle{Progression continue (thème constant)} : Idéal pour l'explication dense d'un même concept (Thème A $\rightarrow$ Thème A $\rightarrow$ Thème A). \textit{Ex: L'école... Elle... Cette institution...}
                \item \cle{Progression linéaire (thème-rhème)} : Le rhème (info nouvelle) de la phrase $n$ devient le thème de la phrase $n+1$. Idéal pour la narration ou l'enchaînement causal. \textit{Ex: L'école forme. \textbf{Cette formation} passe par... \textbf{Ce passage} implique...}
                \item \cle{Progression par énumération (thème multiple)} : Un thème général introduit des sous-thèmes. \textit{Ex: Trois piliers fondent la citoyenneté : le savoir, le vouloir, le pouvoir. \textbf{Le savoir}... \textbf{Le vouloir}... \textbf{Le pouvoir}...}
            \end{itemize}
        \item \textbf{Insérer les connecteurs logiques (Cohésion) :} Choisir selon la relation sémantique : 
            \begin{center}
            \begin{tabularx}{\linewidth}{|l|X|}\hline
            \textbf{Relation} & \textbf{Connecteurs (échantillon)} \\ \hline
            Cause / Explication & \emph{car, en effet, parce que, du fait que} \\ \hline
            Conséquence & \emph{donc, par conséquent, ainsi, c'est pourquoi} \\ \hline
            Opposition / Concession & \emph{mais, toutefois, cependant, pourtant, bien que} \\ \hline
            Ajout / Argumentation & \emph{de plus, en outre, par ailleurs, d'une part... d'autre part} \\ \hline
            Illustration & \emph{par exemple, ainsi, notamment, c'est le cas de} \\ \hline
            \end{tabularx}
            \end{center}
        \item \textbf{Vérifier la cohérence globale :} Le paragraphe répond-il à la thèse annoncée ? L'ordre des arguments est-il logique (chronologique, croissant, dialectique) ?
    \end{enumerate}
    \textbf{Astuce Bac :} Au brouillon, surlignez les thèmes (jaune) et les connecteurs (vert). Si une phrase "flotte" (thème nouveau non annoncé), réécrivez-la ou ajoutez un connecteur de liaison.
\end{methode}

\begin{exoresolu}[Remédiation : Cohésion, Progression et Cohérence]
    \textbf{Énoncé :} Le paragraphe suivant est maladroit (répétitions, ruptures, connecteurs manquants). Réécrivez-le en appliquant la \cle{progression continue} sur le thème « L'exemple colonial » et en insérant les connecteurs logiques adaptés.
    
    \emph{Texte brut :} « L'école coloniale forme des auxiliaires. Le Vieux Nègre reçoit une médaille. Il est humilié. L'école française apprend le français. Les élèves oublient leur culture. Le texte de Mongo Beti montre ça. C'est une critique. L'école ne forme pas le citoyen. »
    
    \tcblower
    \textbf{Solution détaillée :}
    
    \textbf{Analyse des défauts du texte brut :}
    \begin{itemize}
        \item Rupture thématique constante (L'école $\rightarrow$ Le Vieux Nègre $\rightarrow$ Il $\rightarrow$ L'école française $\rightarrow$ Les élèves $\rightarrow$ Le texte $\rightarrow$ C'est $\rightarrow$ L'école).
        \item Phrases nominales/telegraphiques, pas de connecteurs $\rightarrow$ Pas de cohésion.
        \item Incohérence : "L'école française apprend le français" contredit "L'école coloniale" sans nuance.
    \end{itemize}
    
    \textbf{Réécriture corrigée (Progression continue sur "L'école coloniale / Elle / Ce système") :}
    \emph{\textbf{L'école coloniale}, telle que la dépeint Mongo Beti dans \og Le Vieux Nègre et la Médaille \fg{}, \textbf{ne visait pas à former le citoyen} mais à fabriquer des auxiliaires dociles. \textbf{En effet}, \textbf{ce système} instrumentalise l'instruction pour mieux asservir : \textbf{il} décerne des médailles symboliques — comme celle reçue par le Vieux Nègre — qui, \textbf{loin d'honorer le mérite, consacrent l'humiliation et l'aliénation culturelle.} \textbf{Par conséquent}, l'apprentissage imposé de la langue du colonisateur s'accompagne d'un déni systématique des cultures endogènes, \textbf{privant} ainsi l'élève des repères nécessaires à toute autonomie critique. \textbf{Aussi}, \textbf{cette école-là} produit l'exact opposé du citoyen : un sujet déraciné, incapable de penser par lui-même.}
    
    \textbf{Justifications techniques :}
    \begin{itemize}
        \item \textbf{Progression continue} : \cle{L'école coloniale} (P1) $\rightarrow$ \cle{Ce système} (P2) $\rightarrow$ \cle{Il} (P3) $\rightarrow$ \cle{Cette école-là} (P5). Chaîne référentielle solide.
        \item \textbf{Cohésion (Connecteurs)} : \emph{En effet} (explication), \emph{Loin de...} (opposition interne), \emph{Par conséquent} (conséquence logique), \emph{Aussi} (conclusion partielle).
        \item \textbf{Cohérence} : Argument unique (École coloniale = anti-citoyenneté) soutenu par l'analyse textuelle précise (médaille, langue, culture) $\rightarrow$ Preuve textuelle intégrée au service de la thèse dissertation.
        \item \textbf{Style} : Passif/Actif varié, vocabulaire d'analyse (\emph{instrumentalise, asservir, consacre, aliénation, déni, déraciné}).
    \end{itemize}
\end{exoresolu}

\begin{methode}[Commentaire de texte explicatif : « Le Vieux Nègre et la Médaille » (Lectures 2 \& 3)]
    \textbf{Objectif :} Réussir l'exercice de commentaire (ou question de cours) sur l'extrait étudié : identifier la \cle{thèse}, la \cle{structure explicative}, les \cle{procédés rhétoriques} et la \cle{portée critique}.
    \begin{enumerate}
        \item \textbf{Contextualisation précise :} Auteur (Mongo Beti / Alexandre Biyidi Awala), Œuvres (\emph{Le Vieux Nègre et la Médaille}, 1956), Genre (Roman/Récit à thèse, pamphlet anticolonial), Contexte (Cameroun, fin colonisation, loi-cadre Defferre). Personnage : \cle{Meka} (le Vieux Nègre), chef de canton, figure de l'assimilation ratée.
        \item \textbf{Analyse de la structure explicative (Texte explicatif) :} 
            \begin{itemize}
                \item \cle{Thèse} : Dénonciation de la "médaille" comme outil de domination symbolique.
                \item \cle{Arguments} : Inutilité de la médaille (métal sans valeur), ridicule de la cérémonie, incompréhension du récipiendaire, ironie du sort (mort peu après).
                \item \cle{Exemples/Illustrations} : Description physique de la médaille, discours du commandant, réaction de Meka, chute (mort).
            \end{itemize}
        \item \textbf{Étude des procédés stylistiques/rhétoriques (Langue) :}
            \begin{itemize}
                \item \cle{Ironie / Antiphrase} : « \emph{la belle médaille} », « \emph{l'illustre décoré} » (gap entre mot et réalité).
                \item \cle{Champ lexical} : Métal/froid/brillant (médaille) vs Chair/chaud/vieillesse (Meka) $\rightarrow$ Opposition matière/humain.
                \item \cle{Point de vue / Focalisation} : Interne (Meka) $\rightarrow$ Empathie du lecteur ; Narrateur omniscient ironique (jugement sur le colonisateur).
                \item \cle{Registre} : \cle{Pathétique} (vieillesse, mort) + \cle{Satirique} (cérémonie grotesque) + \cle{Tragique} (inéluctabilité).
            \end{itemize}
        \item \textbf{Synthèse (Portée) :} Le texte explicatif devient \cle{texte argumentatif} : l'explication de la cérémonie \emph{est} la démonstration de l'absurdité coloniale. La médaille = \cle{symbole} de l'aliénation.
    \end{enumerate}
    \textbf{Reflexe Bac :} Toujours lier \textbf{Fond} (thème, thèse) et \textbf{Forme} (procédés, structure). Une citation sans analyse de procédé = 0 point.
\end{methode}

\begin{exoresolu}[Commentaire intégré : La médaille, symbole de l'aliénation]
    \textbf{Énoncé :} À partir de l'extrait de la remise de la médaille (Lecture méthodique n°2/3), montrez comment Mongo Beti transforme la description d'une cérémonie officielle en une dénonciation radicale de l'ordre colonial. Vous analyserez la structure du passage, le jeu sur les points de vue et l'ironie du vocabulaire.
    \tcblower
    \textbf{Solution détaillée :}
    
    \textbf{1. Structure explicative détournée : La cérémonie comme procès}
    Le passage suit le rituel officiel (Arrivée $\rightarrow$ Discours $\rightarrow$ Remise $\rightarrow$ Départ), structure classique du \cle{texte explicatif} (chronologique). Mais chaque étape est subvertie :
    \begin{itemize}
        \item \textbf{Le décor} : « \emph{Le soleil darde ses rayons} » (nature indifférente/cruelle) vs solennité attendue.
        \item \textbf{Le discours du Commandant} : Langue de bois administrative (\emph{dévouement, loyauté, services éminents}) vidée de sens par le contexte (Meka a servi l'oppresseur).
        \item \textbf{La remise} : Gestes mécaniques (\emph{épingla, serra la main}) $\rightarrow$ Déshumanisation.
    \end{itemize}
    \textbf{Effet :} L'explication du déroulement \emph{est} l'accusation. La forme explicative sert l'argumentation.
    
    \textbf{2. Le double point de vue : L'ironie dramatique}
    \begin{itemize}
        \item \textbf{Focalisation interne (Meka) :} « \emph{Il ne comprenait pas bien...} », « \emph{Il sentait le poids...} ». Le lecteur partage sa confusion, sa fatigue, sa fierté naïve. \cle{Empathie}.
        \item \textbf{Regard narratif externe (Ironie) :} « \emph{La belle médaille...} », « \emph{l'illustre décoré} ». Le narrateur sait la vérité (métal vil, honneur factice) que le personnage ignore. \cle{Distance critique}.
    \end{itemize}
    \textbf{Résultat :} L'écart entre les deux visions \textbf{crée le sens} : la médaille n'est pas une récompense, c'est un \cle{marque au fer rouge} de la soumission.
    
    \textbf{3. L'ironie lexicale et les champs sémantiques antagonistes}
    \begin{center}
    \begin{tabularx}{\linewidth}{|l|X|X|}\hline
    \textbf{Champ lexical} & \textbf{Médaille / Colonisateur} & \textbf{Meka / Humanité} \\ \hline
    \textbf{Mots} & \emph{métal, froid, brillant, clinquant, épingla, décoré, illustre} & \emph{vieux, ridé, tremblant, chaleur, sueur, chair, vieillesse, mort} \\ \hline
    \textbf{Valeur} & Artificiel, superficiel, mortel & Vivant, profond, digne, éphémère \\ \hline
    \end{tabularx}
    \end{center}
    L'\cle{antiphrase} (« \emph{la belle médaille} » pour un bout de toc) et l'\cle{oxymore} implicite (honneur humiliant) révèlent la \cle{violence symbolique} (Bourdieu) : le colonisateur impose sa lecture du monde.
    
    \textbf{4. Conclusion du commentaire :} Mongo Beti utilise les codes du \cle{texte explicatif} (description, chronologie) pour mieux les faire exploser de l'intérieur. La médaille, objet central, devient le \cle{symbole condensé} de l'aliénation : elle brille pour masquer le vide, elle pèse pour courber l'échine. C'est une \cle{démonstration par l'absurde} de l'impossibilité d'une citoyenneté dans le cadre colonial.
\end{exoresolu}

\begin{methode}[Rédiger la conclusion : Synthèse, Réponse, Ouverture (Production \& Amélioration)]
    \textbf{Objectif :} Clore la dissertation par une conclusion en 3 temps (environ 10-15 lignes) qui ne soit pas une simple répétition, mais une élévation de la réflexion.
    \begin{enumerate}
        \item \textbf{La synthèse (Bilan) :} Reprendre la \cle{progression de la pensée} (pas la liste du plan). « Au terme de cette analyse, il apparaît que... » / « L'étude a permis de mettre en lumière la dialectique suivante : ... » Utiliser des connecteurs de conclusion (\emph{au total, en définitive, force est de constater que}).
        \item \textbf{La réponse à la problématique :} Réponse \cle{nuancée, définitive, formulée en une phrase forte}. Reprendre les termes de la problématique. \textit{Ex: « L'école forme le citoyen \textbf{sous réserve} d'une refonte profonde de ses finalités et de ses méthodes, sans quoi elle ne produit que des sujets scolarisés. »}
        \item \textbf{L'ouverture (Perspective) :} Élargir le débat sans ouvrir un nouveau sujet. 3 types :
            \begin{itemize}
                \item \cle{Ouverture disciplinaire} : Autre auteur, autre œuvre du programme (ex: \emph{L'Enfant noir} de Camara Laye pour l'école coloniale).
                \item \cle{Ouverture transdisciplinaire} : Histoire (lois Ferry), Sociologie (Bourdieu), Droit (Constitution), Actualité (réformes éducatives actuelles au Cameroun).
                \item \cle{Ouverture prospective/philosophique} : Question nouvelle, citation de pensée. \textit{Ex: « Ne faudrait-il pas, comme le suggère Paulo Freire, inventer une "éducation pour la liberté" plutôt qu'une éducation pour l'adaptation ? »}
            \end{itemize}
    \end{enumerate}
    \textbf{Interdits :} « Pour conclure, j'ai dit que... », « J'espère avoir prouvé que... », Introduction d'un \textbf{nouvel argument} non traité, phrase finale banale (« L'avenir nous le dira »).
\end{methode}

\begin{exoresolu}[Rédaction de la conclusion : Synthèse, Réponse, Ouverture]
    \textbf{Énoncé :} Rédigez la conclusion pour le sujet « L'école peut-elle former le citoyen ? » en respectant les 3 temps. Intégrez une ouverture transdisciplinaire (Sociologie/Actualité camerounaise).
    \tcblower
    \textbf{Solution détaillée :}
    
    \emph{Au total, l'analyse conduite révèle que l'école occupe une position \textbf{paradoxale} : elle est simultanément le \textbf{lieu privilégié} de l'accès à la citoyenneté par la transmission des savoirs et des règles collectives, et un \textbf{espace de reproduction} des inégalités sociales et de formatage à l'obéissance qui en contredit l'idéal émancipateur. \textbf{Force est de constater} que l'institution scolaire, en l'état, forme davantage le \textbf{sujet scolarisé} — compétent, certifié, conformiste — que le \textbf{citoyen critique} — autonome, engagé, lucide. \textbf{Néanmoins}, cette tension n'est pas une fatalité : elle appelle une \textbf{refondation éthique et pédagogique} plaçant l'élève au centre comme acteur de sa propre formation civique.}
    
    \emph{\textbf{Dès lors}, la question se déplace : comment transformer l'école de "usine à diplômes" en "atelier de citoyenneté" ? \textbf{L'apport de la sociologie de l'éducation} (Bourdieu, Passeron) nous rappelle que l'école ne peut changer seule sans une transformation des structures sociales qui la nourrissent. \textbf{Dans le contexte camerounais actuel}, marqué par la mise en œuvre de la \cle{Stratégie Nationale de Développement 2020-2030 (SND30)} qui fait de l'éducation un levier de l'émergence, l'enjeu dépasse le cadre scolaire : il s'agit de garantir que l'école forme des \textbf{compétences pour la vie} et non seulement pour l'examen, afin que le \cle{Baccalauréat technique} sanctionne non plus seulement un savoir technique, mais une \textbf{compétence citoyenne} effective.}
    
    \textbf{Analyse des 3 temps :}
    \begin{itemize}
        \item \textbf{Synthèse} : « Paradoxe », « lieu privilégié / espace de reproduction », « sujet scolarisé / citoyen critique ». Vocabulaire conceptuel.
        \item \textbf{Réponse} : « Néanmoins... n'est pas une fatalité... refondation... ». Réponse nuancée (conditionnelle).
        \item \textbf{Ouverture} : Rebond sur « comment transformer » $\rightarrow$ Sociologie (Bourdieu) $\rightarrow$ Actualité nationale précise (SND30, Bac Tech) $\rightarrow$ Perspective forte (compétence citoyenne). Pas de "nouvel argument", mais élévation du débat.
    \end{itemize}
\end{exoresolu}

\begin{methode}[Intégration finale : De la dissertation au commentaire (Cohérence globale)]
    \textbf{Objectif :} Savoir articuler les acquis de la leçon (Dissertation + Commentaire + Langue) pour l'épreuve intégrée du Probatoire/Bac (Sujet unique ou double épreuve).
    \begin{enumerate}
        \item \textbf{Gestion du temps (4h Bac / 3h Probatoire) :}
            \begin{itemize}
                \item 45-60 min : Analyse sujet + Brouillon plan détaillé (Dissertation) \textbf{ET} Lecture active + Repérages (Commentaire).
                \item 1h30-2h : Rédaction Dissertation (Intro, Développement, Conclusion).
                \item 1h-1h30 : Rédaction Commentaire (Intro, 2-3 axes, Conclusion).
                \item 15-30 min : \cle{Relecture croisée} (Orthographe, Cohésion, Respect consignes, Numérotation pages).
            \end{itemize}
        \item \textbf{Transfert de compétences :}
            \begin{itemize}
                \item \cle{Problématique (Dissertation)} $\leftrightarrow$ \cle{Question directrice (Commentaire)}.
                \item \cle{Paragraphe argumenté (Dissertation)} $\leftrightarrow$ \cle{Axe de commentaire (Explication/Analyse)}. Même structure : Idée directrice + Preuve textuelle + Commentaire.
                \item \cle{Progression thématique / Connecteurs} : Communs aux deux exercices.
                \item \cle{Culture littéraire (Le Vieux Nègre...)} : Sert d'exemple dans la dissertation ET objet du commentaire.
            \end{itemize}
        \item \textbf{Check-list finale (Auto-évaluation) :}
            \begin{itemize}
                \item [ ] Sujet collé ? Problématique claire ? Plan annoncé/réalisé ?
                \item [ ] Intro/Conclusion présentes et soignées ?
                \item [ ] 1 argument = 1 paragraphe ? Citations intégrées/commentées ? Transitions ?
                \item [ ] Commentaire : Thèse du texte ? Structure ? Procédés analysés (pas listés) ?
                \item [ ] Langue : Impersonnel, connecteurs, progression thématique, pas de familier ?
                \item [ ] Propreté : Ratures minimes, paragraphes sautés, écriture lisible.
            \end{itemize}
    \end{enumerate}
\end{methode}

\begin{exoresolu}[Simulation Bac : Gestion du temps et Transfert (Cas pratique)]
    \textbf{Énoncé :} Vous avez 4h pour traiter : \textbf{Sujet Dissertation :} « La tradition est-elle un frein au progrès ? » + \textbf{Commentaire :} Extrait de \emph{Le Vieux Nègre et la Médaille} (Scène de la médaille). Établissez votre \textbf{planning minute par minute} et montrez comment l'analyse du texte nourrit la dissertation (argument antithèse : tradition/aliénation).
    \tcblower
    \textbf{Solution détaillée :}
    
    \textbf{Planning type (4h = 240 min) :}
    \begin{center}
    \begin{tabularx}{\linewidth}{|c|X|X|}\hline
    \textbf{Temps} & \textbf{Activité} & \textbf{Objectif / Livrable} \\ \hline
    0h00 - 0h50 & \textbf{Analyse double} & Sujet Dissert : Mots-clés, Problématique, Plan 3 parties (6 args). Texte Com : Lecture 3x, Repérage (Thèse, Structure, 3 axes, Procédés clés). \textbf{Brouillon structuré complet.} \\ \hline
    0h50 - 2h00 & \textbf{Rédaction Dissertation} & Intro (20 min), Développement 3 parties (1h10 = ~23 min/partie), Conclusion (20 min). \textbf{Écrire au net directement si brouillon propre, sinon recopie.} \\ \hline
    2h00 - 3h15 & \textbf{Rédaction Commentaire} & Intro (15 min), 3 Axes développés (1h = 20 min/axe), Conclusion (15 min). \textbf{Intégrer citations courtes.} \\ \hline
    3h15 - 3h45 & \textbf{Relecture CROISÉE} & \textbf{1. Dissert :} Cohérence plan/problématique, transitions, citations, orthographe. \textbf{2. Com :} Réponse à la question, analyse procédés (pas liste), citations exactes. \textbf{3. Langue :} Progression thématique, connecteurs, impersonnel. \\ \hline
    3h45 - 4h00 & \textbf{Mise en page finale} & Numérotation, en-têtes (Nom, Série, Épreuve), séparation nette des 2 copies. \\ \hline
    \end{tabularx}
    \end{center}
    
    \textbf{Transfert concret : Argument Antithèse (Dissertation) alimenté par le Commentaire}
    \begin{itemize}
        \item \textbf{Dissertation (Partie II - Antithèse) :} « La tradition, lorsqu'elle est instrumentalisée par un pouvoir dominant, devient un frein au progrès en figeant les hiérarchies. »
        \item \textbf{Preuve issue du Commentaire (Le Vieux Nègre) :} La "tradition" de la chefferie (Meka) est détournée par le colonisateur. La médaille (symbole moderne) récupère la tradition (fidélité au chef) pour figer l'ordre colonial.
        \item \textbf{Rédaction intégrée :} \emph{« Ainsi, la tradition chefale, au lieu d'être un repère identitaire dynamique, se mue en instrument de conservation de l'ordre établi. L'analyse de la remise de médaille dans \og Le Vieux Nègre et la Médaille \fg{} de Mongo Beti en offre une illustration saisissante : la cérémonie traditionnelle de reconnaissance est détournée par l'administration coloniale pour légitimer la soumission de Meka. Le texte explicatif de la cérémonie devient alors la démonstration de l'aliénation : la "belle médaille" (antiphrase) scelle le pacte de domination, transformant le gardien de la tradition en auxiliaire du progrès colonial. »}
    \end{itemize}
    \textbf{Gain :} Preuve littéraire de haut niveau, cohérence inter-épreuves, gain de temps (argument déjà réfléchi lors du commentaire).
\end{exoresolu}

\begin{attention}[Les 5 erreurs qui coûtent des points au Bac / Probatoire]
    \begin{enumerate}
        \item \textbf{Problématique absente, binaire ou hors-sujet :} Écrire « Qu'est-ce que la tradition ? » ou « La tradition freine-t-elle le progrès ? (Oui/Non) » au lieu d'une question tensionnée (« Dans quelle mesure... »). \textit{Conséquence :} Plan incohérent, note plafond 6/20 en dissertation.
        \item \textbf{Introduction/Conclusion "fantômes" ou "copiées" :} Pas d'accroche, pas de définitions, plan annoncé au « Je », conclusion qui répète mot pour mot l'intro ou qui introduit un nouvel argument. \textit{Conséquence :} Perte de 3 à 4 points sur la note de structure/rédaction.
        \item \textbf{Paragraphe "catalogue" sans argumentation :} Aligner des exemples (Meka, Camara Laye, actualité) sans phrase thèse, sans explication, sans commentaire de citation, sans transition. \textit{Conséquence :} Hors-sujet partiel, note « Contenu » effondrée.
        \item \textbf{Citations "parachutées" ou non analysées :} Mettre une citation longue sans l'introduire, sans la commenter, ou la laisser flotter entre deux paragraphes. \textit{Conséquence :} Preuve non validée, 0 point pour l'exemple.
        \item \textbf{Négliger la langue (Cohésion/Progression) :} Phrases courtes hachées, répétitions du sujet (« L'école... L'école... L'école... »), connecteurs absents ou faux (\emph{mais} au lieu de \emph{donc}), familier (\emph{truc, bidule, genre}). \textit{Conséquence :} Pénalité lourde sur la note de « Langue/Expression » (souvent coefficient 2 ou 3), rend la copie illisible pour le correcteur.
    \end{enumerate}
    \textbf{Règle d'or :} \cle{Le brouillon est votre meilleur ami.} Un brouillon propre (analyse, plan, citations choisies, connecteurs notés) = 50\% de la réussite. Ne commencez jamais la copie nette sans brouillon validé.
\end{attention}

\newpage

\coursec{Exercices}

\courssub{Je vérifie mes connaissances}

\begin{exercice}[QCM : les parties de la dissertation]
Pour chaque question, entoure la (ou les) bonne(s) réponse(s).
\begin{enumerate}
\item Une introduction de dissertation complète comporte :
\begin{enumerate}
\item une amorce (accroche) ;
\item la présentation du sujet et la définition des termes ;
\item la problématique et l'annonce du plan ;
\item une ouverture sur un autre sujet.
\end{enumerate}
\item La problématique est :
\begin{enumerate}
\item un résumé du texte étudié ;
\item une question (ou un ensemble de questions) qui naît de la tension entre les termes du sujet ;
\item la réponse définitive au sujet ;
\item une citation d'auteur.
\end{enumerate}
\item Dans un paragraphe de développement, on trouve dans l'ordre :
\begin{enumerate}
\item argument, exemple, explication de l'exemple, mini-conclusion ;
\item exemple, argument, ouverture ;
\item citation, amorce, bilan ;
\item transition, problématique, conclusion.
\end{enumerate}
\item La conclusion d'une dissertation comporte obligatoirement :
\begin{enumerate}
\item une nouvelle argumentation ;
\item le bilan de la démonstration et la réponse claire à la problématique ;
\item une ouverture ;
\item une citation non commentée.
\end{enumerate}
\end{enumerate}
\end{exercice}

\begin{exercice}[Vrai ou faux — justifie ta réponse]
Réponds par VRAI ou FAUX, puis justifie en une ou deux phrases.
\begin{enumerate}
\item Une citation peut être insérée dans un devoir sans guillemets si elle est courte.
\item Une transition sert uniquement à annoncer la partie suivante.
\item La cohésion d'un texte repose notamment sur les reprises pronominales, les connecteurs logiques et les champs lexicaux.
\item Dans un texte explicatif, l'auteur cherche avant tout à raconter une histoire.
\item Une ouverture doit proposer une piste nouvelle en lien avec le sujet, sans contredire le bilan.
\end{enumerate}
\end{exercice}

\begin{exercice}[Définitions]
Définis en deux ou trois phrases chacune des notions suivantes, telles qu'elles s'appliquent à la dissertation :
\begin{enumerate}
\item la \cle{progression thématique} (thème et propos) ;
\item la \cle{cohérence} textuelle ;
\item l'\cle{amorce} (ou accroche) ;
\item le \cle{texte explicatif}.
\end{enumerate}
\end{exercice}

\begin{exercice}[Repérage dans \emph{Le Vieux Nègre et la Médaille}]
Réponds brièvement à chaque question.
\begin{enumerate}
\item Qui sont Meka et Kelara ? Quel rapport chacun entretient-il avec la médaille annoncée ?
\item Pourquoi dit-on que la médaille est un symbole d'aliénation dans le roman de Ferdinand Oyono ?
\item Relève dans ta mémoire (ou dans le texte) deux procédés ironiques utilisés par le narrateur lors de la cérémonie de remise de la médaille.
\item Quel type de texte (narratif, descriptif, explicatif, argumentatif) domine dans le passage de la cérémonie ? Justifie.
\end{enumerate}
\end{exercice}

\courssub{Je m'entraîne}

\begin{exercice}[Réécrire des introductions « nulles » — 10 min chrono]
Voici deux introductions faibles. Pour chacune : (a) identifie deux défauts précis ; (b) réécris une introduction complète et correcte (amorce, sujet, problématique, annonce du plan).

\textbf{Sujet 1 :} « Le téléphone portable : ami ou ennemi de l'élève ? »

\emph{Introduction nulle :} « Le téléphone portable est un appareil. Tout le monde en a un. Moi j'ai un téléphone portable. Il peut être un ami ou un ennemi. Nous allons voir ça. »

\textbf{Sujet 2 :} « Faut-il apprendre les langues locales à l'école ? »

\emph{Introduction nulle :} « Depuis longtemps les langues existent. Le Cameroun a beaucoup de langues. C'est un sujet intéressant. Je vais donner mon avis. »
\end{exercice}

\begin{exercice}[Rédiger un paragraphe avec citation et transition]
À partir du sujet « Le respect des traditions est-il un frein au développement technique ? », rédige \textbf{un paragraphe de développement complet} (8 à 12 lignes) qui :
\begin{enumerate}
\item commence par un argument clairement formulé ;
\item insère correctement une citation courte (entre guillemets, avec incise) tirée de \emph{Le Vieux Nègre et la Médaille} (par exemple sur l'attitude de Meka face aux traditions du village) ;
\item explique la citation et la relie à l'argument ;
\item se termine par une phrase de transition vers l'argument suivant (que tu annonces sans le développer).
\end{enumerate}
\end{exercice}

\begin{exercice}[Progression thématique imposée]
Voici la chaîne thématique d'un paragraphe (chaque propos devient le thème de la phrase suivante) :

\begin{center}
\begin{tabularx}{\linewidth}{|c|X|}\hline
\rowcolor{popblueL} Phrase & Chaîne imposée \\ \hline
1 & Thème : la médaille / Propos : récompense coloniale \\ \hline
2 & Thème : récompense coloniale / Propos : flatte l'orgueil de Meka \\ \hline
3 & Thème : l'orgueil de Meka / Propos : l'éloigne des siens \\ \hline
4 & Thème : cet éloignement / Propos : révèle l'aliénation \\ \hline
\end{tabularx}
\end{center}

Rédige le paragraphe correspondant en respectant exactement cette chaîne (reprise du propos précédent en tête de phrase, par pronom, synonyme ou groupe nominal). Souligne ensuite dans ton texte chaque reprise thématique.
\end{exercice}

\begin{exercice}[Du récit au texte explicatif]
Le passage de la mort de Meka est un récit. Transforme-le en \textbf{paragraphe explicatif argumentatif} (10 à 15 lignes) intitulé « La médaille coloniale : symbole d'aliénation ». Consignes :
\begin{enumerate}
\item supprime la narration des événements (pas de « puis », « ensuite », « soudain ») ;
\item utilise le présent de vérité générale et des connecteurs explicatifs (\emph{en effet}, \emph{c'est-à-dire}, \emph{ainsi}, \emph{par conséquent}) ;
\item explique le mécanisme de l'aliénation (ce que le personnage croit gagner, ce qu'il perd réellement) ;
\item conclus le paragraphe par une phrase de bilan.
\end{enumerate}
\end{exercice}

\courssub{Je me prépare au Bac}

\begin{exercice}[Sujet type Probatoire — dissertation complète]
\textbf{Sujet :} « Le respect des traditions est-il un frein au développement technique ? »

Tu t'appuieras notamment sur \emph{Le Vieux Nègre et la Médaille} (l'opposition entre les valeurs du village incarnées par Kelara et la fascination de Meka pour les objets et les honneurs coloniaux) et sur ta culture personnelle.
\begin{enumerate}
\item Analyse le sujet : définis les trois termes clés (« respect des traditions », « frein », « développement technique ») et dégage la tension qui fonde la problématique. (3 pts)
\item Formule la problématique sous forme d'une question principale et de deux questions secondaires. (2 pts)
\item Construis le plan détaillé d'une dissertation dialectique (thèse / antithèse / synthèse) : pour chaque partie, donne deux arguments et un exemple précis. (6 pts)
\item Rédige l'introduction complète. (4 pts)
\item Rédige le premier paragraphe de la thèse, avec une citation correctement insérée et une transition finale. (3 pts)
\item Rédige la conclusion complète, avec une ouverture sur le rôle de l'innovation dans les sociétés traditionnelles camerounaises. (2 pts)
\end{enumerate}
\end{exercice}

\begin{exercice}[Sujet type Baccalauréat technique — dissertation]
\textbf{Sujet :} « L'école technique doit-elle former des spécialistes ou des citoyens cultivés ? »
\begin{enumerate}
\item Dégage la problématique et formule-la en une question. (2 pts)
\item Élabore un plan détaillé en deux ou trois parties, en veillant à ne pas opposer artificiellement les deux termes du sujet. (5 pts)
\item Rédige l'introduction (amorce, présentation du sujet, problématique, annonce du plan). (4 pts)
\item Rédige un paragraphe du développement dans lequel tu mobilises un exemple tiré de ta filière technique et une référence littéraire (africaine ou française) correctement citée. (4 pts)
\item Rédige la conclusion : bilan, réponse nuancée à la problématique, ouverture sur la place de l'humanisme dans les métiers techniques. (3 pts)
\item Relis-toi : souligne dans ta copie deux connecteurs logiques, deux reprises thématiques et une citation, et indique en marge leur fonction (cohésion, cohérence, appui de l'argumentation). (2 pts)
\end{enumerate}
\end{exercice}

\begin{exercice}[Sujet de Bac blanc — dissertation avec exigences croisées]
\textbf{Sujet :} « L'intelligence artificielle peut-elle remplacer le savoir-faire artisanal camerounais ? »

\textbf{Consignes générales :} traitement du sujet en trois étapes obligatoires.
\begin{enumerate}
\item \textbf{Analyse et plan.} Définis « intelligence artificielle » et « savoir-faire artisanal » ; montre que le verbe « remplacer » est le cœur du débat ; propose un plan dialectique détaillé (arguments + exemples : sculpture de Foumban, vannerie, couture, mécanique, informatique...). (6 pts)
\item \textbf{Rédaction de l'introduction.} Rédige une introduction complète comportant obligatoirement une citation littéraire ou philosophique pertinente, correctement introduite. (4 pts)
\item \textbf{Rédaction d'un paragraphe explicatif.} Rédige un paragraphe de type explicatif (10 à 15 lignes) expliquant pourquoi un savoir-faire artisanal repose sur une transmission humaine ; respecte une progression thématique continue (chaque phrase reprend un élément de la précédente) et utilise au moins trois connecteurs explicatifs. (5 pts)
\item \textbf{Rédaction de la conclusion.} Rédige la conclusion : bilan équilibré, réponse nuancée, ouverture sur la question éthique (que gagne et que perd une société qui automatise ses gestes techniques ?). (3 pts)
\item \textbf{Langue.} Ton devoir sera évalué sur la correction syntaxique et la cohésion : vérifie les accords, la ponctuation des citations et l'emploi des connecteurs. (2 pts)
\end{enumerate}
\end{exercice}

\newpage

\coursec{Corrigés des exercices}

\begin{corrige}[Exercice 1 — QCM : les parties de la dissertation]
\begin{enumerate}
\item \textbf{Bonnes réponses : a, b, c.} Une introduction complète comporte l'amorce (accroche), la présentation du sujet avec la définition des termes, puis la problématique et l'annonce du plan. La réponse d est fausse : l'ouverture appartient à la \emph{conclusion}, jamais à l'introduction.
\item \textbf{Bonne réponse : b.} La problématique est la question (ou l'ensemble de questions) qui naît de la tension entre les termes du sujet. Ce n'est ni un résumé (a), ni une réponse définitive (c) — la réponse se construit dans le développement et se formule en conclusion — ni une citation (d).
\item \textbf{Bonne réponse : a.} Le schéma du paragraphe de développement est : argument $\rightarrow$ exemple $\rightarrow$ explication de l'exemple $\rightarrow$ mini-conclusion. Les autres propositions mélangent des éléments qui n'appartiennent pas au paragraphe (amorce, problématique, ouverture).
\item \textbf{Bonnes réponses : b, c.} La conclusion comporte obligatoirement le bilan de la démonstration et la réponse claire à la problématique, auxquels s'ajoute classiquement une ouverture. Elle ne doit contenir ni nouvelle argumentation (a), ni citation non commentée (d) : toute citation doit être expliquée.
\end{enumerate}
\textbf{Points de vigilance :} ne pas confondre l'ouverture (conclusion) avec l'amorce (introduction) ; retenir que la problématique est une \emph{question}, pas une réponse.
\end{corrige}

\begin{corrige}[Exercice 2 — Vrai ou faux]
\begin{enumerate}
\item \textbf{FAUX.} Toute citation, même courte, doit être placée entre guillemets et introduite par une incise (par exemple : « selon le narrateur... »). Sans guillemets, il s'agit d'un plagiat ou d'une paraphrase non signalée.
\item \textbf{FAUX.} Une transition a une double fonction : elle \emph{bilanne} la partie qui vient d'être développée \emph{et} annonce la partie suivante. Une transition qui ne fait qu'annoncer est une transition incomplète.
\item \textbf{VRAI.} La cohésion textuelle repose sur des procédés linguistiques : reprises pronominales, connecteurs logiques, champs lexicaux, reprises par synonymes ou groupes nominaux. La cohérence, elle, concerne l'organisation logique des idées.
\item \textbf{FAUX.} Dans un texte explicatif, l'auteur cherche à \emph{expliquer} un phénomène, un mécanisme ou une idée (présent de vérité générale, connecteurs comme \emph{en effet}, \emph{c'est-à-dire}), non à raconter une histoire : raconter est la fonction du texte narratif.
\item \textbf{VRAI.} L'ouverture propose une piste nouvelle (question voisine, élargissement) en lien avec le sujet ; elle ne doit ni contredire le bilan, ni relancer la démonstration.
\end{enumerate}
\end{corrige}

\begin{corrige}[Exercice 3 — Définitions]
\begin{enumerate}
\item \textbf{La progression thématique.} Le \cle{thème} est ce dont on parle dans une phrase (l'élément connu, souvent placé en tête) ; le \cle{propos} est ce qu'on en dit (l'information nouvelle). Dans une progression thématique continue, le propos d'une phrase devient le thème de la phrase suivante, ce qui assure l'enchaînement fluide des idées dans la dissertation.
\item \textbf{La cohérence textuelle.} La cohérence est la qualité d'un texte dont les idées s'organisent logiquement : chaque partie répond à la problématique, les arguments s'enchaînent sans contradiction, et la conclusion découle du développement. Elle se distingue de la cohésion, qui concerne les marques linguistiques de liaison (pronoms, connecteurs).
\item \textbf{L'amorce (ou accroche).} L'amorce est la première phrase de l'introduction : elle capte l'attention du lecteur en situant le sujet dans un cadre général (constat, fait social, référence culturelle). Elle ne doit être ni trop banale (« de tout temps... ») ni hors sujet.
\item \textbf{Le texte explicatif.} Le texte explicatif est un type de texte qui vise à faire comprendre un phénomène, un mécanisme ou une notion. Il emploie le présent de vérité générale, des connecteurs explicatifs (\emph{en effet}, \emph{c'est-à-dire}, \emph{ainsi}, \emph{par conséquent}) et un vocabulaire précis ; il exclut la narration d'événements.
\end{enumerate}
\end{corrige}

\begin{corrige}[Exercice 4 — Repérage dans \emph{Le Vieux Nègre et la Médaille}]
\begin{enumerate}
\item \textbf{Meka} est un ancien porteur, devenu notable de son village ; il attend la médaille avec fierté et fascination : il y voit une consécration personnelle et un honneur suprême. \textbf{Kelara}, son épouse, incarne les valeurs traditionnelles du village ; elle accueille la médaille avec lucidité et distance, car elle mesure ce que Meka a sacrifié (sa terre, donnée aux missionnaires) pour l'obtenir.
\item La médaille est un symbole d'\cle{aliénation} parce que Meka croit gagner la reconnaissance et la dignité, alors qu'il perd réellement l'essentiel : sa terre, sa place dans la communauté, l'estime des siens. La récompense coloniale le coupe de son identité et de son groupe : il admire l'objet qui consacre sa propre dépossession.
\item Procédés ironiques possibles (deux au choix) : le contraste entre la solennité pompeuse de la cérémonie et la misère réelle de Meka ; le discours grandiloquent du commandant qui célèbre le « dévouement » de Meka alors qu'il l'a spolié ; le décalage entre la fierté naïve du vieil homme et le regard lucide ou moqueur des villageois ; l'antiphrase (dire le contraire de ce qu'on pense) dans les éloges officiels.
\item Le passage de la cérémonie est dominé par le texte \textbf{narratif} (avec des pauses descriptives) : on y suit le déroulement d'événements (arrivée des officiels, discours, remise de la médaille, réactions). Mais l'ironie du narrateur lui donne une visée \emph{argumentative} implicite : le récit dénonce la colonisation. (Toute réponse justifiée par des marques textuelles — temps du récit, connecteurs temporels — est acceptée.)
\end{enumerate}
\end{corrige}

\begin{corrige}[Exercice 5 — Réécrire des introductions « nulles »]
\textbf{Sujet 1.} (a) Deux défauts : amorce vide et tautologique (« Le téléphone portable est un appareil ») ; absence de problématique réelle et d'annonce de plan ; on peut ajouter le « moi » personnel, proscrit dans une introduction.

(b) \emph{Proposition de réécriture :} « Outil de communication devenu omniprésent, le téléphone portable s'est imposé jusque dans les salles de classe, au point de transformer les habitudes des élèves. Compagnon quotidien, il donne accès au savoir et facilite les échanges ; mais il peut aussi devenir une source de distraction et de dépendance. Dès lors, le téléphone portable est-il un ami ou un ennemi de l'élève ? Nous verrons d'abord les services qu'il rend à l'élève, puis les dangers qu'il fait courir, avant de montrer que tout dépend de l'usage qu'on en fait. »

\textbf{Sujet 2.} (a) Deux défauts : amorce creuse (« Depuis longtemps les langues existent ») ; aucune définition des termes ni problématique ; annonce de plan absente (« Je vais donner mon avis » ne suffit pas).

(b) \emph{Proposition de réécriture :} « Le Cameroun compte plus de deux cents langues locales, richesse culturelle qui coexiste avec le français et l'anglais, langues officielles de l'école. Apprendre les langues locales à l'école, c'est transmettre une identité ; mais c'est aussi poser la question du temps, des moyens et de l'unité nationale. Faut-il alors enseigner les langues locales à l'école ? Nous examinerons d'abord les arguments en faveur de cet enseignement, puis ses limites pratiques, avant de proposer une solution équilibrée. »

\textbf{Points de vigilance :} une introduction se juge sur quatre éléments repérables : amorce, présentation du sujet (avec définition des termes), problématique interrogative, annonce du plan en deux ou trois parties.
\end{corrige}

\begin{corrige}[Exercice 6 — Rédiger un paragraphe avec citation et transition]
\emph{Proposition de paragraphe (modèle) :}

« Le respect absolu des traditions peut freiner l'adhésion aux techniques nouvelles lorsqu'il enferme l'individu dans le refus de tout changement. Dans \emph{Le Vieux Nègre et la Médaille}, Kelara incarne cette fidélité aux valeurs du village : face aux honneurs coloniaux, elle reste lucide, quand Meka, fasciné, se laisse séduire par la médaille que lui promet l'administration. Le narrateur souligne cette fascination naïve lorsque Meka se réjouit d'être « le premier de son village » à recevoir une telle distinction. Cette réplique montre que le personnage mesure sa valeur à l'aune du regard colonial et non plus à celle de sa communauté : la tradition, qui devrait l'enraciner, devient chez lui un poids dont il cherche à se défaire, sans pour autant gagner une véritable modernité. Ainsi, la tradition mal comprise peut paralyser ; mais, comme nous le verrons maintenant, elle peut aussi, bien comprise, servir de fondement au développement technique. »

\textbf{Barème indicatif (sur 5 pts) :} argument clair en tête (1 pt) ; citation entre guillemets avec incise (1 pt) ; explication de la citation reliée à l'argument (1 pt) ; mini-conclusion (1 pt) ; transition finale annonçant l'argument suivant (1 pt).

\textbf{Points de vigilance :} la citation doit être \emph{courte}, \emph{exacte} et \emph{commentée} ; la transition doit à la fois boucler le paragraphe et annoncer la suite.
\end{corrige}

\begin{corrige}[Exercice 7 — Progression thématique imposée]
\emph{Proposition de paragraphe (les reprises thématiques sont indiquées en gras dans ce corrigé ; l'élève les souligne) :}

« \textbf{La médaille} est d'abord perçue par Meka comme une \textbf{récompense coloniale} prestigieuse. \textbf{Cette récompense coloniale} flatte \textbf{l'orgueil de Meka}, qui se croit enfin reconnu par l'administration. \textbf{Cet orgueil} l'éloigne \textbf{des siens} : il se coupe peu à peu de Kelara et des villageois. \textbf{Cet éloignement} révèle enfin \textbf{l'aliénation} du personnage, qui perd son identité en croyant gagner un honneur. »

\textbf{Points de vigilance :} chaque phrase doit reprendre en tête (thème) le propos de la phrase précédente, par pronom, synonyme ou groupe nominal ; ne pas introduire d'élément nouveau en tête de phrase ; vérifier que la chaîne complète forme un raisonnement cohérent (de la récompense à l'aliénation).
\end{corrige}

\begin{corrige}[Exercice 8 — Du récit au texte explicatif]
\emph{Proposition de paragraphe explicatif (modèle) :}

\textbf{La médaille coloniale : symbole d'aliénation.} « La médaille coloniale fonctionne comme un symbole d'aliénation parce qu'elle inverse les valeurs aux yeux de celui qui la reçoit. En effet, Meka croit gagner la dignité et la reconnaissance, alors qu'il a déjà perdu l'essentiel : sa terre, cédée aux missionnaires, et sa place dans la communauté du village. C'est-à-dire que l'honneur affiché masque une dépossession réelle. L'objet brillant agit comme un leurre : il fait admirer au bénéficiaire le système même qui l'a spolié. Ainsi, le vieil homme se coupe de Kelara et des siens, dont le regard lucide mesure l'étendue de sa perte. Par conséquent, la médaille ne récompense pas un mérite ; elle consacre une soumission. Elle transforme la victime en admirateur de son propre malheur : telle est la définition même de l'aliénation. »

\textbf{Barème indicatif (sur 5 pts) :} suppression de toute narration (1 pt) ; présent de vérité générale (1 pt) ; au moins trois connecteurs explicatifs correctement employés (1 pt) ; mécanisme de l'aliénation clairement expliqué (gain illusoire / perte réelle) (1 pt) ; phrase de bilan finale (1 pt).

\textbf{Points de vigilance :} bannir « puis », « ensuite », « soudain » ; ne pas raconter la mort de Meka mais \emph{expliquer} ce qu'elle signifie.
\end{corrige}

\begin{corrige}[Exercice 9 — Sujet type Probatoire : traditions et développement technique]
\textbf{1. Analyse du sujet (3 pts).} Le « respect des traditions » désigne la fidélité aux coutumes, aux valeurs et aux savoirs transmis par les anciens. Un « frein » est ce qui ralentit ou empêche un mouvement. Le « développement technique » renvoie au progrès des outils, des techniques et des conditions de vie matérielle. Tension : la fidélité au passé semble s'opposer au mouvement vers l'avenir ; le sujet demande si cette opposition est réelle.

\textbf{2. Problématique (2 pts).} Question principale : le respect des traditions est-il incompatible avec le développement technique ? Questions secondaires : les traditions peuvent-elles s'opposer au progrès ? Peuvent-elles au contraire le nourrir et l'orienter ?

\textbf{3. Plan dialectique détaillé (6 pts).}
\begin{itemize}
\item \textbf{Thèse — la tradition peut freiner le progrès.} Arg. 1 : le refus du changement au nom des coutumes (ex. : certaines communautés refusant les techniques agricoles nouvelles). Arg. 2 : la fascination inverse pour le « moderne » détruit la tradition sans construire (ex. : Meka, qui sacrifie sa terre pour une médaille).
\item \textbf{Antithèse — la tradition peut fonder le développement.} Arg. 1 : les savoirs traditionnels sont des techniques (ex. : la pharmacopée, la métallurgie et la sculpture de Foumban). Arg. 2 : Kelara incarne une sagesse qui jauge lucidement le faux progrès colonial.
\item \textbf{Synthèse — un développement enraciné.} Arg. 1 : le progrès durable s'appuie sur les valeurs locales (ex. : les tontines, modèle d'épargne traditionnelle devenu outil économique moderne). Arg. 2 : moderniser sans se renier (ex. : artisans camerounais utilisant Internet pour vendre leurs œuvres).
\end{itemize}

\textbf{4. Introduction (4 pts) — modèle :} « Toute société se construit dans le dialogue entre l'héritage des anciens et les inventions du présent. Au Cameroun, les traditions structurent encore la vie des villages, tandis que le développement technique transforme les villes et les campagnes. Le respect des traditions désigne la fidélité aux valeurs transmises ; le développement technique, le progrès des outils et des savoir-faire. Mais ces deux exigences sont-elles contradictoires ? Le respect des traditions est-il un frein au développement technique ? Nous verrons d'abord que la tradition peut entraver le progrès, puis qu'elle peut aussi le fonder, avant de montrer qu'un véritable développement s'enracine dans la culture qui le porte. »

\textbf{5. Premier paragraphe de la thèse (3 pts) — modèle :} « Le respect absolu des traditions peut d'abord apparaître comme un obstacle au progrès, lorsqu'il impose l'immobilisme. Dans \emph{Le Vieux Nègre et la Médaille}, Ferdinand Oyono montre un village partagé entre la fidélité aux coutumes et l'attrait des nouveautés coloniales : Meka, séduit, attend sa médaille comme « un grand honneur » qui le distinguera des autres. Cette attente révèle que la tradition, vécue comme un poids, pousse paradoxalement l'individu à tout sacrifier — sa terre, sa dignité — pour un progrès illusoire. Ainsi, la fidélité mal comprise ou abandonnée au profit du mirage moderne bloque tout développement authentique. Mais la tradition est-elle nécessairement un obstacle ? C'est ce qu'il faut maintenant examiner. »

\textbf{6. Conclusion (2 pts) — modèle :} « En définitive, le respect des traditions n'est un frein au développement technique que lorsqu'il devient refus de tout changement ; bien compris, il offre au contraire des valeurs et des savoirs qui fondent un progrès durable. La médaille de Meka nous avertit : courir après une modernité qui dépossède, c'est perdre deux fois. Reste à se demander si l'innovation, dans les sociétés traditionnelles camerounaises, n'a pas toujours été la condition même de leur survie. »

\textbf{Points de vigilance :} définir les trois termes avant la problématique ; éviter le hors-sujet « pour ou contre les traditions » ; la synthèse ne doit pas être un simple compromis mais un dépassement ; citer le roman au moins une fois avec guillemets et incise.
\end{corrige}

\begin{corrige}[Exercice 10 — Sujet type Baccalauréat technique : spécialistes ou citoyens cultivés]
\textbf{1. Problématique (2 pts).} L'école technique doit-elle se limiter à former des spécialistes compétents, ou doit-elle aussi former des citoyens cultivés — et ces deux missions sont-elles vraiment opposables ?

\textbf{2. Plan détaillé (5 pts) — sans opposition artificielle :}
\begin{itemize}
\item \textbf{Partie 1 — la formation spécialisée est nécessaire.} Arg. 1 : l'insertion professionnelle exige des compétences précises (ex. : un technicien en froid doit maîtriser les circuits frigorifiques). Arg. 2 : le développement du pays a besoin de techniciens qualifiés.
\item \textbf{Partie 2 — la culture générale est indispensable au technicien.} Arg. 1 : le technicien communique, rédige, argumente (ex. : rédiger un rapport de stage ou un devis). Arg. 2 : la littérature forme le jugement et la conscience (ex. : \emph{Le Vieux Nègre et la Médaille} apprend à lire les mécanismes de domination).
\item \textbf{Partie 3 — les deux missions se complètent.} Arg. 1 : la culture rend la technique humaine et responsable. Arg. 2 : le citoyen cultivé est un meilleur spécialiste (ex. : l'ingénieur qui comprend les besoins de sa communauté innove mieux).
\end{itemize}

\textbf{3. Introduction (4 pts) — modèle :} « À l'heure où le Cameroun mise sur la formation technique pour son industrialisation, une question traverse les ateliers et les salles de classe : que doit enseigner l'école technique ? Former des spécialistes, c'est transmettre un savoir-faire précis ; former des citoyens cultivés, c'est donner à chacun les moyens de penser, de juger et de vivre en société. Ces deux missions s'excluent-elles ? L'école technique doit-elle former des spécialistes ou des citoyens cultivés ? Nous montrerons d'abord la nécessité de la spécialisation, puis celle de la culture générale, avant de prouver que les deux se complètent. »

\textbf{4. Paragraphe du développement (4 pts) — modèle :} « La culture générale est d'abord indispensable au technicien lui-même, car aucun métier ne s'exerce dans l'isolement. Dans ma filière, un élève en génie civil doit rédiger des rapports de chantier clairs et convaincre des clients : sans maîtrise de la langue, la meilleure compétence technique reste muette. La littérature, elle, aiguise le jugement : dans \emph{Le Vieux Nègre et la Médaille}, Ferdinand Oyono montre Meka ébloui par « l'honneur » que lui promet l'administration coloniale, au point de ne plus voir sa propre dépossession. Cette lucidité que le roman nous transmet — distinguer l'apparence de la réalité — est précieuse pour tout professionnel confronté à des choix. Ainsi, la culture n'est pas un ornement : elle est un outil de travail. »

\textbf{5. Conclusion (3 pts) — modèle :} « En définitive, opposer le spécialiste et le citoyen cultivé serait artificiel : l'école technique forme nécessairement les deux. La compétence sans culture produit des exécutants ; la culture sans compétence, des rêveurs. Le technicien accompli réunit les deux. C'est toute la place de l'humanisme dans les métiers techniques : rappeler que derrière chaque machine, chaque ouvrage, il y a des hommes à servir. »

\textbf{6. Relecture (2 pts).} Exemples attendus : connecteurs logiques soulignés (\emph{d'abord}, \emph{ainsi}, \emph{en définitive} — fonction : cohésion et organisation du raisonnement) ; reprises thématiques (\emph{cette lucidité}, \emph{les deux} — fonction : cohésion) ; citation (« l'honneur » — fonction : appui de l'argumentation).

\textbf{Points de vigilance :} ne pas transformer le plan en « pour / contre » sec ; la référence littéraire doit être citée et commentée ; la conclusion doit répondre \emph{nuancée} à la question (ni « ou » sec, ni « et » mou sans justification).
\end{corrige}

\begin{corrige}[Exercice 11 — Sujet de Bac blanc : intelligence artificielle et savoir-faire artisanal]
\textbf{1. Analyse et plan (6 pts).} L'« intelligence artificielle » désigne des programmes capables d'accomplir des tâches qui demandaient l'intelligence humaine (calcul, reconnaissance, production). Le « savoir-faire artisanal camerounais » renvoie aux techniques manuelles transmises de maître à apprenti : sculpture de Foumban, vannerie, couture, forge, mécanique de rue. Le verbe « remplacer » est le cœur du débat : il suppose une équivalence entre la machine et l'homme, qu'il faut interroger.

Plan dialectique :
\begin{itemize}
\item \textbf{Thèse — l'IA peut remplacer une part du savoir-faire.} Arg. 1 : rapidité et précision de la machine (ex. : découpe automatisée, couture industrielle). Arg. 2 : coût réduit et production en série (ex. : meubles fabriqués en usine).
\item \textbf{Antithèse — l'artisanat est irremplaçable.} Arg. 1 : le geste artisanal porte une culture et une âme (ex. : chaque statue de Foumban est unique). Arg. 2 : la transmission humaine crée lien social et emploi (ex. : l'apprentissage chez un maître-artisan).
\item \textbf{Synthèse — l'IA au service de l'artisan.} Arg. 1 : la machine comme outil, non comme remplaçant (ex. : l'artisan qui vend ses œuvres en ligne). Arg. 2 : former des artisans capables d'utiliser le numérique (ex. : mécanicien utilisant le diagnostic électronique).
\end{itemize}

\textbf{2. Introduction (4 pts) — modèle avec citation :} « De l'atelier du sculpteur de Foumban au garage du quartier, le savoir-faire artisanal fait vivre des milliers de Camerounais. Or l'intelligence artificielle, capable désormais de dessiner, calculer et fabriquer, semble menacer ces métiers anciens. Un proverbe rappelle : « La main est la fenêtre de l'esprit » : peut-on alors confier à la machine ce que la main humaine enseigne depuis des siècles ? L'intelligence artificielle peut-elle remplacer le savoir-faire artisanal camerounais ? Nous verrons d'abord ce que la machine peut faire à la place de l'artisan, puis ce qui, dans l'artisanat, lui échappe, avant de montrer que l'avenir appartient à leur alliance. »

\textbf{3. Paragraphe explicatif (5 pts) — modèle :} « Un savoir-faire artisanal repose sur une transmission humaine, parce que le geste ne s'apprend pas dans un manuel. En effet, l'apprenti sculpteur observe son maître, imite ses mouvements et corrige les siens sous son regard. Cette relation maître-élève, c'est-à-dire un accompagnement quotidien et patient, transmet bien plus qu'une technique : elle communique un sens du beau, des valeurs et une histoire. Ainsi, le savoir-faire devient un héritage vivant, que chaque génération adapte à son époque. Par conséquent, une machine, même perfectionnée, peut reproduire un objet mais non transmettre cette culture du geste. En somme, l'artisanat est une école de l'homme autant qu'une production d'objets. »

\textbf{4. Conclusion (3 pts) — modèle :} « En définitive, l'intelligence artificielle peut assister l'artisan, accélérer certaines tâches et ouvrir de nouveaux marchés ; elle ne saurait remplacer le savoir-faire camerounais, qui est transmission humaine, culture et identité. La réponse est donc nuancée : oui au progrès, non à la substitution. Reste une question éthique : une société qui automatise ses gestes techniques gagne en efficacité, mais que perd-elle en mémoire, en lien et en humanité ? »

\textbf{5. Langue (2 pts) : points de vigilance.} Vérifier les accords sujet-verbe et participe passé ; ponctuer les citations (deux-points ou virgule avant les guillemets, incise) ; employer des connecteurs variés et justes (\emph{en effet} = cause, \emph{c'est-à-dire} = reformulation, \emph{par conséquent} = conséquence) ; assurer les reprises thématiques pour la cohésion.

\textbf{Barème indicatif global :} analyse et plan 6 pts ; introduction 4 pts ; paragraphe explicatif 5 pts ; conclusion 3 pts ; langue 2 pts ; total 20 pts.
\end{corrige}

\newpage

\coursec{Fiche bilan}

\begin{popfigure}
\begin{tikzpicture}[
    mindmap,
    concept color=popblue,
    level 1 concept/.append style={level distance=4.5cm, sibling angle=360/7},
    level 2 concept/.append style={level distance=3cm},
    every node/.style={font=\small, text=popdark},
    branch1/.style={concept color=poporange, grow=0},
    branch2/.style={concept color=popgreen, grow=51},
    branch3/.style={concept color=poppurple, grow=102},
    branch4/.style={concept color=poppink, grow=153},
    branch5/.style={concept color=popblue, grow=204},
    branch6/.style={concept color=popgold, grow=255},
    branch7/.style={concept color=popmuted, grow=306},
    text width=2.8cm, align=center
]
\node[concept, text width=3.2cm] {\textbf{La\\ Dissertation}}
    child[branch1] { node[concept, align=center]{\textbf{1. Analyse\\ du sujet}\\
        \textit{Mots-clés, type,\\ tension implicite}} }
    child[branch2] { node[concept, align=center]{\textbf{2. Introduction}\\
        \textit{Accroche, sujet posé,\\ problématique, plan}} }
    child[branch3] { node[concept, align=center]{\textbf{3. Développement}\\
        \textit{Plan dialectique ou\\ thématique, paragraphe\\ P-E-E, citations,\\ transitions}} }
    child[branch4] { node[concept, align=center]{\textbf{4. Conclusion}\\
        \textit{Bilan synthétique,\\ réponse à la problématique,\\ ouverture}} }
    child[branch5] { node[concept, align=center]{\textbf{5. Cohésion /\\ cohérence}\\
        \textit{Progression thématique,\\ connecteurs logiques,\\ reprises anaphoriques}} }
    child[branch6] { node[concept, align=center]{\textbf{6. Étude intégrée}\\
        \textit{\emph{Le Vieux Nègre\\ et la Médaille}\\ (F. Oyono)\\ Lectures 2 \& 3,\\ texte explicatif}} }
    child[branch7] { node[concept, align=center]{\textbf{7. Rédaction /\\ remédiation}\\
        \textit{Amélioration du brouillon,\\ auto-évaluation,\\ grille de critères}} };
\end{tikzpicture}
\legende{Carte mentale : architecture de la dissertation en Terminale technique (7 piliers).}
\end{popfigure}

\courssub{Définitions clés}

\begin{itemize}
    \item \cle{Dissertation} : exercice d'argumentation structuré répondant à une \cle{problématique} issue de l'\cle{analyse du sujet}, organisé en \cle{introduction}, \cle{développement} (plan dialectique : thèse/antithèse/synthèse, ou plan thématique) et \cle{conclusion}.
    \item \cle{Problématique} : question ouverte, issue de la tension contenue dans le sujet, qui guide tout le devoir. Elle ne se subit pas, elle se \textbf{construit} à partir des mots-clés.
    \item \cle{Introduction} (schéma de l'entonnoir) : \textbf{Accroche} (citation, fait d'actualité, définition) $\rightarrow$ \textbf{Sujet posé} (reformulation) $\rightarrow$ \textbf{Problématique} $\rightarrow$ \textbf{Annonce du plan} (mots-clés des grandes parties).
    \item \cle{Paragraphe type (P-E-E)} : \textbf{P}hrase-topic (idée directrice) $\rightarrow$ \textbf{E}xplication / arguments $\rightarrow$ \textbf{E}xemple / citation / donnée $\rightarrow$ \textbf{mini-transition} vers le paragraphe suivant.
    \item \cle{Transition} : phrase-charnière liant deux parties ou deux paragraphes : \textit{rappel de l'idée achevée} + \textit{annonce de l'idée à venir} + \textit{marque de la logique du lien} (concession, approfondissement, opposition).
    \item \cle{Conclusion} (entonnoir inversé) : \textbf{bilan synthétique} (réponse aux grandes parties) $\rightarrow$ \textbf{réponse directe à la problématique} $\rightarrow$ \textbf{ouverture} (autre sujet, autre œuvre, actualité, perspective).
    \item \cle{Progression thématique} : organisation de l'information (thème $\rightarrow$ propos) : \textit{linéaire} (le propos d'une phrase devient le thème de la suivante), \textit{à thème constant} (un même thème est repris), \textit{à thème dérivé} (les thèmes dérivent d'un hyperthème). Elle est garante de la \cle{cohérence} du texte.
    \item \cle{Cohésion} : ensemble des procédés de liaison de surface du texte : connecteurs logiques, reprises anaphoriques (pronoms, démonstratifs), substitutions lexicales, champs lexicaux.
    \item \cle{Texte explicatif} : texte qui vise à faire comprendre un phénomène ou un mécanisme. Structure : définition $\rightarrow$ explication (causes, étapes, conséquences) $\rightarrow$ illustration / exemple. Marqueurs : connecteurs logiques (car, donc, par conséquent, en effet), présent de vérité générale, vocabulaire précis.
\end{itemize}

\courssub{Méthodes express (à maîtriser par cœur)}

\begin{enumerate}
    \item \textbf{Analyser le sujet (10 min)} : souligner les mots-clés $\rightarrow$ définir et paraphraser $\rightarrow$ identifier le type de sujet (citation, question, affirmation) $\rightarrow$ dégager la tension (pour/contre, ancien/nouveau) $\rightarrow$ formuler la problématique.
    \item \textbf{Rédiger l'introduction (15 min)} : 1. accroche pertinente ; 2. sujet posé (reformulation personnelle) ; 3. problématique (une phrase interrogative) ; 4. plan annoncé (verbes d'action : \textit{analyser, montrer, démontrer, nuancer}).
    \item \textbf{Construire le développement (1h30)} : 1. fiche de brouillon : idées directrices + arguments + citations (précises, courtes, intégrées) ; 2. respecter l'architecture (2 ou 3 grandes parties, 2 à 3 paragraphes par partie) ; 3. rédiger des paragraphes P-E-E ; 4. soigner les transitions (une par changement de grande partie, une mini-transition entre paragraphes).
    \item \textbf{Rédiger la conclusion (15 min)} : 1. bilan (reprise des mots-clés des parties) ; 2. réponse nette à la problématique ; 3. ouverture pertinente (éviter les formules creuses du type « pour conclure je dirais que... »).
    \item \textbf{Relire et corriger (15 min)} : 1. cohérence (sujet respecté ? progression logique ?) ; 2. cohésion (connecteurs, pronoms, progression thématique) ; 3. langue (orthographe, accords, syntaxe, ponctuation, vocabulaire varié).
\end{enumerate}

\courssub{Repères : \emph{Le Vieux Nègre et la Médaille} (Ferdinand Oyono) — Lectures 2 \& 3}

\begin{center}
\begin{tabularx}{\linewidth}{|l|X|}\hline
\rowcolor{popblueL} \textbf{Aspect} & \textbf{Points saillants pour la dissertation} \\ \hline
\cle{Thèmes majeurs} & Condition coloniale, illusion et désillusion, dignité bafouée, ingratitude du colonisateur, dérision des récompenses (la médaille), éveil de la conscience. \\ \hline
\cle{Structure narrative} & Récit à la première personne (témoignage de Meka), progression vers la chute ironique (la médaille dérisoire face aux humiliations subies). \\ \hline
\cle{Style / Rhétorique} & Ironie et satire, registre pathétique puis lucide, oralité, images fortes de l'humiliation, peinture critique de l'administration coloniale et de l'Église. \\ \hline
\cle{Portée argumentative} & Exemple type pour illustrer : l'ingratitude et l'hypocrisie du système colonial, la valeur du témoignage littéraire, la littérature africaine comme arme de dénonciation et de résistance. \\ \hline
\cle{Texte explicatif} & Analyser comment Oyono \textbf{explique} la condition du colonisé : présentation du personnage et de son statut $\rightarrow$ explication des causes de sa désillusion (promesses non tenues) $\rightarrow$ illustration par le vécu (la remise de la médaille). \\ \hline
\end{tabularx}
\end{center}

\begin{attention}[Erreurs fréquentes à éviter]
\begin{itemize}
    \item Attribuer l'œuvre au mauvais auteur : \emph{Le Vieux Nègre et la Médaille} est un roman de \textbf{Ferdinand Oyono} (écrivain camerounais), et non de Léopold Sédar Senghor.
    \item Réciter le cours au lieu de répondre au sujet posé : toute idée doit être reliée à la problématique.
    \item Juxtaposer des citations sans les intégrer ni les commenter : une citation s'introduit, s'encadre de guillemets et s'exploite.
    \item Annoncer un plan puis ne pas le suivre : la cohérence entre l'annonce et le développement est évaluée.
    \item Terminer sans répondre explicitement à la problématique dans la conclusion.
\end{itemize}
\end{attention}

\begin{aretenir}[Checklist avant le Bac — Dissertation]
\begin{itemize}
    \item $\square$ Je sais \textbf{analyser un sujet} en 5 étapes (mots-clés, type, tension, reformulation, problématique).
    \item $\square$ Je connais la \textbf{structure canonique de l'introduction} (accroche, sujet posé, problématique, plan annoncé) et je la rédige sans hors-sujet.
    \item $\square$ Je maîtrise l'\textbf{architecture du développement} (plan dialectique : thèse / antithèse / synthèse, \textbf{ou} plan thématique cohérent) et j'annonce mes parties clairement.
    \item $\square$ Je rédige des \textbf{paragraphes P-E-E} : phrase-topic claire, arguments développés, citation ou exemple \textbf{intégré} (guillemets, référence, commentaire).
    \item $\square$ J'écris des \textbf{transitions explicites} entre les grandes parties (rappel + annonce + logique) et des mini-transitions entre les paragraphes.
    \item $\square$ Ma \textbf{conclusion} fait un bilan synthétique, répond \textbf{précisément} à la problématique et propose une ouverture pertinente.
    \item $\square$ J'utilise la \textbf{progression thématique} (linéaire, à thème constant, à thème dérivé) et des \textbf{connecteurs logiques} variés pour assurer cohésion et cohérence.
    \item $\square$ Je sais exploiter \emph{\textbf{Le Vieux Nègre et la Médaille}} de Ferdinand Oyono comme référence argumentative (thèmes, style, portée) ou comme corpus pour le texte explicatif.
    \item $\square$ Je gère mon \textbf{temps} : analyse (10 min) + plan détaillé (15 min) + rédaction (1h45) + relecture (15 min), soit environ 3 heures.
    \item $\square$ Je relis \textbf{systématiquement} : respect du sujet, logique du raisonnement, qualité de la langue (accords, syntaxe, vocabulaire), propreté de la copie.
    \item $\square$ Je connais les \textbf{critères d'évaluation} : pertinence (sujet/problématique), structure et organisation, argumentation et exemples, expression et langue, présentation.
\end{itemize}
\end{aretenir}

\findecours

\end{document}
